%% last revised by wdj on 2025-06-08

\documentclass[12pt]{article}
\usepackage{graphicx}
\usepackage{rotating}
\usepackage{url}
\usepackage{amsfonts}
\usepackage{amsmath}
\usepackage{amsthm}
\usepackage{makeidx}
\usepackage{latexsym}
\usepackage{oplotsymbl}
\usepackage{adjustbox}
\usepackage{tcolorbox}
\usepackage{overpic}

% Define some colors
\definecolor{RithmoBlue}{RGB}{46, 86, 137}
\definecolor{RithmoGold}{RGB}{218, 165, 32}
\definecolor{RithmoWhite}{RGB}{245, 245, 245}
\definecolor{RithmoBlack}{RGB}{50, 50, 50}

% Section command for poster
\newcommand{\postersection}[1]{%
    % \tcbmaketheletter removed as it's often paired with the advanced title styles
    % and might not be necessary for the simplified version.
    \begin{tcolorbox}[title=#1, noparskip]
}
\newcommand{\postersectionend}{\end{tcolorbox}}

\newtheorem{theorem}{Theorem}
\newtheorem{corollary}[theorem]{Corollary}
\newtheorem{conjecture}[theorem]{Conjecture}
\newtheorem{lemma}[theorem]{Lemma}
\newtheorem{proposition}[theorem]{Proposition}
\newtheorem{definition}[theorem]{Definition}
\newtheorem{example}[theorem]{Example}
\newtheorem{axiom}{Axiom}
\newtheorem{remark}{Remark}
\newtheorem{question}{Question}
\newtheorem{exercise}{Exercise}[section]

\def\ppp{{\mathbb{P}}}
\def\hhh{{\mathbb{H}}}
\def\aaa{{\mathbb{A}}}
\def\fff{{\mathbb{F}}}
\def\qqq{\mathbb{Q}}
\def\rrr{\mathbb{R}}
\def\ccc{\mathbb{C}}
\def\zzz{\mathbb{Z}}
\def\nnn{\mathbb{N}}
\def\pf{{\bf proof}:\ }
\def\qed{$\Box$}
\def\rithmomachia{$\frak{rithmomachia}\ $}

\makeindex

\begin{document}



\author{\normalsize by David Joyner\thanks{email:
    {\tt wdjoyner@gmail.com}.}}
\title{\begin{overpic}[height=8.3cm,width=12.6cm]{rithmomachia-initial-state_cover-image}
    \put(15,40){\bf {\it The recreational arithmetic of}}
    \put(30,25){\huge{\bf \rithmomachia}}
\end{overpic}}
%The recreational arithmetic of\\ Rithmomachia}
\date{2025-06-08}
  
\maketitle
%\pagestyle{empty}
\pagenumbering{arabic}

\begin{quotation}
Rithmomachia (or \rithmomachia if you prefer:-)
is a two-player board game invented in the middle ages.
According to medieval mathematics scholars,
using the arithmetically-based rules governing the
rithmomachia pieces helps students with
logical and spiritual reasoning, allowing them to
attain a closer connection the cosmos.
Indeed, some cathedral-based schools at the time {\it banned} chess
which they believed could lead to vices (such as gambling!),
but allowed their students to play rithmomachia.

In the sections below, the reader will find an introduction to the game,
including a detailed description of the rules regarding the board and pieces.
In particular, we describe the interesting way the pieces are
set up initially, how they move, how they capture,
and most importantly, how to win!
%\footnote{Unfortunately, how to {\it defend} is not so easily described.}!
Chess-like notation and
easily readible board diagrams are used for examples and
even a move-by-move presentation of two complete games.

\end{quotation}

\vskip.1in
\postersection{Versions}
Given the different rules among the different versions of rithmomachia,
we narrow our focus here 
to the rules of the game described by
Claude de Boissi\`ere in 1556 \cite{Ri46}.
%See \S \ref{sec:rules} below for details.
\index{de Boissi\`ere, Claude}
\postersectionend

\vskip .3in
%\newpage

\tableofcontents

%\newpage

\vskip .3in

As you can tell from the table of contents, the organization is pretty simple.
First, we talks about the history of the game,
which involves an unusual combination of mathematics, religion and
philosophy. After that background is sketched out, the game rules are
covered, with examples. In the next section a complete game,
using the
notation introduced earlier, is given. Finally, there is a summary with
some open questions. Fun exercises and problems
(and cartoons:-) are scattered throughout.


\begin{figure}[h!] 
\centering 
\includegraphics[height=2.0cm,width=15.0cm]{Rithmomachia-game-pieces/white-triangles-green-border-in-line} 
\caption{Rithmomachia White triangles pieces in a line. {\small{(The graphics for the pieces were made using
{\tt SageMath} and Python.)}}}
\label{fig:white-triangles-green-border-in-line}
\end{figure}


Let's start in the beginning, by explaining the fascinating Pythagorean
philosophical origins of the game, the intruiguiging
numerical relationships among its pieces,
and how victory is seen as spiritual enlightenment.

\section{The Philosopher's Game}

It is known that the game goes back at least until the 11th century,
but it may even go back further. As with chess, which has versions
going back to 7th century Persia and India, this board game models
a battle of sorts. Indeed, the {\it Rithmomachia}
name itself means, roughly speaking, ``battle with numbers.''
\index{chess}

Rithmomachia is indeed called the {\it Battle of Numbers} but it's also
called the {\it Philosopher's Game}, and one medieval
author even called it the {\it Pythagorean Game}, believing
(perhaps incorrectly)
that it dated back to the time of the ancient Greeks.
What spiritual connections
did religious medieval scholars believe this game had with 
the world as they saw it?

%\vskip .3in

To monastic scholars, rithmomachia was more than a medieval math
game — it was a ritual of contemplation. The Pythagorean
philosophers during the middle ages believed
\begin{itemize}
\item
numbers are sacred,
\item
harmony is divine, and
\item
through the process of playing rithmomachia, the player
gains a better understanding of the metaphysics of the soul
and the cosmos.
\end{itemize}

\subsection{Pythagorean philosophy}

This board game was popular in many parts of Europe from the 1100s
to the 1600s \cite{Mo01}.
While a reasonable hypothesis, there is no historical
evidence that the origin of the game goes back to
Pythagoras (who lived around 600-500 BCE). What is
known is that the Greek scholar Nicomachus, a Pythagorean who
lived around 60-120 CE, produced a mathematical text that
eventually inspired the rules for rithmomachia. Nicomachus is best known as
the author of {\it Introduction to Arithmetic} and (the musical text)
{\it Manual of Harmonics}.
A semantical word of caution: in the time of Nicomachus,
``arithmetic'' had a different meaning that it does now.
Then, arithmetic was a discipline of study
-- primarily a philosophical discipline concerned with
the essence and underlying principles of numbers,
such as classification and their mystical significance.
Developing skill to perform practical calculations was a separate
discipline of study.
Nicomachus more often viewed numbers as objects having mystical
properties rather than mathematical properties. His work was
translated into Latin and popularized by the Roman scholar Boethius
around 500 CE. More on this important character in our story can be
found in \S \ref{sec:math-edu} below.
\index{Nicomachus}

While not spiritual in the sense of religious doctrine, rithmomachia
is deeply embedded in the philosophical and mathematical
worldview of the Middle Ages, particularly Neoplatonism and Boethian
number theory. It was conceived as more than just a game; it was an
educational tool designed to instill an understanding of cosmic order
through numbers.
Playing it well was seen as both an intellectual achievement and
a spiritual practice.

The starting configuration on the $8\times 16$ board (split between White and
Black halves) reflects a
moral and metaphysical battlefield:
\begin{itemize}
\item
White (Even): Often associated with order, reason, and
light — placed at the bottom, symbolizing the soul's ascent from material
beginnings.
\item
Black (Odd): Associated with dynamism, passion, and will — beginning
at the top, perhaps symbolizing the descent of spiritual forces into the
material.
\end{itemize}
The symmetry and numerical elegance of the starting rows
mimic a cosmic harmony — much like a music scale (which Pythagoreans
also based on ratios). The game’s goal, especially proper victories
(by placing harmonically valued pieces in the enemy’s half), mirrors the soul
achieving harmony within the world.
\index{White/Even}
\index{Black/Odd}


\begin{figure}[h!] 
\centering 
\includegraphics[height=5.5cm,width=5.5cm]{Rithmomachia-game-pieces/ChatGPT-white-square-metaphor.png} 
\caption{\rithmomachia pieces can get metaphysical. \small{Thanks chatGPT:-).}}
\label{fig:ChatGPT-white-square-metaphor} 
\end{figure}


\subsection{Pieces and their spiritual significance}

Rithmomachia features pieces with distinct shapes and numerical values
that convey deep philosophical and spiritual meaning:
\begin{itemize}
\item
The game uses four types of pieces: Circles, Triangles, Squares, and Pyramids.
Each piece has a specific value assigned to it, with black and white
forces having asymmetrical numerical values. This asymmetry was
intentional and reflected medieval philosophical concepts about
numerical harmony and cosmic order.

\item
The numerical values assigned to the pieces aren't arbitrary but represent
mathematical proportions that help explain the fabric of the cosmos.
\end{itemize}
These numerical values and their relationships arose from the
study of Boethian mathematical philosophy, which emphasized the natural
harmony and perfection of number systems.
According to the medieval interpretation,
the world was created out of a knowledge of number and proportion
(Boethius following Nicomachus, Pythagoras, and Plato),
so these values are symbolic of spiritual laws that govern both
the physical and metaphysical realms.
The different shapes of pieces carried the range of Pythagorean symbolism:

\begin{itemize}

\item
Triangles: Represent the trinity in Christian thought and the three-fold
nature of existence (mind/body/spirit).
Their two space movement pattern symbolized
movement between different stages. Triangles stand for the mind or reason,
able to reach farther than the (circular) soul.
These pieces represent developing intellectual capacity.

\vskip.1in
\postersection{Triangular numbers}
... count equally spaced dots inside an equilateral
triangle ($6$, $10$, $15$, \dots).
There is a simple algebraic formula for the $n$th triangular number:
$n(n+1)/2$. The Pythagoreans believed they 
represent growth, virtue, and human
striving. They are associated with moral development.
\postersectionend
\index{triangle}
\index{triangular number}

\item
Squares: Moving three spaces in any orthogonal direction was reminder
of the four elements (earth/wind/fire/water) 
and of the four seasons (summer/fall/winter/spring).
Representing wisdom or even divine authority, their
movement is the longest range — symbolizing fully realized power
or divine insight.

\vskip.1in
\postersection{Square numbers}
... count equally spaced dots inside a square ($4$, $9$, $16$, \dots).
  They represent stability and justice in the cosmos.
Pythagoreans saw squares as the foundation of order.
\postersectionend
\index{square}
\index{square!numbers}

\item
  Pyramids: This is the most complicated and powerful piece.
It also moves like a square, but one that
contains several component sub-pieces of different
values. Indeed, their total value is a sums of squares
-- White's pyramid has value

\[
1^2 + 2^2 + 3^2 + 4^2 + 5^2
+ 6^2 = 1 + 4 + 9 + 16 + 25 + 36 = 91,
\]
while Black's is
\[
4^2+5^2+6^2+7^2+8^2=16+25+36+49+64=190.
\]
It's a spiritually significant piece, representing ascension
and the hierarchy of knowledge. This piece is symbolic of
higher wisdom.
\index{pyramid}
The pyramid piece moves like a square and
contains several component sub-pieces.
\begin{itemize}%[leftmargin=*, label=\textcolor{blue}]
    \item \textbf{White's Pyramid (Value 91)}: Composed of six white sub-pieces with values $1^2, 2^2, 3^2, 4^2, 5^2, 6^2$ (i.e., 1, 4, 9, 16, 25, 36).
    \item \textbf{Black's Pyramid (Value 190)}: Composed of five black sub-pieces with values $4^2, 5^2, 6^2, 7^2, 8^2$ (i.e., 16, 25, 36, 49, 64).
    %\item Sub-piece components can be involved with captures. If a sub-piece is captured, its value is deducted from the pyramid.
\end{itemize}

\vskip.1in
\postersection{Pyramid numbers}
... count equally spaced dots inside a square-based pyramid
($1$, $5$, $14$, \dots). 
There is a formula for the $n$th such number:
$Pyr_n = n(n+1)(2n+1)/6$. The first $9$ pyramid numbers are:
\[
1, 5, 14, 30, 55, 91, 140, 204.
\]
Note that the $7$th pyramid number is the value of the white pyramid.
While the value $190$ of the black
pyramid is not a pyramid number, it is a difference of them:
$190=204-14$. Imagine that this way: if you take the $8\times 8$ pyramid
and remove the $3\times 3$ pyramid from its top, you get the black pyramid. 

The pyramid numbers occured in the works of Nicomachus \cite{Ni00}. 
\postersectionend
\index{pyramid!numbers}


\item
  Circles (or rounds): Symbolized perfection and unity,
  representing the divine whole.
Their limited movement (one space) reflected how fundamental principles
operate within constrained natural laws. The one space they move is the
smallest units — like the circle signifying the
soul's first steps in understanding.
\index{circle}
\index{rounds}


\vskip.1in
\postersection{``Circular'' numbers?}
Unlike triangular numbers, square numbers, and pyramid numbers,
the corresponding definition for the circle doesn't work the same.
{\it If} we defined a ``circular number'' to be those numbers
that count equally spaced dots inside a circle ($1$, $5$, $13$, ...)
then it is {\it much} harder to find a formula for these.
Indeed, the problem of estimating such numbers is
called the {\it Gauss circle problem}, named after Carl
Frederich Guass who first made progress in the early 1800s.
\index{Gauss circle problem}
Indeed, for such numbers, the term ``circular number'' is not
traditionally used. Instead, it is called the ``counting
function for the Gauss circle problem.'' Much better:-)
\index{``circular'' numbers}
\postersectionend
\end{itemize}

The initial arrangement of pieces on the board reflected medieval concepts
of universal order. This orderly arrangement mirrored the medieval concept
of the universe as divinely ordered and mathematically harmonious.
The numerical sequences used in the game followed mathematical progressions
that were considered spiritually significant, such as arithmetic,
geometric, and harmonic sequences. To the medieval students playing
rithmomachia, these victorious patterns were not just game
mechanics but reflections of divine proportions.



\subsection{Victory as spiritual enlightenment}

The various victory conditions in rithmomachia had spiritual significance:
Players could achieve ``common victories'' by capturing certain numbers
of pieces, or ``proper victories'' by arranging pieces in arithmetic,
geometric, or harmonic proportions.

The three proper victories — small, big, and great — require
arranging at least 3 of your pieces in a ``coordinating progression''
explained in detail later. For now, we simply list their names and
spiritual meanings.
\begin{itemize}
\item
Arithmetic progression (e.g., 3, 5, 7): representing justice and equity.
\item
Geometric progression (e.g., 2, 4, 8): symbolizing balance and proportion.
\item
Harmonic progression (e.g., 3, 5, 15): invoking celestial music
and the soul’s harmony. Some medieval authors called
this a ``musical progression''
but harmonic progression is the most commonly used terminology.
\end{itemize}
All of this will be discussed in more detail later. For now, note
that the Pythagoreans believed that achieving a ``great victory''
is not just a tactical feat but a spiritual one as well,
akin to uniting body, soul, and spirit in perfect balance.
These higher victory conditions
represented spiritual enlightenment - the ability to perceive and
create cosmic harmony.



\vskip .2in
\begin{center}
    \adjustimage{max size={1.0\linewidth}{0.9\paperheight}}{Rithmomachia-game-pieces/chatGPT_rithmomachia-logo15c}
\end{center}

\subsection{In mathematics education}
\label{sec:math-edu}


As mentioned, rithmomachia emerged from monastery and cathedral
schools in the 11th
century. One purpose was to teach Boethian number
theory, a cornerstone of the medieval {\it quadrivium} (the four
mathematical arts: arithmetic, geometry, music, and
astronomy) which was the foundation of the educational system at the time.

Ironically, Boethius lived in the final years of the Western Roman Empire.
After its fall, it took hundreds of years to rebuild the
system of mathematical education that was in place.
Monasteries and cathedral schools contributed to the
development of the new educational
system. Of course, after the fall the influence of Islamic scholarship
in the Iberian Peninsula was significant as well. In particular
al-Kwarizimi (780-850 C.E.), who introduced algebra and the Hindu-Arabic
numeral system, had an influence lasting until today,
while the quadrivium curriculum supported by Boethius faded away
by the end of the middle ages.
None-the-less, Boethius' text (the above-mentioned Latin
translation of \cite{Ni00})
was used for hundreds of years after his death -- not just in these
schools but many others in Europe.
Indeed, one of the earliest universities to arise from the dust after the
fall of the Western Roman Empire was the University of Paris, which
coalesced, about 600 years after Boethius' lived,
from an assimilation of the cathedral schools centered around
Notre Dame. 
\index{quadrivium}

Many historians believe the game was first described
(and possibly invented by) a German monk named
Asilo around 1030 C.E.
some believe that this monk was none other than Adalbero,
nominated by Henry III in 1045 to be the highly influential
Bishop Adelbero of W\"urzburg \cite{Mo01}.
Since he was elevated to sainthood by Pope Leo XIII
centuries later in 1883, he is now known as {\it Saint Adalbero}.
\index{Asilo}
\index{Saint Adalbero}

While Asilo is a fascinating character, he was aided in the
development of rithmomachia by a contemporary monk,
Herman the Lame (1013-1054),
who may even be more fascinating. Also known as
Hermannus Contractus and as Blessed Hermann of Reichenau,
Herman the Lame was an 11th-century Benedictine monk and scholar. 
According to wikipedia, he was born with a cleft palate, cerebral palsy, and
possibly had either amyotrophic lateral sclerosis or spinal
muscular atrophy. Since the age of 20, he lived at the
Abbey of Reichenau. In spite of his difficulties in movement
and even in speaking, Herman became literate in several languages,
and contributed scholarship to all four arts of the quadrivium.
Herman is believed to have refined Asilo's original rules,
writing one of the earliest treaties on the game
\cite{MacT}. Despite the fact that these two scholars lived only
about 200 miles from each other, it's not known if they ever met in person.
In fact, some historians suggest Herman the Lame invented rithmomachia
before Asilo \cite{SE11}.
\index{Herman the Lame}
\index{Blessed Hermann of Reichenau}

\begin{figure}[h!] 
\centering
\begin{tabular}{cc}  
  \includegraphics[height=5.5cm,width=5.5cm]{adalbero-of-wurzburg_crop} &
  \includegraphics[height=5.5cm,width=4.5cm]{herman-the-lame} \\
\end{tabular}  
\caption{Who invented \rithmomachia? Asilo as Bishop Adalbero of W\"urzburg (left), Blessed Hermannn of Reichenau  (right)}
\label{fig:asilo-herman} 
\end{figure}
\index{Bishop Asilo}

Who exactly invented the game may be forever
lost to the sands of time,
but it is known that medieval monastic teachers believe that
by playing the game of rithmomachia the student engages with
the mathematical principles that govern the universe.

%\fbox{%
%    \parbox{\linewidth}{%
\postersection{Virtue quotation}
Indeed, de Boissi\`ere emphasized how the \dots

 \begin{quotation}
\dots game appeals to the weaker sides of human nature
that seek diversion and entertainment; at the same time
it leads the player to a more virtuous character.
  
\hskip 1.5in  - Ann Moyer \cite{Mo01}, page 91.
\end{quotation}
%    }%
%}
\postersectionend
\index{de Boissi\`ere, Claude}
\index{Moyer, Ann}
\vskip .1in

Influenced by the works of Boethius (like {\it De institutione arithmetica}),
medieval scholars took a Neoplatonic view of mathematics.
To them, numbers weren't just abstract quantities, they
were fundamental principles governing divine order,
harmony, and the structure of reality. Understanding
numerical relationships was seen as a way to understand
the mind of God and the cosmos.

To help the student improve this numerical understanding,
the teachers encourage students to play rithmomachia.
Consider the rithmomachia rules for capture.
The various capture methods (numbering, addition,
subtraction, multiplication, division, siege) are direct
exercises in calculation and recognizing numerical
relationships (equality, sum, difference, product, ratio,
spatial enclosure). They force the player to constantly
analyze the numerical properties of the pieces and their
spatial relationships.

Rithmomachia was a widely used instructional
tool to teach arithmetic in monasteries, universities and cathedral-based
grade schools in the later half of the middle ages.
Unfortunately, it appears that the
numerous versions of the game, with variations from
one kingdom to the next, may have contributed to its slow acceptance,
while chess was becoming more popular. Ultimately, the decline of rithmomachia
was more a result of wide-spread change in the curriculum for mathematics
education in universities.
Indeed, around the 1600s, university
curriculum we decided more by subject matter experts
and less by authority figures (see chapter 5 in \cite{Mo01}).
Without the need to teach Boethian number theory
and Pythagorean philosophy, there was no reason for teachers to
promote the virtues of rithmomachia to their students.


\vskip .1in


\begin{figure}[h!] 
\centering 
\includegraphics[height=2.0cm,width=14.0cm]{Rithmomachia-game-pieces/black-circles-blue-border-in-line} 
\caption{Rithmomachia Black circles pieces in a line}
\label{fig:black-circles-blue-border-in-line} 
\end{figure}

\section{The battlefield}
\label{sec:board}


We shall adopt the {\it game board} as a $8\times 16$ board.
It may be checkered like a chess board or plain like a go board.
Following military terminology first introduced in the 1500s, 
there are eight {\it files} and 16 {\it ranks}.
\index{file (on gameboard)}
\index{rank (on gameboard)}
On a rithmomachia board, the ranks are lettered (a-p)
while the files are numbered (1-8).
The ranks {\tt a-h} comprise {\it Even/White territory}, while the
ranks {\tt i-p} comprise {\it Odd/Black territory}.

\vskip .1in

{\small{
\begin{center}
\begin{tabular}{||c|c|c|c|c|c|c|c|c|c|c|c|c|c|c|c||c} \hline \hline
\quad & \quad & \quad  & \quad & \quad  & \quad & \quad & \quad & \quad & \quad  & \quad & \quad  & \quad & \quad & \quad & \quad & 1\\ \hline
\quad & \quad & \quad  & \quad & \quad  & \quad & \quad & \quad & \quad & \quad  & \quad & \quad  & \quad & \quad & \quad & \quad & 2\\ \hline
\quad & \quad & \quad  & \quad & \quad  & \quad & \quad & \quad & \quad & \quad  & \quad & \quad  & \quad & \quad & \quad & \quad & 3\\ \hline
\quad & \quad & \quad  & \quad & \quad  & \quad & \quad & \quad & \quad & \quad  & \quad & \quad  & \quad & \quad & \quad & \quad & 4\\ \hline
\quad & \quad & \quad  & \quad & \quad  & \quad & \quad & \quad & \quad & \quad  & \quad & \quad  & \quad & \quad & \quad & \quad & 5\\ \hline
\quad & \quad & \quad  & \quad & \quad  & \quad & \quad & \quad & \quad & \quad  & \quad & \quad  & \quad & \quad & \quad & \quad & 6\\ \hline
\quad & \quad & \quad  & \quad & \quad  & \quad & \quad & \quad & \quad & \quad  & \quad & \quad  & \quad & \quad & \quad & \quad & 7\\ \hline
\quad & \quad & \quad  & \quad & \quad  & \quad & \quad & \quad & \quad & \quad  & \quad & \quad  & \quad & \quad & \quad & \quad & 8\\ \hline \hline
a    &   b    & c     &   d   &   e    &   f   &    g   &   h   &   i   &   j    &   k   &   m    &   m   &   n  &    o  &   p  &  \\ 
\end{tabular}
\end{center}
}}
\index{battlefield}
\index{territory!Odd}
\index{territory!Even}

\vskip .2in

%\newpage
%{\small{
%%%%%% box
\postersection{Coordinate notation}
An alternative to the algebraic notation is the
coordinate-based notation shown in Figures \ref{fig:latex-table} and
\ref{fig:rithmomachia-initial-state-with-pyramid3}.
Imagine the rithmomachia board as being a grid in the $1$st quadrant
of the Cartesian $xy$ plane ranging from $0$ to $15$ in the $x$-direction
and $0$ to $7$ in the $y$-direction. 
%(In this paper, the rithmomachia board
%is usually presented horizontally for typographical reasons.)
\index{notation!coordinate}
We order the coordinates in this grid as though they were entries
in an $8\times 16$ matrix, with the count starting at $0$ (instead
of starting at $1$).
\newline
{\it Example}: A White circle of value $6$ located in the
upper left-hand corner is denoted
$C6{:}(0,0)$ or, in exponential notation, $C^6(0,0)$.
Another option is to use the circle symbol font:
\circlet{\,}$^6(0,0)$.
\postersectionend
%}}
\index{coordinate notation}

A schematic for how the armies are initially arranged on the battlefield
before battle commences is given below. 

\vskip .1in
\postersection{Initial Setup}
\vskip .1in
\begin{center}
\textbf{White/Even's Back Ranks:}
\newline
{\small{
\renewcommand{\arraystretch}{1.5}
\begin{tabular}{c|cccccccc} \hline
  Rank d & xxxx &  xxxx &   C008 &  C006 &  C004 &  C002 &  xxxx &  xxxx \\
  Rank c & T081 &  T072 &  C064 &  C036 &  C016 &  C004 &  T006  & T009 \\
  Rank b & S153 &  P091 &  T049 &  T042 &  T020 &  T025 &  S045 &  S015 \\
  Rank a & S289 &  S169  &  xxxx &  xxxx &  xxxx &   xxxx  &  S081 & S025\\ \hline
\end{tabular}
}}
\end{center}
\textit{(P denotes pyramid, T triangle, C circle, S square)}

\vskip .1in
Black's setup mirrors this on the opposite side with its own set of values.
\vskip .1in

\begin{center}
\textbf{Black/Odd's Back Ranks:}
\newline
{\small{
\renewcommand{\arraystretch}{1.5}
\begin{tabular}{c|cccccccc} \hline
  Rank m &  xxxx &  xxxx &  c009 &  c007 &  c005 &  c003 &  xxxx &  xxxx \\
  Rank n &   t100 &  t090 &  c081 &  c049 &  c025 &  c009 &   t012 &  t016 \\
  Rank o &   p190 &  s120 &  t064 &  t056 &  t030 &  t036 &   s066 &  s028 \\
  Rank p &   s361 &  s225 &  xxxx &  xxxx &  xxxx &  xxxx &   s121 &  s049 \\ \hline
\end{tabular}
}}
\end{center}
\textit{(p denotes pyramid, t triangle, c circle, s square)}
\postersectionend

\vskip .1in

The symbols for the pieces will be described below.

\section{The armies}
\label{sec:pieces0}

Initially, both Black/Odd and White/Even are provided with a company of
twenty-four pieces, each divided up into $8$ platoon-like sets
(to borrow the battlefield lingo):
\index{platoon}
\index{company}
\begin{itemize}
\item
eight {\it circles} (denoted $C$ or {\circlet} for a white circle,
$c$ or {\circletfill} for a black circle,
when the color must be specified),
\item
eight {\it triangles} (denoted $T$ or {\trianglepa} for white, $t$ or
{\trianglepafill} for black), and
\item
eight {\it squares} (denoted $S$ or {\squad} for white, or $s$ or
{\squadfill} for black),
\item
For each side, exactly one of the square pieces is distinguished
and called a {\it pyramid}, denoted $P$ = {\rhombus} (or $p$ = \rhombusfill).
This type of piece is more powerful, and more complicated, than the
others but will be explained in complete detail later.
\end{itemize}
The position of the pieces during a particular moment in a game
will be called the (current) {\it state} or the {\it game state}.
\index{piece!circle}
\index{circle!piece}
\index{piece!triangle}
\index{triangle!piece}
\index{piece!square}
\index{square!piece}
\index{piece!pyramid}
\index{state}
\index{game state}

In chess, the files are named after the back-rank pieces initially
occupying them:
the queen rook file, the queen knight file, and so on. Likewise, in
rithmomachia, we can name the file after the {\it most valuable} piece
initially occupying it:

\begin{center}
{\tt 289},\quad {\tt 169},\quad  {\tt 49},\quad  {\tt 42},\quad {\bf 56},\quad {\bf 64},\quad {\bf 225},\quad {\bf 361}.
\end{center}
Note: the number is in {\bf bold} font if it is a Black piece,
{\tt Courier} font if it is a White piece.
Half of these are White/Even and half are Black/Odd!

%\begin{center}
%\begin{tabular}{|ccc|} \hline
%The {\tt $289$ file} & \dots & \dots \\ \hline
%The {\tt $169$ file} &  \dots & \dots \\ \hline
%The {\tt $49$ file} &  \dots & \dots \\ \hline
%The {\tt $42$ file} &  \dots & \dots \\ \hline
% \dots & \dots &  The {$\mathbf{56}$ file}\\ \hline
% \dots & \dots &  The {$\mathbf{64}$ file}\\ \hline
% \dots & \dots   &  The {$\mathbf{225}$ file}\\ \hline
% \dots & \dots   &  The {$\mathbf{361}$ file}\\ \hline
%\end{tabular}
%\end{center}

Some other fun facts about the file numbers:

\begin{enumerate}
\item
$289$ is
\begin{itemize}
\item
  %(i)
  a perfect square ($=17^2$),
\item
  %(ii)
  a sum of perfect cubes ($1^3+2^3+4^3+6^3$), and
\item
  %(iii)
  is the sum of the first 5 numbers raised to themselves
($0^0+1^1+2^2+3^3+4^4$).
\end{itemize}

\item
$169$ is
\begin{itemize}
\item
  %(i)
  a perfect square ($=13^2$),
\item  
  %(ii)
  a centered hexagonal number.
\end{itemize}

\index{centered hexagonal number}
\index{hexagonal number!centered}

\begin{center}
    \adjustimage{max size={0.9\linewidth}{0.9\paperheight}}{centered_hexagonal_numbers.jpg}
\end{center}
%
%\begin{figure}[h] 
%\centering 
%\includegraphics[height=4.5cm,width=12.5cm]{centered_hexagonal_numbers.jpg} 
%\caption{
A centered hexagonal number $H_n$ counts the number of
  equally spaced dots inside a hexagon. There is a formula for them:
  $H_n = n^3-(n-1)^3$.
 %}
%\label{fig:centered_hexagonal_numbers}
%\end{figure}

\item
$49$ is
\begin{itemize}
\item
  %(i)
  a perfect square ($=7^2$),
\item
  %(ii)
  the smallest triple of three squares in arithmetic succession
  $(1,25,49)$.
\end{itemize}
\item
$42$ is 
\begin{itemize}
\item
  not a perfect square, but ``near'' one: it's the product of consecutive
  integers, $6\cdot 7$, one of which is perfect!
\item
  Anyone who has heard of Douglas Adams' book
{\it The Hitchhiker's Guide to the Galaxy} knows all that is needed to know
about the number 42!
\end{itemize}

\item
$56$ is 
\begin{itemize}
\item
  %(i)
  the sum of the first six triangular numbers (so counts the equally spaced dots
  in a regular tetrahedron whose base has length $6$),
\item  
  %(ii)
  twice a perfect number ($=2\cdot 28$).
\end{itemize}
Plutarch \cite{Pl00} reported the Pythagoreans associated
a regular polygon of $56$ sides with Typhon
  (one of the deadliest creatures in Greek mythology).
\index{Typhon}

\item
$64$ is 
\begin{itemize}
\item
  %(i)
  a perfect square,
\item
  %(ii)
  a perfect cube,
\item
  % (iii)
  a centered triangular number.
\end{itemize}


\begin{center}
    \adjustimage{max size={0.9\linewidth}{0.9\paperheight}}{centered_triangular_numbers.jpg}
\end{center}
%\begin{figure}[h] 
%\centering 
%\includegraphics[height=4.5cm,width=12.5cm]{centered_triangular_numbers.jpg} 
  %\caption{
The centered triangular numbers are constructed by placing a
dot in the center and all its other dots surrounding the center in successive
equilateral triangular layers.
%}
%\label{fig:centered_triangular_numbers}
%\end{figure}

\vskip.1in
%%%%%%%%%%%% box
\postersection{centered triangular number}
%hese numbers are very similar to a triangular
%numbers, but there is a subtle distinction.
The triangular numbers are constructed inductively
in a different way:
by adding a slightly larger base (depicted as dots along a line segment).

If $C_n$ is the $n$th centered triangular number and $T_n$ is
the $n$th triangular number then
    there is a simple formula relating them: $C_n = 3T_n+1$.
\postersectionend
\index{centered triangular number}
\index{triangular number!centered}

\item
Note $225$ is
(i) a perfect square ($=15^2$), (ii) the square of a triangular number.

\item
 Note $361$ is (i) a perfect square ($=19^2$), (ii) a perfect cube,
  (iii) a centered triangular number.
\end{enumerate}


\section{The rules}
\label{sec:rules}

%\begin{quotation}

This history and philosophical underpinnings of the game presented
in the previous section makes it clear that there were several
design decisions made by the inventor (or inventors) or
rithmomachia that are very spiritual in nature.
Do today's players believe that by playing rithmomachia they
become closer to God? Maybe they don't, but it seems the game
was designed to be entertaining to those who did.


Of all the many versions of rithmomachia, this
section only records the rules and setup for the version described by
Claude de Boissi\`ere in 1556.
Not a lot is known about Claude de Boissi\`ere,
  also known as Claudius Buxerius. Born around 1500 near Grenoble, he
  was an influential French scholar active in promoting mathematical education.
  He wrote and published texts on various subjects,
  including poetry, music, and astronomy. His most renown work is
\cite{dB56}, translated by Richards in \cite{Ri46}.
\index{Claudius Buxerius}
\index{de Boissi\`ere, Claude}
%\end{quotation}

%\tableofcontents

%\newpage

\vskip .3in

%The focus of the discussion here is on the initial setup and rules for one
%relatively common version of rithmomachia. However, the
%fascinating history of the game will be mostly ignored.

The current section delves into the fundamental rules governing
the game of rithmomachia, based on the version described by
de Boissi\`ere.  Sections \S\S \ref{sec:setup}-\ref{sec:pieces}
detail the game's physical
components, covering the specifics of the board and the various
pieces involved.  Following this, \S \ref{sec:setup-pieces}
explains the initial arrangement of these pieces on the board.
The subsequent sub-sections will describe the distinct ways in which each
piece can move (\S \ref{sec:move-pieces}) then the various methods for
capturing an opponent's pieces (\S \ref{sec:capture-pieces}).
Finally, \S \ref{sec:winning} describes the conditions that determine
victory in a rithmomachia game.


\subsection{Initial setup}
\label{sec:setup}

Like chess or checkers,
rithmomachia is a board game played between two people, called
Black (or Odd) and White (or Even). These different
names (White and Even, for example) will be used interchangably.
The two players alternate turns, each turn consisting of a move of a friendly
piece and possibly captures of enemy pieces.
According to de Boissi\`ere, White/Even has the first
turn. Other versions of rithmomachia state Black/Odd moves first.
Here, we give an overview of the game play process.

To describe these
rules more clearly, we distinguish between a move
and a capture (where no friendly piece moves
but an enemy piece is taken off the board).
Details of how the pieces move will be described in
\S \ref{sec:move-pieces} and how they capture will be described
in \S \ref{sec:capture-pieces}.


\vskip.1in
\postersection{What is a player's turn?}
A {\it turn} consists of three stages.
\index{victory!conventions}

\begin{itemize}
\item
Inital captures, if any.

For instance, such captures could be the result of the opponent's previous
move. In any case, more than one capture is allowed.

\item
A legal move.

Exactly one move is made per turn, onto an unoccupied position on the board.
This space could be the one occupied by an enemy piece just captured.

\item
Final captures, if any.

This could be a capture that was made possible by move just made.
\end{itemize}
\postersectionend
\index{turn}

If it is a player's turn and no legal move can be made, then that
player {\it loses} the game. This is in contrast to chess,
where a stalemate results in a {\it draw}.
\index{stalemate}

There are no draws in this version of rithmomachia.

\subsection{Board and pieces}
\label{sec:pieces}


Each piece has a value. These values are tabulated in Figures
\ref{fig:whitepieces} and \ref{fig:blackpieces}.

\begin{figure}[h!] 
\begin{center}
\renewcommand{\arraystretch}{1.5}
\begin{tabular}{c|cccc|c} 
piece & & & & & formula \\ \hline
circles \quad               &  9                  &   7  &   5 &   3        & $x$ \\
circles {\circletfill}       &  81               &  49 &  25 &  9        & $x^2$ \\ \hline
triangles \quad               & 90               &  56 &  30 &  12        & $x(x+1)$ \\
triangles {\trianglepafill}  & 100                &  64  &  36 & 16        & $(x+1)^2$ \\ \hline
squares \quad               & {\bf 190}\rhombusfill  &  120  & 66 & 28 & $(x+1)(2x+1)$ \\
squares \squadfill \ \      & 361                 &  225 & 121 & 49         & $(2x+1)^2$ \\ \hline
\end{tabular}
\caption{\textcolor{gray}{Black pieces and their values, which total $1752$.
    The black pyramid is the square $190$. You can imagine stretching this table
    over the Black side of the board to see how the pieces are initially placed.
    Note all the values of the back rank of the Black/Odd pieces are odd numbers,
while all the values in the second-to-the-last rank of pieces are {\it even} integers.}}
\label{fig:blackpieces}
\end{center}
\end{figure}


%The 24 black pieces are tabulated, along with their values, in Figure \ref{fig:blackpieces}.
%These 24 pieces are sometimes called the {\it Odd pieces},
%even though the values of the pieces are not necessarily odd integers.
%  \index{Odd pieces}
  
%Likewise, the 24 white pieces (sometimes called the
%{\it Evens}) are in Figure \ref{fig:whitepieces} \cite{Ne98}.
%Most of the white and black pieces have distinct values, but
%repetitions are noted in the table caption.
%\index{Even pieces}

\begin{figure}[h!] 
\begin{center}
\renewcommand{\arraystretch}{1.5}
\begin{tabular}{c|cccc|c}
piece                 &     & & & & formula \\ \hline
circles               &  8  &   6   &   4              &  2   & $y$ \\
circles  {\circlet}   &  64  &  36   &  16             &  4   & $y^2$ \\ \hline
triangles              &  72 &  42   &  20              &  6   & $y(y+1)$ \\
triangles  \trianglepa  &  81 &  49  &  25              &  9  & $(x+1)^2$ \\ \hline
squares                & 153  &  15  &  {\bf 91}\rhombus & 45 & $(y+1)(2y+1)$ \\
squares {\squad}       & 289  &  169  & 81              & 25  & $(2y+1)^2$ \\ \hline
\end{tabular}
\caption{\textcolor{gray}{White pieces and their values, which total $1312$.
    The White pyramid is the square $91$. Note all the values of the back rank of the
    Even/White pieces are {\it odd} integers.}}
\label{fig:whitepieces}
\end{center}
\end{figure}


Both
sides have piece values $9$, $16$, $25$, $36$, $49$, $64$, $81$ in common.
These values, all square numbers, are visible in the diagram of
the initial board setup in Figure \ref{fig:latex-table} and
in Figure \ref{fig:rithmomachia-initial-state-with-pyramid3}.
%\newpages
\index{piece!value(s)}


\renewcommand{\arraystretch}{1.5}
\begin{center}
\begin{figure}[h!]
\centering
\begin{tabular}{c|c|c|c|c|c|c|c|c|}\hline
\textcolor{gray}{$p$} &  \squadfill$_1$49 & \squadfill121 &\  & \   & \   & \   & \squadfill225  & \squadfill361  \\ \hline
\textcolor{gray}{$o$} &  \squadfill28  & \squadfill66 & \trianglepafill36 & \trianglepafill30 & \trianglepafill56 & \trianglepafill64 & {\squadfill}120 & \rhombusfill190   \\ \hline
\textcolor{gray}{$n$} &  \trianglepafill16  & \trianglepafill12 & \circletfill$_2$9 & \circletfill25 & \circletfill49 & \circletfill81 & \trianglepafill90 & \trianglepafill100  \\ \hline
\textcolor{gray}{$m$} &   \ & \ & \circletfill3 & \circletfill5 & \circletfill7 & \circletfill$_1$9 & \ & \  \\ \hline
\textcolor{gray}{$l$} &   \ & \ & \ & \ & \ & \ & \ & \  \\ \hline
\textcolor{gray}{$k$} &   \ & \ & \ & \ & \ & \ & \ & \  \\ \hline
\textcolor{gray}{$j$} &   \ & \ & \ & \ & \ & \ & \ & \  \\ \hline
\textcolor{gray}{$i$} &   \ & \ & \ & \ & \ & \ & \ & \  \\ \hline
\textcolor{gray}{$h$} &   \ & \ & \ & \ & \ & \ & \ & \  \\ \hline
\textcolor{gray}{$g$} &   \ & \ & \ & \ & \ & \ & \ & \  \\ \hline
\textcolor{gray}{$f$} &   \ & \ & \ & \ & \ & \ & \ & \  \\ \hline
\textcolor{gray}{$e$} &   \ & \ & \ & \ & \ & \ & \ & \  \\ \hline
\textcolor{gray}{$d$} &   \ & \ & \circlet8 & \circlet6 & \circlet$_1$4 & \circlet2 & \ & \  \\ \hline
\textcolor{gray}{$c$} &   \trianglepa 81 & {\trianglepa}72 & \circlet64 & \circlet36 & \circlet16 & \circlet$_2$4 &  \trianglepa 6 & {\trianglepa}9  \\ \hline
\textcolor{gray}{$b$} &   {\squad}153 & {\rhombus}91  & {\trianglepa}49 & {\trianglepa}42 & {\trianglepa}20 & {\trianglepa}25 & {\squad}45 & {\squad}15  \\ \hline
\textcolor{gray}{$a$} &  {\squad}289  & {\squad}169 & \ & \ & \ & \ & {\squad}81 & {\squad}$_1$25  \\ \hline
\ \ &  \textcolor{gray}{$1$} & \textcolor{gray}{$2$} & \textcolor{gray}{$3$} & \textcolor{gray}{$4$} & \textcolor{gray}{$5$} & \textcolor{gray}{$6$} & \textcolor{gray}{$7$} & \textcolor{gray}{$8$}  \\ 
\end{tabular}
\newline
\caption{{\small{
      \textcolor{gray}{The initial state of rithmomachia, with \rhombus (at $b2$, in algebraic notation) and \rhombusfill
        (at $o8$) as the pyramids.
        Note there are duplicate pieces (eg, two White circles of value $4$ denoted \circlet$_1$4 and \circlet$_2$4),
        distinguished using subscripts when necessarly. All these pieces, including the duplicates, will be listed in
      \S \ref{sec:enum-pieces} below.}
}}}
\label{fig:latex-table}
\end{figure}
\end{center}

%\begin{figure}[h] 
%\centering 
%\includegraphics[height=7.5cm,width=13.5cm]{rithmomachia-initial-state.jpg} 
%\caption{Rithmomachia inital position, where $91$ is the value of the white pyramid and $190$ is that of the black pyramid}
%\label{fig:rithmomachia-initial-state} 
%\end{figure}

These are the pieces and their values. 
Next, we learn about some of the numerical relationships among
a piece and its neighbors in their initial placement on the board.


\subsection{Setting up the pieces}
\label{sec:setup-pieces}

%Fortunately, the initial setup of the pieces is more standardized across
%the various versions of rithmomachia.

The chess-like notation for the rows and columns used in
Figure \ref{fig:latex-table} is borrowed from the book
Codenotti-Resta \cite{CR22}.
Unlike chess, there is no standardized notation yet for rithmomachia
but this Codenotti-Resta notation is a place to good start.
%In chess, the ranks are numbered (1-8) while the files are lettered (a-h).
%In rithmomachia, the ranks are lettered (a-p) while the files
%are numbered (1-8).

It's natural to wonder how these piece values were selected.
%There is an interesting, if somewhat numerological,
%answer but for more details, please see the references in the papers
%included in the bibliography below.
%On the other hand, we do want to point out
By design, the game contains some interesting patterns in the
initial setup in Figure \ref{fig:rithmomachia-initial-state-with-pyramid3}.
The arithmetical patterns in the setup of the pieces is displayed in the
next two tables, presented next.

Start with the top row of pieces in the initial position -- namely, $8, 6, 4, 2$
for White and $9, 7, 5, 3$ for Black. For simplicity, we'll denote this row by
$a,b,c,d$. The following lemma basically says
very other piece value is easily determined from $a,b,c,d$.

\begin{lemma}
{\rm
Each initial piece value can be computed using 
a polynomial in the values of the forward-most piece values.
}
\end{lemma}  

\begin{proof}
This follows from Figures \ref{fig:blackpieces} and \ref{fig:whitepieces},
but will also be established in the tables below.
\end{proof}


To get the value of the pieces in the second row, compute these
(mostly) squares:

\renewcommand{\arraystretch}{1.5}
\begin{center}
%\begin{figure}[h]
%\centering
\begin{tabular}{|c|c|c|c|c|c|c|c|}%\hline
 &\  & \   & \   & \   &  &\  & \     \\ \hline
 & \   & $\quad a\quad $  & $\quad b\quad $  &  $\quad c\quad $  &  $\quad d\quad $ & \   & \   \\ \hline
$(a+1)^2$ & $a+a^2$  & $a^2$  & $b^2$  & $c^2$  &  $d^2$ & $d+d^2$   & $(d+1)^2$   \\ \hline
\dots &\dots  & \   & \   & \   &  &\dots & \dots     \\ 
\end{tabular}
%\caption{The second row circles are squares of those on the top row.}
%\label{fig:second-row-circles}
%\end{figure}

\vskip .1in
{\bf Table: top two rows}
\end{center}

Rename these values of the 2nd row as $e, f, \dots, \ell$ for simplicity.
The numerical relationships in the third row are also easy to describe:

\renewcommand{\arraystretch}{1.5}
\begin{center}
%\begin{figure}[h]
%\centering
\begin{tabular}{|c|c|c|c|c|c|c|c|}%\hline
\quad &\quad  & \   & \   & \   &  &\quad  & \quad     \\ \hline
 & \   & $\quad a\quad $  & $\quad b\quad $  & $\quad c\quad $  &  $\quad d\quad $ & \   & \   \\ \hline
$e$ & $f$  & $g$  & $h$  & $i$  &  $j$ & $k$   & $\ell$   \\  \hline
$m$ & $n$  & $o$  & $p$  & $q$  &  $r$ & $s$   & $t$   \\ \hline
$u$ & $v$  & 0   & 0   & 0   & 0  & $w$  & $x$    \\ \hline
\end{tabular}
%\caption{The second row circles are squares of those on the top row.}
%\label{fig:second-row-circles}
%\end{figure}
\vskip .1in
{\bf Table: all four rows}
\end{center}
We know already,

\[
e=(a+1)^2,\quad  f=a^2+a, \quad  g= a^2,\quad  h= b^2,
\]
\[
i = c^2,\quad  j = d^2,\quad  k = d^2+d,\quad \ell = (d+1)^2,
\]
but the remaining expressions are easy to compute as well:

\[
m = (a+1)(2a+1), \quad n = (b+1)(2b+1),\quad o = (b+1)^2,\quad p = b^2+b,
\]
\[
\quad q = c^2+c,\quad r = (c+1)^2,\quad s = (c+1)(2c+1), \quad t = (d+1)(2d+1).
\]
We leave the values of the pieces at the bottom to the reader.
%\[
%u = (2a+1)^2,\quad v = (2b+1)^2,\quad w = (2c+1)^2,\quad x = (2d+1)^2.
%\]
\begin{question}
\label{question:1}  
Can you find the arithmetical relationships that determine the numbers
$u,v,w,x$ in the bottom row of the above table?
Here, for example, $u=289, \dots, x = 25$
if you are White and $u=361, \dots, x = 49$ if you are Black.
What algebraic formula determines them from the other entries?
(Hint: they are all perfect squares! The answer is at the end,
just before the references.) 
\end{question} 

The symmetry and arithmetical properties of the setup of
the rithmomachia pieces is remarkably beautiful, isn't it?
The interested reader can find another discussion of the numerical
relationships between the
piece values on pages 74-75 of \cite{SE11}.


\vskip.1in
{\small{
%%%%%%%%%%%% box
\postersection{Algebraic notation}
Each of the $128$ {\it coordinate positions}
on the board can be described
by a pair, which we
juxtapose for simplicity: $a1$, $a2$, \dots, $p8$.
\index{coordinates}
The shape of the pieces can be described by its first letter:
$C, P, S, T$ for the Even/White pieces and $c, p, s, t$ for the Odd/Black
ones. The value of a piece can be described by a list of positive
integers\footnote{Except the case of the multi-valued pyramid piece
  and its subpieces, which is described in detail in
  \S \S \ref{sec:pyramid1}, \ref{sec:pyramid2} below,
  all these lists have only one element,
  na,ely the value of the piece itself.}.
\newline
{\it Example}: A Black circle of value $3$ at location $h5$ is denoted
$c3h5$, but we sometimes use exponential notation, $c^3h5$.
Another option is to use the filled circle symbol font:
\circletfill{\,}$^3h5$.
\postersectionend
}}
\index{algebraic notation}


The pieces are all listed in \S \ref{sec:enum-pieces} below, including all duplicates and
component subpieces of the pyramid pieces.

\begin{figure}[h!] 
\centering 
\includegraphics[height=7.5cm,width=13.5cm]{rithmomachia-initial-state-with-pyramid3.jpg} 
\caption{\textcolor{gray}{Rithmomachia initial position, where
    $91$ (at $(1,1)$ in coordinate notation) is the value of the white
    pyramid and $190$ (at $(7, 14)$ in coordinate notation) is that of the black pyramid.
    This diagram matches that of de Boissi\`ere's, as shown in
\cite{Ri46}, Figure 8.} }
\label{fig:rithmomachia-initial-state-with-pyramid3} 
\end{figure} 



\subsubsection{Enumerating the pieces}
\label{sec:enum-pieces}

If we count the White pieces, including the duplicates, the pyramid,
and all its sub-pieces, we get $30$:

\[
S^{289},\ \  S^{169}, \ \ S^{153},\ \  P^{91},\ \  S^{81}, \ \ T^{81},\ \  T^{72},\ \  C^{64},\ \  T^{49},\ \  S^{45}, \ \ T^{42},
\]
\[
C^{36},\ \  S^{36}, \ \ S_1^{25},\ \  S_2^{25}, T^{25},\ \  \ \ T^{20},\ \  C^{16},\ \  S^{16},\ \  S^{15},\ \  T^{9}, \ \ S^{9},\ \  C^{8},
\]
\[
T^{6},\ \  C^{6},\ \  C_1^{4},\ \  C_2^{4},\ \  S^{4},\ \  C^{2}, \ \ S.
\]
(We use subscripts to notationally distinguish the two White circles of
value $4$, and the two White square of value $25$.)
On the other hand, the total number of distinct values used by the Even/White pieces is, due to the duplicates,
only $20$.


\vskip .1in
\postersection{Duplicate values for the Even/White pieces}
White has three 4s, two 6s, two 9s, two 16s,
two 25s, two 36s, and two 81s.
Other piece values are unique to the color.
 Here is a histogram plot:
 \index{duplicate values}

\[
\begin{array}{l|ccccccc}
                    & \circ &      &        &       &       &       &      \\
                    & \circ & \circ & \circ & \circ & \circ & \circ & \circ \\
                    & \circ & \circ & \circ & \circ & \circ & \circ & \circ \\  \hline
{\rm piece\ value}  &  4    &  6    & 9     & 16    & 25    &  36   &  81 \\ 
\end{array}
\]
\postersectionend
\vskip .1in

A way to avoid using subscripts in the notation, yet still eliminate ambiguities caused by the repetition of
piece values, is to include the initial coordinate of the piece:
\begin{itemize}
\item  
There are $14$ {\squad}{'}s (counting duplicates and the pyramid as a square):

\[
S^{289}a1,\ \  S^{169}a2, \ \ S^{153}b1,\ \  P^{91}b2,\ \  S^{81}a7,\ \  S^{45}b7,
\]
\[
S^{36}b2, \ \ S^{25}a8, \ \ S^{25}b2, \ \ S^{16}b2, \ \ S^{15}b8, \ \ S^9b2,\ \  S^4b2,\ \  S^1b2.
\]
\item
There are $8$ $\triangle$s:

\[
T^{81}c1,\ \  T^{72}c2,\ \  T^{49}b3, \ \ T^{42}b4,\ \ T^{25}b6,\ \  T^{20}b5,\ \  T^{9}c8,\ \  T^6c7.
\]
\item
There are $8$ of the \circlet:

\[
C^{64}c3,\ \  C^{36}c4,\ \  C^{16}c5,\ \  C^{8}d3,\ \  C^{6}d4, \ \ C^{4}c6,\ \  C^{4}d5,\ \  C^{2}d6.
\]
\end{itemize}


If we count the Black pieces, including the duplicates, the pyramid,
and all its sub-pieces, we get $29$:
%Using subscripts to eliminate the ambiguities of the Odd/Black pieces:

\begin{itemize}
\item  
There are $14$ {\squad}{'}s

\begin{center}
\squadfill{\,361},\ \  \squadfill{\,225}, \ \ \squadfill{\,190},\ \  \squadfill{\,121},\ \  \squadfill{\,120},\ \  \squadfill{\,66}, \ \ \squadfill{\,64}, \ \ \\
\squadfill$_1${\,49}, \ \ \squadfill$_2${\,49},  \ \ \squadfill{\,36},\ \ \squadfill{\,28},\ \ \squadfill{\,25},\ \ \squadfill{\,16},\\

\end{center}

Note there are $2$ black squares with value $49$.

\item
  There are $8$ of the \trianglepafill{'}s:

\begin{center}
\trianglepafill{\,100},\ \ \trianglepafill{\,90}, \ \ \trianglepafill{\,64}, \ \ \trianglepafill{\,56},
\trianglepafill{\,36}, \ \ \trianglepafill{\,30},\ \ \trianglepafill{\,16}, \ \ \trianglepafill{\,12},
\end{center}

\item
  There are $8$ of the \circletfill{'}s:
  
\begin{center}
  \circletfill{\,81},\ \  \circletfill{\,49}, \ \ \circletfill{\,25},\ \ \circletfill$_1${\,9},\ \
  \circletfill$_2${\,9},\ \ \circletfill{\,7}, \ \ \circletfill{\,5}, \ \ \circletfill{\,3}.
\end{center}

Note there are $2$ black circles with value $9$.
\end{itemize}
On the other hand, the total number of distinct values used by the Black/Odd pieces is, due to the duplicates,
only $22$.


\vskip .1in
\postersection{Duplicate values for the Odd/Black pieces}
Black has two 9s, two 16s,
two 25s, two 36s, three 49s, and two 64s.
Other piece values are unique to the color.
Here's a histogram plot:
 \index{duplicate values!Black/Odd}
 
\[
\begin{array}{l|cccccc}
                    &       &      &        &       & \circ &        \\
                    & \circ & \circ & \circ & \circ & \circ & \circ  \\
                    & \circ & \circ & \circ & \circ & \circ & \circ  \\  \hline
{\rm piece\ value}  &  9    &  16    & 25     & 36  & 49    &  64    \\ 
\end{array}
\]
\postersectionend

%\begin{question}
%Can you list the initial coordinates of all 29 Black pieces?
%\end{question}



\begin{figure}[h!] 
\centering 
\includegraphics[height=4.0cm,width=14.0cm]{Rithmomachia-game-pieces/black-squares-red-border-in-2lines} 
\caption{Rithmomachia Black square pieces}
\label{fig:black-squares-red-border-in-2lines} 
\end{figure}





\subsection{How the pieces move}
\label{sec:move-pieces}

The rules for how pieces move are pretty simple.

\subsubsection{Movement notation}

During a player's turn, only one move (as described in this section)
can be made, but multiple captures are possible, and these can be performed
before and/or after the move. Captures are described in a later section.
\index{piece!moves}

\vskip.1in
%%%%%%%%% box
\postersection{Move notation}
Based on notation used for chess,
one simple coordinate-based notation for a move 
would be something like this:

%{\small{
\begin{center}
{\tt  move = (shape+value):(starting coord.)-(ending coord.)}.
\end{center}
%}}
For example, in coordinate notation
$C2{:}(5,3)-(6,3)$ describes moving the white circle piece with
value $2$ from the position/coordinate $(5, 3)$ to $(6, 3)$.
In the more compact algebraic notation, this move would be $C2d6{-}d7$,
or $C^2d6d7$, or even $\bigcirc^2d6d7$.
\postersectionend

For another example of a move in coordinate notation, 
see Figure \ref{fig:rithmomachia-after-S81-60-30-with-pyramid2}
for the move $S^{81}{:}(6,0){-}(3,0)$.
That same move in the more chess-like notation of
Figure \ref{fig:latex-table},
would be written as $S^{81}a7{-}a4$ or $\square^{81}a7a4$.
See the caption in Figure \ref{fig:rithmomachia-initial-state-with-pyramid3}
for more coordinate examples.
%The algebraic notation used here is the same
%as that used in the book by Codenotti et al \cite{CR22}.


\subsubsection{Movement rules}

While it may seem complicated at first sight, the rules are pretty
simple and easy to understand.

\begin{itemize}
\item  
The players move alternately, each move must land on an empty space
on the game board.

Other versions of rithmomachia list this rule: ``No move
  can return the game to a previous state.'' This is called
  the {\it no repetition rule}.
  This one is not used by de Boissi\'ere, so it is ignored here.
\index{no repetition rule}

\item  
No piece is allowed to jump over another, a move
possible in some other versions of rithmomachia. Howeever, de Boissi\`ere
did not approve, calling the chess knight which can jump over other pieces,
``the mad warrior'' (see page 199 in \cite{Ri46}).

\item  
Each shape has a different movement (although, the pyramid piece
is skipped over for the time being).
\item
The movement of a piece doesn't depend on its value.
\item  
A circle piece moves exactly one space vertically
or horizontally (but not diagonally) in any direction.
\index{piece!circle}

{\footnotesize{
\begin{center}
\begin{tabular}{c|c|c|c|c}
 &            &             &          & \\ \hline
 &             & \circlet  &         & \\ \hline
 & \circlet    &    *     & \circlet & \\ \hline
 &             & \circlet  &         & \\ \hline
 &            &             &          & \\ \hline
\end{tabular}
\vskip .1in
{\bf Table: Possible \circlet, \circletfill\, moves from *}
\end{center}
}}
%provided, in doing so,
%\begin{itemize}
%\item
%  there is no enemy piece in the intended position, and
%\item
%  the piece doesn't step off the edge of the board. 
%\end{itemize}
\item
A triangular piece moves two spaces in any horizontal or vertical
direction.
\index{piece!triangle}

{\footnotesize{
\begin{center}
\begin{tabular}{c|c|c|c|c}
  \quad \    &     \quad \         &  \trianglepa           &   \quad  \      &   \quad \   \\ \hline
 &             &   &         & \\ \hline
 \trianglepa  &    &    *     &  &  \trianglepa  \\ \hline
 &             &   &         & \\ \hline
 &        \quad         &    \trianglepa           &          & \\ \hline
\end{tabular}
\vskip .1in
{\bf Table: Possible \trianglepa, \trianglepafill\, moves from *}
\end{center}
}}
  
\item  
A square piece moves exactly three squares but only vertically or horizontally.
\index{piece!square}

{\footnotesize{
\begin{center}
\begin{tabular}{c|c|c|c|c|c|c}
  \quad  \    &  \quad \   &      \quad   \       &  \squadfill           &   \quad     \      & \quad  \   &   \quad   \   \\ \hline
 &  &             &                 &          &   &      \\ \hline
 &  &             &                 &          &   &      \\ \hline
 \squadfill  &  &            &    *          &            &   & \squadfill \\ \hline
 &  &             &                 &          &   &      \\ \hline
 &  &             &                 &          &   &      \\ \hline
 &  &            &  \squadfill           &          &  & \\ \hline
\end{tabular}
\vskip .1in
{\bf Table: Possible \squad, \squadfill\, moves from *}
\end{center}
}}
  
\item
For clarity we repeat: in this version of Rithmomachia no piece can move (or capture) diagonally.
\end{itemize}


\begin{remark}
\label{rmk:parity1}
{\rm
A remark on a parity condition that these rules determine:

\begin{itemize}
\item  
  %1)
  As a consequence of these conditions, a triangle piece of either color
with initial position in column 1 can never move to an even column
(since it must move 2 spaces, hence this triangle can only move
to an odd column).

\item  
  %2)
  Likewise, if it's initially in column 2 it can never move to an odd
  column.
\end{itemize}
}
\end{remark}
\index{parity conditions}

\begin{question}
  \label{question:parity1}
A similar condition to that in Remark \ref{rmk:parity1}
holds for the squares. Can you find it?
\end{question}  


%\begin{figure}[h!] 
%\centering 
%\includegraphics[height=7.5cm,width=13.5cm]{rithmomachia-after-S81-60-30-with-pyramid2.jpg} 
%\caption{\textcolor{gray}{Rithmomachia position after the move $\square^{81}$a7a4.}}
%\label{fig:rithmomachia-after-S81-60-30-with-pyramid2} 
%\end{figure}
%%      sage: GSp = board_initial_matrix(pyramid_decomposition = True)
%%      sage: GS = move_piece(GSp, (6, 0), (3, 0))
%%      sage: board_plot2(GS, pyramid=True)

\begin{figure}[h!] 
\centering 
\begin{tabular}{|c|c|c|c|c|c|c|c||c} \hline
\squadfill{\,}49 & \squadfill{\,}121 & \quad  & \quad  & \quad  & \quad  & \squadfill{\,}225 & \squadfill{\,}361 & $p$ \\ \hline
\squadfill{\,}28 & \squadfill{\,}66 & \trianglepafill{\,}36 & \trianglepafill{\,}30 & \trianglepafill{\,}56 & \trianglepafill{\,}64 & \squadfill{\,}120 & \rhombusfill{\,}$p^{190}$ & $o$ \\ \hline
\trianglepafill{\,}16 & \trianglepafill{\,}12 & \circletfill{\,}9 & \circletfill{\,}25 & \circletfill{\,}49 & \circletfill{\,}81 & \trianglepafill{\,}90 & \trianglepafill{\,}100 & $n$ \\ \hline
\quad  & \quad  & \circletfill{\,}3 & \circletfill{\,}5 & \circletfill{\,}7 & \circletfill{\,}9 & \quad  & \quad  & $m$ \\ \hline
\quad  & \quad  & \quad  & \quad  & \quad  & \quad  & \quad  & \quad  & $l$ \\ \hline
\quad  & \quad  & \quad  & \quad  & \quad  & \quad  & \quad  & \quad  & $k$ \\ \hline
\quad  & \quad  & \quad  & \quad  & \quad  & \quad  & \quad  & \quad  & $j$ \\ \hline
\quad  & \quad  & \quad  & \quad  & \quad  & \quad  & \quad  & \quad  & $i$ \\ \hline
\quad  & \quad  & \quad  & \quad  & \quad  & \quad  & \quad  & \quad  & $h$ \\ \hline
\quad  & \quad  & \quad  & \quad  & \quad  & \quad  & \quad  & \quad  & $g$ \\ \hline
\quad  & \quad  & \quad  & \quad  & \quad  & \quad  & \quad  & \quad  & $f$ \\ \hline
\quad  & \quad  & \quad  & \quad  & \quad  & \quad  & \quad  & \quad  & $e$ \\ \hline
\quad  & \quad  & \circlet{\,}8 & \circlet{\,}6 & \circlet{\,}4 & \circlet{\,}2 & \quad  & \quad  & $d$ \\ \hline
\trianglepa{\,}81 & \trianglepa{\,}72 & \circlet{\,}64 & \circlet{\,}36 & \circlet{\,}16 & \circlet{\,}4 & \trianglepa{\,}6 & \trianglepa{\,}9 & $c$ \\ \hline
\squad{\,}153 & \rhombus{\,}$P^{91}$ & \trianglepa{\,}49 & \trianglepa{\,}42 & \trianglepa{\,}20 & \trianglepa{\,}25 & \squad{\,}45 & \squad{\,}15 & $b$ \\ \hline
\squad{\,}289 & \squad{\,}169 & \quad  & \squad{\,}81 & \quad  & \quad  & \quad  & \squad{\,}25 & $a$ \\ \hline
1 &  2 &  3 &  4 &  5 &  6 &  7 &  8  &  \\
\end{tabular} 
\caption{\textcolor{gray}{Rithmomachia position after the move $\square^{81}$a7a4.}}
\label{fig:rithmomachia-after-S81-60-30-with-pyramid2} 
\end{figure}
%%   sage: pclst = game_state_to_piece_list(GS)
%%   sage: piece_list_to_latex(pclst, verbose = False)

When a piece cannot move (due to the positions of other pieces)
then it is {\it blocked}.
\index{piece!blocked}
        
\subsubsection{The pyramid}
\label{sec:pyramid1}

In the de Boissi\`ere's set up, the pyramid piece moves just like the square piece.
In all versions of rithmomachia,
the pyramid piece can move to an unoccupied coordinate/position
if and only if one of its component sub-pieces can legally move
to that position.
Since some versions allow circle and triangular
subpieces of the pyramid (e.g., \cite{Ne98} and \cite{DG21}),
such versions endow more flexibility
of movement on the pyramid than we do here.

The {\it white pyramid}, is initially set to be $S^{91}$,
the white square with (total) value $91$. 
Create the following six white {\it sub-pieces}, depicted in
Figure \ref{fig:chatGPT_rithmomachia-logo12b}.

\begin{center}
{\squad}36, {\squad}25, {\squad}16,
{\squad}9, {\squad\,}4, {\squad\,}1,
\end{center}
which are called the {\it sub-piece components} of the pyramid.
These are not separate pieces, but can be involved with
captures, as we will explain in more detail in \S \ref{sec:pyramid2}.
\index{piece!pyramid}
\index{pyramid!component sub-pieces}
\index{piece!value(s)}
%The de Boissi\'ere decomposition of the pyramid piece
%displayed here differs from that in other versions of rithmomachia.


\begin{figure}[h!] 
\centering 
\includegraphics[height=10.5cm,width=10.0cm]{Rithmomachia-game-pieces/chatGPT_rithmomachia-logo12b.png} 
\caption{\rithmomachia pyramids are multi-valued because they have
  subpieces. Thanks chatGPT for the visual humor!}
\label{fig:chatGPT_rithmomachia-logo12b}
\end{figure}


Imagine a layer of $36$ equally spaced dots arranged in a $6\times 6$ grid.
Centered over that, imagine a layer of $25$ equally spaced dots arranged in a $5\times 5$ grid.
Continuing in this fashion, on to of that stack a $4\times 4$ layer,
a $3\times 3$ layer, a $2\times 2$ layer, and finally a dot on top.
This is one way to imagine the White pyramid.

%\begin{remark}
%
%\end{remark}

The {\it black pyramid}, is initially set to be $s^{190}$, the
black square $190$.
Create the following five black {\it sub-pieces}:
%\footnote{Again,
%  these are not separate pieces, but only can be used (and potentially
%  captured) as part of the pyramid.}:
\begin{center}
{\squadfill}64, {\squadfill}49, {\squadfill}36, {\squadfill}25,
{\squadfill}16,
\end{center}
which are called the {\it components} of the pyramid.
\index{components}
In the polynomial ring with variables $c$, $C$, $t$, $T$, $s$, $S$,

\[
R = \zzz[c,C,t,T,s,S],
\]
A sub-piece component only can be used (and potentially captured by White)
as part of the pyramid. The list of values of $p$ is, by design,
the sequence of five consecutiive squares whose total is $190$
(namely, $16, 25, 36, 49, 64$).
Imagine a layer of $64$ equally spaced dots arranged in an $8\times 8$ grid.
Centered over that, imagine a layer of $49$ equally spaced dots arranged in a $7\times 7$ grid.
Continuing in this fashion, on to of that stack a $6\times 6$ layer,
a $5\times 5$ layer, and finally a $4\times 4$ layer.
This is one way to imagine the Black pyramid.

\vskip .1in
\postersection{Pyramid decompositions}
Algebraically, White's pyramid can be initially represented
as

\[
P = S^{36}+S^{25}+S^{16}+S^9+S^4+S^1,
\]
or
\[
\rhombus = \square^{36} + \square^{25} +\square^{16} +\square^{9}
+\square^{4} +\square^{1},
\]
and its (list of) values as $[1, 4, 9, 16, 25, 36]$ and its total $91$.

Black's pyramid can be initially represented as
a polynomial belonging to $R$,
\[
p = s^{64}+s^{49}+s^{36}+s^{25}+s^{16},
\]
or
\begin{center}
\rhombusfill = {\squadfill}$^{{\,}64}$ + {\squadfill}$^{{\,}49}$ +
{\squadfill}$^{{\,}36}$ + {\squadfill}$^{{\,}25}$ + {\squadfill}$^{{\,}16}$,
\end{center}
and its (list of values) as $[16, 25, 36, 49, 64]$ and its total $190$.
\postersectionend
\vskip .1in

The Pythagoreans believed that elucidating such harmonious arithmetical relationships
is spiritually significant.

\index{pyramid!moves}
The rules on capturing a pyramid, or one of its component
sub-pieces, will be explained in \S \ref{sec:pyramid2}.



\begin{figure}[h!] 
\centering 
\includegraphics[height=2.5cm,width=14.0cm]{Rithmomachia-game-pieces/white-circles-blue-border-in-line} 
\caption{Rithmomachia White circles pieces in a line}
\label{fig:white-circles-blue-border-in-line} 
\end{figure}



  
\subsection{How the pieces capture}
\label{sec:capture-pieces}

As mention, a player's turn consists of captures and
a piece move. How do captures work?

We follow de Boissi\'ere’s set-up, as discussed in
\cite{Ri46}.


\vskip.1in
\postersection{How Pieces Capture}
%{\small{
Capture of an enemy piece does not mandate a friendly piece move but
requires permanent removal of the captured piece or subpiece.
In particular, a captured piece can never be placed back on the board,
even as the opposite color.
If a pyramid sub-piece is captured, then that subpiece is removed
and the value of the pyramid is deducted accordingly.
%The capturing player need not move their capturing piece to replace the captured piece.
%Chain captures are allowed.
Captures are not mandatory.
%}}
\postersectionend


A natural question arises: What about repeated captures?
They are allowed in checkers but not in chess.
What about for rithmomachia?
\index{checkers}

\begin{question}
%To the best of my knowledge, no mention has been found of
%repeated captures in the medieval rithmomachia game manuals.
%(Not unreasonable, as repeated captures in an actual game
%is probably a rare occurrence.)
%None-the-less, there is no rule that blocks such captures.
Assume you are playing a game and the following scenerio (or game state to use the
technical term) arises:

\begin{itemize}
\item[(a)]
you capture an opponent piece, and

\item[(b)]
its removal from the board exposes a second opponent piece for capture
(by the same friendly piece or by some other).
\end{itemize}
Can you take it?
\end{question}

There is no rithmomachia rule that stops you from capturing that second opponent
piece as well. If you are lucky enough to land in such a game state, the captures
are allowed.
 \index{capture!chain}
 When at least two pieces are captured in this way,
it is called a {\it chain capture}.

Like chess, captures in rithmomachia are not mandatory. There are several types:
\begin{itemize}
\item
capture by {\it numbering} (or {\it encounter}),
\item
by {\it addition} (or {\it ambush}),
\item
{\it subtraction} (also called an {\it ambush}),
\item
{\it multiplication} (or {\it assault}),
\item
{\it division} (also called an {\it assault}),
\item
and {\it siege} (or {\it surrounding}).
\end{itemize}
We will go into more detail about each of these below.

If we ignore the pyramid pieces for the moment,
a captured piece is {\it removed} from the board (and cannot be
returned or replaced or swapped). To be perfectly clear:
the capturing player need not move
their capturing piece (e.g., to replace the captured piece,
as in chess).
Captures involving the pyramids are treated separately below,
but to give a hint: if a subpiece of a pyramid is captured
then that subpiece is removed but the pyramid is not removed.
\index{capture!pyramid}
\index{capture!numbering}
\index{capture!addition}
\index{capture!division}
\index{capture!subtraction}
\index{capture!multiplication}

The captures by numbering, by addition, and by subtraction all
have the special property called {\it local captures}.
That is, such a capture is performed ``up close'' in the sense
that the capturing piece(s) can, in one move, land on captured piece.
\index{capture!local}


\vskip.1in
\postersection{Main ``local'' capture types}
\begin{itemize}%[leftmargin=*, label=\textcolor{RithmoGold}{\faCrosshairs}]
    \item \textbf{Numbering:} A piece captures an enemy piece if it can land on the enemy's square and both pieces have the \textbf{same value}.
    \item \textbf{Addition/Subtraction:} Two friendly pieces capture an enemy piece if the \textbf{sum/difference} of their values equals the enemy's value, and each friendly piece can legally land on the enemy's square.
    %\item \textbf{Subtraction:} Two friendly pieces capture an enemy if their \textbf{difference} in values equals the enemy's value, AND each friendly piece can land on it.
\end{itemize}
\postersectionend

On the other hand, capturing by multiplication
or by division has a different property, called {\it non-local}:
such a capture can be performed at a distance not equal to the
move distance.
\index{capture!at distance}


\vskip.1in
\postersection{Main ``non-local'' capture types}
\begin{itemize}%[leftmargin=*, label=\textcolor{RithmoGold}{\faCrosshairs}]
\item
  \textbf{Multiplication:} A friendly piece (value $v$) captures an enemy piece
  (value $m \cdot v$) if there are exactly $m$ vacant squares in a straight orthogonal line between them ($m>0$).
\item
  \textbf{Division:} A friendly piece (value $d \cdot v$) captures an enemy piece
  (value $v$) if there are exactly $d$ vacant squares in a straight orthogonal line between them ($d>0$).
\item
  \textbf{Siege:} An enemy piece is captured if it's surrounded by opposing pieces
  (or board edges) such that it has no legal moves.
%    \item \textbf{Pyramid Captures:}
%    \begin{itemize}
%        \item The pyramid can be captured using its total value.
%        \item A component sub-piece can be captured; its value is deducted from the pyramid's total.
%        \item The pyramid can capture using its total value or using the value of one of its component sub-pieces.
%    \end{itemize}
\end{itemize}
\postersectionend

An example of what this means is this: Suppose, during White's turn,
one or more of Black's pieces is captured.
White removes the capture(s) and now {\it must} make a move
of some piece to an unoccupied board position
%possibly the position that 
%was previously occupied by a captured piece but can move elsewhere
%as well
(see \cite{Ri46}, pp 202-203).
This is in contrast with chess - in chess, a capturing piece
{\it must} replace the captured piece.
%\footnote{With perhaps a small exception in the case of
%the {\it en passant} capture of a chess pawn.},
%However, in the de Boissi\`ere version of rithmomachia,
%such replacement is not required 
\index{capture!chess-like}

%Notation for a capture could be something like this:
%
%{\small{
%\begin{center}
%{\tt  capture = (piece(s)+value(s):coord(s))$\times$(capture piece+value:coord.)}.
%\end{center}
%}}

{\footnotesize{
\begin{center}
\begin{tabular}{|c|c|} \hline  
capture type   & description \\ \hline    
numbering  & 1 friendly takes 1 enemy, local \\
addition   & 2 friendlies take 1 enemy, local \\
subtraction & 2 friendlies take 1 enemy, local \\
multiplication & 1 friendly takes 1 enemy, non-local \\
division      &1 friendly takes 1 enemy, non-local \\
siege &  many friendlies takes 1 enemy, local \\  \hline
\end{tabular}
\vskip .1in
{\bf Table: piece capture types}
\end{center}
}}

\subsubsection{Numbering}

If a piece can move to a position where it lands on an enemy piece
with the same number, then the enemy piece may be captured.
%To emphasize what was already stated above:
%{\it the capturing piece does not move.}

For a numbering capture to occur, the values of two opposing pieces must
be equal. As a matter of fact, here are the only
values for which there is a piece of each color with that value:
$9$, $16$, $25$, $36$, $49$, $64$, $81$.

\vskip.1in
\postersection{Capture by numbering notation}
A capture by numbering could be written something like this:
%
{\small{
\begin{center}
{\tt  (attacking piece+value:coord.)$=$(captured piece+value:coord.)}.
\end{center}
}}
However, often the context is clear and we simply write
{\small{
\begin{center}
{\tt  (attacking piece+value)x(captured piece+value)}.
\end{center}
}}
\postersectionend

\begin{lemma}
\label{lemma:numbering-count}  
Given all pieces are on the board, there are $21$ possible
White captures Black by numbering. 
\end{lemma}  

\begin{proof}
The simplest verification is to simply list all such captures.

First we list White piece/Black piece pairs that satisfy in theory
capturing by numbering. Then we take parity conditions into account.
The initial coordinates of the pieces are included below to avoid
any ambiguities.

{\footnotesize{
\[
S^{81} (6, 0)  =  c^{81} (5, 13),\quad
T^{81} (0, 2)  =  c^{81} (5, 13),\quad
C^{64} (2, 2)  =  s^{64} (7, 14),
\]
\[
C^{64} (2, 2)  =  t^{64} (5, 14),\quad
T^{49} (2, 1)  =  s^{49} (0, 15),\quad
T^{49} (2, 1)  =  c^{49} (4, 13),
\]
\[
T^{49} (2, 1)  =  s^{49} (7, 14),\quad
C^{36} (3, 2)  =  s^{36} (7, 14),\quad
C^{36} (3, 2)  =  t^{36} (2, 14),
\]
\[
S^{36} (1, 1)  =  s^{36} (7, 14),\quad
S^{36} (1, 1)  =  t^{36} (2, 14),\quad
S^{25} (1, 1)  =  s^{25} (7, 14),
\]
\[
S^{25} (1, 1)  =  c^{25} (3, 13),\quad
S^{25} (7, 0)  =  s^{25} (7, 14),\quad
S^{25} (7, 0)  =  c^{25} (3, 13),
\]
\[
T^{25} (5, 1)  =  s^{25} (7, 14),\quad
T^{25} (5, 1)  =  c^{25} (3, 13),\quad
C^{16} (4, 2)  =  s^{16} (7, 14),
\]
\[
C^{16} (4, 2)  =  t^{16} (0, 13),\quad
S^{16} (1, 1)  =  s^{16} (7, 14),\quad
S^{16} (1, 1)  =  t^{16} (0, 13),
\]
\[
T^{9} (7, 2)  =  c^{9} (2, 13),\quad
T^{9} (7, 2)  =  c^{9} (5, 12),\quad
\]
\[
S^{9} (1, 1)  =  c^{9} (2, 13),\quad
S^{9} (1, 1)  =  c^{9} (5, 12).
\]
}}
%sage: wtb = white_takes_black_by_numbering()
%sage: for x in wtb:
%....:     print(latex(x[0][0]), x[0][2]," = ", latex(x[1][0]), x[1][2])

The parity conditions mentioned in Remark \ref{rmk:parity1} and
Question \ref{question:parity1} tell us that the following are impossible,

%{\small{
\[
S^{25} (7, 0)  =  s^{25} (7, 14),\quad
S^{16} (1, 1)  =  s^{16} (7, 14),\quad
\]
\[
S^{36} (1, 1)  =  s^{36} (7, 14),\quad
S^{25} (1, 1)  =  s^{25} (7, 14).
\]
%}}
but don't rule out the remaining ones.
This leaves 21.
\end{proof}

\begin{question}
%{\rm
%Ponder this:
%
%\noindent
%Problem:
Suppose it's White turn and a White triangle is able to
capture an Black square by numbering. By the movement rules, this implies
there is exactly {\it one} empty coordinate position between these
two opposing pieces.

Does this also means that, if it were instead Black's turn to move,
the enemy square can capture the triangle by numbering?
(Hint, hint\footnote{Indeed, the movement rules require {\it two} empty coordinate positions between them.}:-)
%
%}
\end{question}

None-the-less,
it turns out there are also $21$ Black captures White by numbering.

Before continuing, we need more notation.

\vskip.1in
\postersection{Modulus notation $\equiv$}
The modulus notation, introduced by
Carl Frederick Gauss in the 1800s, looks a little strange
\index{Gauss, Carl Frederick}
at first sight but it's easy to work
with. For integers $a,b,d>0$, the expression

\[
a \equiv b\pmod d
\]
is pronounced ``$a$ is congruent to $b$ modulo $d$''
and simply means that $d$ divides $b-a$. It's
treated in practice more as an equality rather than a
divisibility condition. For example, $13\equiv 1 \pmod{12}$
simply says that ``$13$ is congruent to $1$ modulo $12$''.
In non-technical language, that's the same as saying
$13$ o'clock is the same as $1$ o'clock. Think of
$a \equiv b\pmod d$ as ``$=$'' on the (analog) clock
face with $d$ numbers spaced equally clockwise around
the face.
\postersectionend

\begin{remark}
\label{rmk:parity2}
{\rm
A remark on symmetry and parity considerations expanding on
Remark \ref{rmk:parity1}.

Since a triangle moves by 2, if one starts at a coordinate
position $(i,j)$ then the only other positions it can possibly
land on during the game are $(i+2a, j+2b)$, for some integers $a$ and $b$.
To get a visual of this, imagine the Rithmomachia board is checkered,
like a chess board. The above identity implies that if a
triangle started on a black square, it can never land on a white square.
}
\end{remark}
\index{parity conditions}
\index{symmetry principle}

Therefore, we have the following {\it Triangle Parity Lemma}.
\index{Triangle Parity Lemma}
  
\begin{lemma}
\label{lemma:TPL}  
If we have one triangle at $(i_1,j_1)$
and another triangle at $(i_2,j_2)$ then one can land on
the other if and only if

\[
i_1\equiv i_2\pmod{2}\ \ \ \ \ {\rm and}\ \ \ \ j_1\equiv j_2\pmod{2}.
\]
\end{lemma}

\begin{remark}
\label{rmk:parity3}
{\rm
Similarly, if we have one square at $(i_1,j_1)$
and another square at $(i_2,j_2)$ then one can land on
the other if and only if

\[
i_1\equiv i_2\pmod{3}\ \ \ \ \ {\rm and}\ \ \ \ j_1\equiv j_2\pmod{3}.
\]
}
\end{remark}

Therefore, we have the following {\it Square Parity Lemma}.
\index{Square Parity Lemma}

\begin{lemma}
\label{lemma:SPL}  
If we have one square at $(i_1,j_1)$
and another square at $(i_2,j_2)$ then one can land on
the other if and only if

\[
i_1\equiv i_2\pmod{3}\ \ \ \ \ {\rm and}\ \ \ \ j_1\equiv j_2\pmod{3}.
\]
\end{lemma}
In particular, the sub-piece $S^{36}$ of the white pyramid
can never capture the sub-piece $s^{36}$ of the black pyramid
since they have different coordinate parities.
\index{capture!symmetry}


\subsubsection{Addition}

In this capture, two pieces on the same side must satisfy
these conditions:
\begin{itemize}
\item
there is an enemy piece that either capturing piece can legally land on,
\item
the value of the enemy piece is the sum of the values of the two
capturing pieces.
\end{itemize}
Then the enemy piece can be captured.

In particular, if White moves any piece of value $x$
to a coordinate that is 
\begin{itemize}
\item[(a)] orthogonally next to a black circle piece with value $y$,
\item[(b)] orthogonally next to another black circle piece with value $z$,
\item[(c)] $x=y+z$, 
\end{itemize}
then Black can capture this White piece on the next move.



{\small{
\begin{center}
\begin{tabular}{c|c|c|c|}
  \quad \    &  \circletfill$^{\, x}$         &    \quad \   &                        \\ \hline
              &  \squad$^{x+y}$             &           &  \trianglepafill$^y$\\ \hline
             &                           &            &                         \\ %\hline
\end{tabular}
\vskip .1in
{\bf Table: Addition: Black/Odd pieces of value $x, y$ capturing a White piece of value $x+y$}
\end{center}
}}


\vskip.1in
\postersection{Capturing by addition notation}
A capture by addition could be written something like this:
%
{\small{
\begin{center}
{\tt  [(attacking piece1+value1:coord.)+(attacking piece2+value2:coord.))$=$(captured piece3+value3:coord.)}.
\end{center}
}}
However, often the context is clear and we simply write
{\small{
\begin{center}
{\tt  [(attacking piece1+value1)+(attacking piece2+value2)]x(captured piece3+value3)}.
\end{center}
Example: In the table above, we have $[c^x(i-1, j)+t^y(i+2,j)]\times S^{x+y}(i,j)$, where $i$ is the rank
(row number) of the captured piece and $j$ is the file (column number).
}}
\postersectionend

Another example of Black capturing by addition is depicted in
Figure \ref{fig:rithmomachia-capture-by-addition-subtraction-with-pyramid3}.
It is written in coordinate notation as,

{\small{
\begin{center}
{\tt [$c^3$(1,3) + $c^5$(2,4)]$\times C^8$(2,3)} (addition)
\end{center}
}}

\noindent
and in algebraic notation, as

{\small{
\begin{center}
{\tt [$c^3$d7 + $c^5$e6]$\times C^8$d6} (addition)\ .
\end{center}
}}

There are at most $118$ ways for White to capture Black by
addition. To be perfectly clear, the list of the value
combinations that allow for a capture by addition has length
$118$. As in the proof of Lemma \ref{lemma:numbering-count},
one must check whether parity conditions rule any of these captures out.

Some examples are listed below (using subscripts when there is
ambiguity):

\[
[S^{289}   +  T^{72}]  \times  s^{361},\quad
[S^{153}   +  T^{72}]  \times  s^{225},\quad
[P^{91}   +  T^{9}] \times  t^{100},\quad
\]
\[
[P^{91}   +  S^{9}] \times  t^{100},\quad
[S^{81}  +  T^{9}] \times  t^{90},\quad
[S^{81}  +  S^{9}] \times  t^{90},
\]
\[
[T^{81}  +  T^{9}] \times  t^{90},\quad
[T^{81}  +  S^{9}] \times  t^{90}, \quad\dots \ .
\]

There are at most $82$ ways for Black to capture White by addition.
Some examples are listed below:

\[
[s^{225}  +  s^{64}] \times  S^{289},\quad
[s^{225}  +  t^{64}] \times  S^{289},\quad
[s^{120}  +  c^{49}] \times  S^{169},
\]
\[
[s^{120}  +  s_1^{49}] \times  S^{169},\quad [s^{120}  +  s_2^{49}] \times  S^{169},\quad
[s^{66}  +  s^{25}] \times  P^{91}, \quad\dots \ .
\]

\begin{question}
What is the exact number of ways for Black to capture White by addition?
Same question for White to capture Black by addition.
\end{question}  

\subsubsection{Subtraction}

Very similar to addition: two pieces on the same side must satisfy
these conditions:
\begin{itemize}
\item
there is an enemy piece that either capturing piece can legally land on,
\item
the value of the enemy piece is the difference of the values of the two
capturing pieces.
\end{itemize}
Then the enemy piece can be captured.

\vskip.1in
\postersection{Capturing by subtraction notation}
A capture by subtraction could be written something like this:
%
{\small{
\begin{center}
{\tt  [(attacking piece1+value1:coord.)-(attacking piece2+value2:coord.))$=$(captured piece3+value3:coord.)}.
\end{center}
}}
However, often the context is clear and we simply write
{\small{
\begin{center}
{\tt  [(attacking piece1+value1)-(attacking piece2+value2)]x(captured piece3+value3)}.
\end{center}
}}
\postersectionend

\vskip .2in

\noindent
Example: In Figure \ref{fig:rithmomachia-capture-by-addition-subtraction-with-pyramid3},
Black captures the piece $\bigcirc^2$ by subtraction,
{\small{
\begin{center}
{\tt [$c^9$e3 - $c^7$e5]$\times C^2$d4} (subtraction).
\end{center}
}}
or $(c^9-c^7)\times C^2$.

\begin{figure}[h!] 
\centering 
%\includegraphics[height=7.0cm,width=13.5cm]{rithmomachia-capture-by-addition-subtraction-with-pyramid3.jpg} 
\begin{tabular}{|c|c|c|c|c|c|c|c||c} \hline
\squadfill{\,}49 & \squadfill{\,}121 & \quad  & \quad  & \quad  & \quad  & \squadfill{\,}225 & \squadfill{\,}361 & $p$ \\ \hline
\squadfill{\,}28 & \squadfill{\,}66 & \trianglepafill{\,}36 & \trianglepafill{\,}30 & \trianglepafill{\,}56 & \trianglepafill{\,}64 & \squadfill{\,}120 & \rhombusfill{\,}$p^{190}$ & $o$ \\ \hline
\trianglepafill{\,}16 & \trianglepafill{\,}12 & \circletfill{\,}9 & \circletfill{\,}25 & \circletfill{\,}49 & \circletfill{\,}81 & \trianglepafill{\,}90 & \trianglepafill{\,}100 & $n$ \\ \hline
\quad  & \quad  & \quad  & \quad  & \quad  & \quad  & \quad  & \quad  & $m$ \\ \hline
\quad  & \quad  & \quad  & \quad  & \quad  & \quad  & \quad  & \quad  & $l$ \\ \hline
\quad  & \quad  & \quad  & \quad  & \quad  & \quad  & \quad  & \quad  & $k$ \\ \hline
\quad  & \quad  & \quad  & \quad  & \quad  & \quad  & \quad  & \quad  & $j$ \\ \hline
\quad  & \quad  & \quad  & \quad  & \quad  & \quad  & \quad  & \quad  & $i$ \\ \hline
\quad  & \quad  & \quad  & \quad  & \quad  & \quad  & \quad  & \quad  & $h$ \\ \hline
\quad  & \quad  & \quad  & \quad  & \quad  & \quad  & \quad  & \quad  & $g$ \\ \hline
\quad  & \quad  & \quad  & \quad  & \quad  & \quad  & \quad  & \quad  & $f$ \\ \hline
\quad  & \quad  & \circletfill{\,}5 & \quad  & \quad  & \circletfill{\,}9 & \quad  & \quad  & $e$ \\ \hline
\quad  & \circletfill{\,}3 & \circlet{\,}8 & \circlet{\,}6 & \circlet{\,}4 & \circlet{\,}2 & \circletfill{\,}7 & \quad  & $d$ \\ \hline
\trianglepa{\,}81 & \trianglepa{\,}72 & \circlet{\,}64 & \circlet{\,}36 & \circlet{\,}16 & \circlet{\,}4 & \trianglepa{\,}6 & \trianglepa{\,}9 & $c$ \\ \hline
\squad{\,}153 & \rhombus{\,}$P^{91}$ & \trianglepa{\,}49 & \trianglepa{\,}42 & \trianglepa{\,}20 & \trianglepa{\,}25 & \squad{\,}45 & \squad{\,}15 & $b$ \\ \hline
\squad{\,}289 & \squad{\,}169 & \quad  & \quad  & \quad  & \quad  & \squad{\,}81 & \squad{\,}25 & $a$ \\ \hline
1 &  2 &  3 &  4 &  5 &  6 &  7 &  8  &  \\
\end{tabular} 
\caption{\textcolor{gray}{Rithmomachia position before a capture by addition and a capture by subtraction.}}
\label{fig:rithmomachia-capture-by-addition-subtraction-with-pyramid3} 
\end{figure}
% sage: GSp = board_initial_matrix(pyramid_decomposition = True)
% sage: GS = copy(GSp)
% sage: GS[5, 4] = c^9; GS[5, 12] = 0; GS[6, 3] = c^7; GS[4, 12] = 0; GS[1, 3] = c^3; GS[2, 12] = 0; GS[2, 4] = c^5; GS[3, 12] = 0
% sage: board_plot(GS)
%%sage: pclst = game_state_to_piece_list(GS)
%%sage: piece_list_to_latex(pclst, verbose = False)


There are at most $141$ ways for White to capture Black by subtraction.
Some examples are listed below:

\[
[S^{289}   -  S^{169}]  \times  s^{120},\quad
[S^{289}   -  C^{64}]  \times  s^{225},\quad
[S^{169}   -  S^{153}]  \times  s^{16} (7, 14),
\]
\[
[S^{169}   -  S^{153}]  \times  t^{16},\quad
[S^{169}   -  T^{49}]  \times  s^{120},\quad
[S^{153}   -  T^{72}]  \times  c^{81},
\]
\[
[P^{91}   -  T^{42}]  \times  s_1^{49},\quad
[P^{91}   -  T^{42}]  \times  s_2^{49},\quad
[P^{91}   -  T^{42}]  \times  c^{49}, \quad \dots 
\]

There are at most $119$ ways for Black to capture White by subtraction.
Some examples are listed below:

\[
[s^{225}   -  t^{56}]  \times  S^{169},\quad
[s^{121}   -  s^{120}]  \times  S^{1} ,\quad
[s^{121}   -  c^{49}] \times  T^{72},
\]
\[
[s^{121}   -  s_1^{49}] \times  T^{72},\quad
[s^{121}   -  s_2^{49}] \times  T^{72},\quad
[s^{121}   -  t^{30}]  \times  P^{91}, \quad \dots
\]


\begin{question}
What is the exact number of ways for Black to capture White by subtraction?
Same question for White to capture Black by subtraction.
\end{question}

\subsubsection{Multiplication}

A friendly piece of value $v$ can capture an enemy piece of value
$m {\cdot} v$ {\it if} the number of vacant coordinates
along an unoccupied straight line (vertical, horizontal)
drawn between these two pieces is exactly $m>0$.
In the special case when
$m=1$ then capturing by multiplication coincides with capture by
{\it numbering}, in which case
it {\it is} necessary for the friendly capturing piece to
be able to land on the enemy captured piece.
However, outside of this special case, it is not generally
necessary for the friendly capture-by-multiplication piece
to be able to land on the enemy captured piece on their next move.


{\footnotesize{
\begin{center}
\begin{tabular}{c|c|c|c|c}
 &            &             &          & \\ \hline
\trianglepafill\, $x$  &  $\leftarrow$ \dots &    $m$     & \dots $\rightarrow$ & \squad\, $m\cdot x$\\ \hline
 &            &             &          & \\ %\hline
\end{tabular}
\vskip .1in
{\bf Table: Multiplication: Black piece of value $x$ capturing from a distance a White piece of value $mx$}
\end{center}
}}


\vskip.1in
\postersection{Capturing by multiplication notation}
A capture by multiplication could be written something like this:
%
{\small{
\begin{center}
{\tt  [(attacking piece1+value1:coord.)*distance))$\times$(captured piece2+value2:coord.)}.
\end{center}
}}
However, often the context is clear and we simply write
{\small{
\begin{center}
{\tt  [(attacking piece1+value1)*distance]x(captured piece2+value2)}.
\end{center}
Example: Figure \ref{fig:rithmomachia-capture-by-multiplication-division-with-pyramid} depicts an
example of Black capturing the piece $S^{15}$ by multiplication, denoted

\begin{center}
{\tt [$c^5$m4 * 3]$\times S^{15}$i4 (multiplication)}
\end{center}
}}
\postersectionend
\index{capture!by multiplication}


\begin{figure}[h!] 
\centering 
%\includegraphics[height=7.0cm,width=13.5cm]{rithmomachia-capture-by-multiplication-division-with-pyramid.jpg} 
\begin{tabular}{|c|c|c|c|c|c|c|c||c} \hline
\squadfill{\,}49 & \squadfill{\,}121 & \quad  & \quad  & \quad  & \quad  & \squadfill{\,}225 & \squadfill{\,}361 & $p$ \\ \hline
\squadfill{\,}28 & \squadfill{\,}66 & \trianglepafill{\,}36 & \trianglepafill{\,}30 & \trianglepafill{\,}56 & \trianglepafill{\,}64 & \squadfill{\,}120 & \rhombusfill{\,}$p^{190}$ & $o$ \\ \hline
\trianglepafill{\,}16 & \trianglepafill{\,}12 & \circletfill{\,}9 & \circletfill{\,}25 & \circletfill{\,}49 & \circletfill{\,}81 & \trianglepafill{\,}90 & \trianglepafill{\,}100 & $n$ \\ \hline
\quad  & \quad  & \circletfill{\,}3 & \circletfill{\,}5 & \circletfill{\,}7 & \circletfill{\,}9 & \quad  & \quad  & $m$ \\ \hline
\quad  & \quad  & \quad  & \quad  & \quad  & \quad  & \quad  & \quad  & $l$ \\ \hline
\quad  & \quad  & \quad  & \quad  & \quad  & \quad  & \quad  & \quad  & $k$ \\ \hline
\quad  & \quad  & \quad  & \quad  & \quad  & \quad  & \quad  & \quad  & $j$ \\ \hline
\quad  & \quad  & \quad  & \squad{\,}15 & \quad  & \quad  & \quad  & \quad  & $i$ \\ \hline
\quad  & \quad  & \quad  & \quad  & \quad  & \quad  & \quad  & \quad  & $h$ \\ \hline
\quad  & \quad  & \quad  & \quad  & \quad  & \quad  & \quad  & \quad  & $g$ \\ \hline
\quad  & \quad  & \quad  & \quad  & \quad  & \quad  & \quad  & \quad  & $f$ \\ \hline
\quad  & \quad  & \quad  & \quad  & \quad  & \quad  & \quad  & \quad  & $e$ \\ \hline
\quad  & \quad  & \circlet{\,}8 & \circlet{\,}6 & \circlet{\,}4 & \circlet{\,}2 & \quad  & \quad  & $d$ \\ \hline
\trianglepa{\,}81 & \trianglepa{\,}72 & \circlet{\,}64 & \circlet{\,}36 & \circlet{\,}16 & \circlet{\,}4 & \trianglepa{\,}6 & \trianglepa{\,}9 & $c$ \\ \hline
\squad{\,}153 & \rhombus{\,}$P^{91}$ & \trianglepa{\,}49 & \trianglepa{\,}42 & \trianglepa{\,}20 & \trianglepa{\,}25 & \squad{\,}45 & \squad{\,}15 & $b$ \\ \hline
\squad{\,}289 & \squad{\,}169 & \quad  & \quad  & \quad  & \quad  & \quad  & \squad{\,}25 & $a$ \\ \hline
1 &  2 &  3 &  4 &  5 &  6 &  7 &  8  &  \\
\end{tabular} 
\caption{\textcolor{gray}{Rithmomachia position before a Black capture by multiplication: $c^5$ takes $S^{15}$, or $c^5*3\times S^{15}$}}
\label{fig:rithmomachia-capture-by-multiplication-division-with-pyramid} 
\end{figure}
%sage: GSp = board_initial_matrix(pyramid_decomposition = True)
%sage: GS = copy(GSp)
%sage: GS[3, 8] = S^(15)
%sage: GS[6, 0] = 0
%sage: board_plot2(GS, pyramid=True)
%%sage: pclst = game_state_to_piece_list(GS)
%%sage: piece_list_to_latex(pclst, verbose = False)

If you think about it, there can't be a ``huge'' number of possible captures by
multiplication. The reason why is the {\it multiplication capture value
constraint},

\begin{equation}
\label{eqn:mult-capture-value-constraint}
{\tt  value}_1*{\tt distance} = {\tt value}_2,
\end{equation}
\index{multiplication capture value constraint}
along with the the fact that both sides have only about 20 or so values
to begin with means we can list (with the aid of a computer to keep track
of the terms) all the possible cases in which
equation (\ref{eqn:mult-capture-value-constraint}) holds.
Some captures might be ruled out due to the parity
conditions in Lemmas \ref{lemma:TPL} and \ref{lemma:SPL},
but at least we will have an upper bound which is hopefully
stll pretty close.
These upper bounds are described below.

There are at most $96$ ways for White to capture Black by multiplication.
Some examples that are not capture by numbering are listed below:

\[
[S^{45}  *  5]  \times  s^{225}, \quad
[S^{45}  *  2]  \times  t^{90}, \quad
[S_1^{25}   *  9]  \times  s^{225}, \quad
\]
\[
[S_2^{25}   *  9]  \times  s^{225}, \quad
[S_1^{25}  *  4]  \times  t^{100}, \quad
[S_2^{25}  *  4]  \times  t^{100}, \dots
\]

There are at most $61$ ways for Black to capture White by multiplication.
Some examples that are not capture by numbering are listed below:

\[
[s^{16}  *  4]  \times  C^{64}, \quad
[t^{12}  *  6]  \times  T^{72},\quad
[c_1^{9}   *  9 ] \times  T^{81}, \quad
\]
\[
[c_2^{9}   *  9 ] \times  T^{81}, \quad
[c_1^{9}   *  9 ] \times  S^{81},\quad
[c_2^{9}   *  9 ] \times  S^{81},
\]
\[
[c^{7}  *  13]  \times  P^{91}, \quad
[c^{7}  *  7 ] \times  T^{49},\quad
[c^{3}  *  14]  \times  T^{42}, \quad
[c^{3}  *  12]  \times  C^{36}, \dots
\]

\begin{question}
Suppose on the very first move, White/Even plays
{\tt $T^{72}$c2e2}. Can you find a move by Black/Odd that will
enable the capture of that White \trianglepa?
\end{question}  
\index{opening trap}


\begin{figure}[h!] 
\centering 
\includegraphics[height=9.5cm,width=6.5cm]{Rithmomachia-game-pieces/Gemini_rithmomachia-battle5.png} 
\caption{\rithmomachia pieces can get intense. {\small{(Thanks gemini:-)}}}
\label{fig:Gemini_rithmomachia-battle5} 
\end{figure}


\begin{question}
What is the exact number of ways for Black to capture White by multiplication?
Same question for White to capture Black by multiplication.
\end{question}


\subsubsection{Division}

Similar to multiplication: A piece of value $d {\cdot} v$ can
capture an enemy piece of value $v$ {\it if} the number
of vacant coordinates
along an unoccupied straight line (vertical, horizontal)
drawn between these two pieces is exactly $d>0$.
Like capture by multiplication, this capture by division boils down to
capture by numbering if $d=1$.

\vskip.1in
\postersection{Capturing by division notation}
A capture by division could be written something like this:
%
{\small{
\begin{center}
{\tt  [(attacking piece1+value:coord.)$\div$ distance))$\times$(captured piece2+value2:coord.)}.
\end{center}
}}
However, often the context is clear and we simply write
{\small{
\begin{center}
{\tt  [(attacking piece1+value1)$\div$ distance]x(captured piece2+value2)}.
\end{center}
}}
Example: On White's turn, the white square $S^{15}$
can capture $c^5$ by division ($\div$) in 
Figure \ref{fig:rithmomachia-capture-by-multiplication-division-with-pyramid}.
In coordinate  and algebraic notation:

{\small{
\begin{center}
{\tt [S15{:}(3,8)$\div$ 3] $\times$ c5{:}(3,12) (division)} or $[S^{15}i4\div 3]\times c^5m4$.
\end{center}
}}
\postersectionend
\index{capture!symmetry}
\index{capture!by division}

There are symmetry considerations for this capture as well.
  
\begin{lemma}
\label{lemma:M/DSP}
  (multiplication/division symmetry principle)
A White piece can capture a Black piece by division if and only if that
Black piece can capture that White piece by multiplication.
\end{lemma}
\index{symmetry principle for captures}

Looking at the counts for capturing by multiplication, the multiplication/division symmetry principle
(Lemma \ref{lemma:M/DSP} above) tells us:
\begin{itemize}
\item
there are at most $96$ ways for Black to capture White by division, and
\item
there are at most $61$ ways for White to capture Black by division.
\end{itemize}
\index{symmetry principle for captures}


\begin{question}
The game state displayed in the Rithmomachia diagram

{\small{
\renewcommand{\arraystretch}{1.5}
\[
\begin{array}{|c|c|c|c|c|c|c|c|} \hline
 s^{49} & s^{121} & \ & \ & \ & \ & s^{225} & s^{361} \\ \hline 
 s^{28} & s^{66} & t^{36} & t^{30} & t^{56} & t^{64} & s^{120} & p^{190} % + s^64 + s^49 + s^36 + s^25 + s^16
 \\ \hline 
 t^{16} & t^{12} & c^9 & c^{25} & c^{49} & c^{81} & t^{90} & t^{100} \\ \hline
 \ & \ & c^3 & c^5 & c^7 & c^9 & \ & \ \\ \hline 
 \ & \ & \ & \ & \ & \ & \ & \ \\ \hline
 \ & \ & \ & \ & \ & T^9 & \ & \ \\ \hline 
 \ & \ & \ & \ & \ & \ & \ & \ \\ \hline
 \ & \ & \ & \ & \ & \ & \ & \ \\ \hline 
 \ & \ & \ & \ & \ & \ & \ & \ \\ \hline
 \ & \ & \ & \ & \ & \ & \ & \ \\ \hline 
 \ & \ & \ & \ & \ & \ & \ & \ \\ \hline
 \ & \ & \ & \ & \ & \ & \ & \ \\ \hline 
 \ & \ & C^8 & C^6 & C^4 & C^2 & \ & \ \\ \hline
 T^{81} & T^{72} & C^{64} & C^{36} & C^{16} & C^4 & T^6 & \ \\ \hline 
 S^{153} & P^{91} %+ S^{36} + S^{25} + S^{16} + S^9 + S^4 + S
          & T^{49} & T^{42} & T^{20} & T^{25} & S^{45} & S^{15} \\ \hline 
 S^{289} & S^{169} & \ & \ & \ & \ & S^{81} & S^{25} \\ \hline
\end{array}
\]
}}

\vskip .1in
\noindent
has the curious property that a single piece can capture another by three
of the six different methods.
Can you name them?
\end{question}

{\footnotesize{
\begin{center}
\begin{tabular}{c|c|c|c|c}
 &            &             &          & \\ \hline
\circlet\, $m\cdot x$  &  $\leftarrow$ \dots &    $m$     & \dots $\rightarrow$ & \squadfill\, $x$\\ \hline
 &            &             &          & \\ %\hline
\end{tabular}
\vskip .1in
{\bf Table: Division: White piece of value $mx$ capturing from a distance a Black piece of value $x$}
\end{center}
}}

\begin{question}
Suppose on the very first move, White/Even played
{\tt $T^{72}$c2e2} and then Black/Odd played {\tt $c^9$m6l6},
but no captures. Can you find White's (second) move that will
enable the capture of a Black \trianglepafill?
\end{question}


\subsubsection{Siege}

We recall the definition from \cite{DG21}:
If an enemy piece is surrounded by opposing pieces in such a way that
the enemy piece has no legal moves then the ``besieged''
enemy piece may be captured. For example, this capture can
take place by an edge of the board, with the enemy piece partially surrounded by the enemy,
or away from the edge, with the piece entirely surrounded by enemy pieces.
\index{capture!siege}

Figure \ref{fig:rithmomachia-capture-by-siege-with-pyramid3} depicts an
example of White capturing the piece $c^{3}$ by siege.

\begin{figure}[h!] 
\centering
\begin{tabular}{|c|c|c|c|c|c|c|c||c} \hline
\squadfill{\,}49 & \squadfill{\,}121 & \quad  & \quad  & \quad  & \quad  & \squadfill{\,}225 & \squadfill{\,}361 & $p$ \\ \hline
\squadfill{\,}28 & \squadfill{\,}66 & \trianglepafill{\,}36 & \trianglepafill{\,}30 & \trianglepafill{\,}56 & \trianglepafill{\,}64 & \squadfill{\,}120 & \rhombusfill{\,}$p^{190}$ & $o$ \\ \hline
\trianglepafill{\,}16 & \trianglepafill{\,}12 & \circletfill{\,}9 & \circletfill{\,}25 & \circletfill{\,}49 & \circletfill{\,}81 & \trianglepafill{\,}90 & \trianglepafill{\,}100 & $n$ \\ \hline
\quad  & \quad  & \quad  & \circletfill{\,}5 & \circletfill{\,}7 & \circletfill{\,}9 & \quad  & \quad  & $m$ \\ \hline
\quad  & \quad  & \quad  & \quad  & \quad  & \quad  & \quad  & \quad  & $l$ \\ \hline
\quad  & \quad  & \quad  & \quad  & \quad  & \quad  & \quad  & \quad  & $k$ \\ \hline
\quad  & \quad  & \quad  & \quad  & \quad  & \quad  & \quad  & \quad  & $j$ \\ \hline
\quad  & \quad  & \quad  & \quad  & \quad  & \quad  & \quad  & \quad  & $i$ \\ \hline
\quad  & \quad  & \quad  & \quad  & \quad  & \quad  & \quad  & \quad  & $h$ \\ \hline
\quad  & \quad  & \quad  & \quad  & \quad  & \quad  & \quad  & \quad  & $g$ \\ \hline
\quad  & \quad  & \quad  & \quad  & \quad  & \quad  & \quad  & \quad  & $f$ \\ \hline
\quad  & \quad  & \quad  & \quad  & \quad  & \quad  & \quad  & \quad  & $e$ \\ \hline
\quad  & \quad  & \circlet{\,}8 & \circlet{\,}6 & \circlet{\,}4 & \circlet{\,}2 & \quad  & \quad  & $d$ \\ \hline
\trianglepa{\,}81 & \trianglepa{\,}72 & \circletfill{\,}3 & \circlet{\,}36 & \circlet{\,}16 & \circlet{\,}4 & \trianglepa{\,}6 & \trianglepa{\,}9 & $c$ \\ \hline
\squad{\,}153 & \rhombus{\,}$P^{91}$ & \trianglepa{\,}49 & \trianglepa{\,}42 & \trianglepa{\,}20 & \trianglepa{\,}25 & \squad{\,}45 & \squad{\,}15 & $b$ \\ \hline
\squad{\,}289 & \squad{\,}169 & \quad  & \quad  & \quad  & \quad  & \squad{\,}81 & \squad{\,}25 & $a$ \\ \hline
1 &  2 &  3 &  4 &  5 &  6 &  7 &  8  &  \\
\end{tabular}  
%\includegraphics[height=7.0cm,width=13.5cm]{rithmomachia-capture-by-siege-with-pyramid3.jpg} 
\caption{\textcolor{gray}{Rithmomachia position before a capture by siege, $c^3$ is surrounded by White and can't move.}}
\label{fig:rithmomachia-capture-by-siege-with-pyramid3} 
\end{figure}
%% sage: GSp = board_initial_matrix(pyramid_decomposition = True)
%% sage: GS = copy(GSp)
%% sage: GS[2, 2] = c^3
%% sage: GS[2, 12] = 0
%% sage: board_plot2(GS, pyramid=True)
%%sage: pclst = game_state_to_piece_list(GS)
%%sage: piece_list_to_latex(pclst, verbose = False)

%{\small{
%\begin{center}
%{\tt {White}\,$\times$\,c3:(2,0) (siege)}
%\end{center}
%}}

\vskip.1in
\postersection{Capturing by siege notation}
A capture by siege could be written something like this:
if {\tt attacking piece1} is the last piece to move to help surround and block
{\tt captured piece2} then
%
{\small{
\begin{center}
{\tt  (attacking piece1+value:coord.)$\cap$(captured piece2+value2:coord.)}.
\end{center}
}}
However, often the context is clear and we simply write
{\small{
\begin{center}
{\tt  (attacking piece1+value1)$\cap$ (captured piece2+value2)}.
\end{center}
}}
{\it Example}: If White just moved his circle $C^8$ to $d3$ then this notation becomes

{\small{
\begin{center}
{\tt $C^8d3$\,$\cap$\,$c^3c3$ (siege)}
\end{center}
}}
\postersectionend


\subsubsection{Captures of/with the pyramid}
\label{sec:pyramid2}

The pyramid piece can be captured using its total value. 
In addition, any one of its component sub-pieces can be
captured.%\footnote{There are variations on this rule but this version is the
%simplest.}, as in Figure \ref{fig:rithmomachia-captureS25-by-c25numbering-with-pyramid}.
\index{capture!pyramid}
\index{pyramid!capture}
In case only a sub-piece of the pyramid is captured,
only that subpiece is removed from the board, never to return,
and the value of that sub-piece is deducted from that of the pyramid piece.
In the case the pyramid is captured using its full value,
the entire piece (and all its component sub-pieces) must be
removed from the board.

\begin{example}
For the original pyramid piece $p^{190}$ (see \S \ref{sec:pyramid1} above),
Figure \ref{fig:rithmomachia-capture-C36-with-pyramid2} depicts
a game state in which the component sub-piece $s^{36}$, the black
square sub-piece of $p$ with value $36$, can capture the white
circle $C^{36}$. 
\end{example}

\begin{question}
What other captures (for both sides) can you find in the study
depicted in Figure \ref{fig:rithmomachia-capture-C36-with-pyramid2}?
\end{question}
  

The pyramid can capture an enemy piece, using its total value.
Alternatively, the pyramid can capture an enemy piece including a
component sub-piece of the enemy's pyramid, using the value of any
one of its component sub-pieces.

\begin{example}
If $P=S^{91}$ was against the edge of the board and surrounded on the three
other sides by $c^3$, $c^7$, and $c^{81}$,
then Black can capture $P$ (as a whole) by siege. However, it cannot capture $P$
by addition.
\end{example}

\begin{figure}[h!] 
\centering 
%\includegraphics[height=7.0cm,width=13.5cm]{rithmomachia-capture-C36-with-pyramid2.jpg}
\begin{tabular}{|c|c|c|c|c|c|c|c||c} \hline
\squadfill{\,}49 & \squadfill{\,}121 & \quad  & \quad  & \quad  & \quad  & \squadfill{\,}225 & \squadfill{\,}361 & $p$ \\ \hline
\squadfill{\,}28 & \squadfill{\,}66 & \trianglepafill{\,}36 & \trianglepafill{\,}30 & \trianglepafill{\,}56 & \trianglepafill{\,}64 & \squadfill{\,}120 & \quad  & $o$ \\ \hline
\trianglepafill{\,}16 & \trianglepafill{\,}12 & \quad  & \circletfill{\,}25 & \circletfill{\,}49 & \circletfill{\,}81 & \trianglepafill{\,}90 & \trianglepafill{\,}100 & $n$ \\ \hline
\quad  & \quad  & \circletfill{\,}3 & \circletfill{\,}5 & \circletfill{\,}7 & \circletfill{\,}9 & \quad  & \quad  & $m$ \\ \hline
\quad  & \quad  & \quad  & \circletfill{\,}9 & \quad  & \quad  & \quad  & \quad  & $l$ \\ \hline
\quad  & \quad  & \quad  & \quad  & \quad  & \quad  & \quad  & \quad  & $k$ \\ \hline
\quad  & \quad  & \quad  & \quad  & \quad  & \quad  & \quad  & \quad  & $j$ \\ \hline
\quad  & \quad  & \quad  & \rhombus{\,}$P^{91}$ & \quad  & \quad  & \quad  & \quad  & $i$ \\ \hline
\quad  & \quad  & \quad  & \quad  & \quad  & \quad  & \quad  & \quad  & $h$ \\ \hline
\quad  & \quad  & \quad  & \quad  & \quad  & \quad  & \quad  & \quad  & $g$ \\ \hline
\quad  & \quad  & \quad  & \rhombusfill{\,}$p^{190}$ & \quad  & \quad  & \quad  & \quad  & $f$ \\ \hline
\quad  & \quad  & \quad  & \quad  & \quad  & \quad  & \quad  & \quad  & $e$ \\ \hline
\quad  & \quad  & \circlet{\,}8 & \quad  & \circlet{\,}4 & \circlet{\,}2 & \quad  & \quad  & $d$ \\ \hline
\trianglepa{\,}81 & \trianglepa{\,}72 & \circlet{\,}64 & \circlet{\,}36 & \circlet{\,}16 & \circlet{\,}4 & \trianglepa{\,}6 & \trianglepa{\,}9 & $c$ \\ \hline
\squad{\,}153 & \quad  & \trianglepa{\,}49 & \trianglepa{\,}42 & \trianglepa{\,}20 & \trianglepa{\,}25 & \squad{\,}45 & \squad{\,}15 & $b$ \\ \hline
\squad{\,}289 & \squad{\,}169 & \quad  & \quad  & \quad  & \quad  & \squad{\,}81 & \squad{\,}25 & $a$ \\ \hline
1 &  2 &  3 &  4 &  5 &  6 &  7 &  8  &  \\
\end{tabular} '
\caption{\textcolor{gray}{A rithmomachia position before the triple capture using component
    subpieces of the pyramid: $p^{190}f4\times S^{16}i4$, $p^{190}f4\times S^{25}i4$, $p^{190}f4\times S^{36}i4$.
    This chain capure is an example of a series of capture by numberings.}}
\label{fig:rithmomachia-capture-C36-with-pyramid2} 
\end{figure}
\index{pyramid!capture}
%sage: GSp = board_initial_matrix(pyramid_decomposition = True)
%sage: GS = copy(GSp)
%sage: GS[7, 14] = 0; GS[3,5] = p^(190) + s^(64) + s^(49) + s^(36) + s^(25) + s^(16)
%sage: GS[1, 1] = 0; GS[3, 8] = P^(91) + S^(36) + S^(25) + S^(16) + S^9 + S^4 + S; GS[3,3] = 0
%sage: GS[2,13] = 0; GS[3, 11] = c^9
%%sage: pclst = game_state_to_piece_list(GS)
%%sage: piece_list_to_latex(pclst, verbose = False)

Here is a partial game involving a pyramid capture:

\[
\begin{array}{l|ll|ll|} \hline
  \quad & Even\ moves & captures & Odd\ moves & captures \\ \hline
  1 & T^9{\tt c8e8} &  & t^{12}{\tt n2l2} &   \\
\end{array}
\]
This last move by Black makes a threat that White ignores.

\[
\begin{array}{l|ll|ll|} \hline
  \quad & Even\ moves & captures & Odd\ moves & captures \\ \hline  
  2 & T^6{\tt c7e7} &  & t^{12}{\tt l2j2} & t^{12}\times T^{72}  \\
  3 & T^9{\tt e8g8} &  & t^{12}{\tt j2h2} &   \\
\end{array}
\]
This last move by Black makes a threat that White also ignores.

\[
\begin{array}{l|ll|ll|} \hline
  \quad & Even\ moves & captures & Odd\ moves & captures \\ \hline  
  4 & S^{15}{\tt b8e8} &  & t^{12}{\tt h2f2} & t^{12}\times S^{36}  \\
\end{array}
\]
This last partured piece $S^{36}$ was a subpiece of the 
pyramid $P^{91}$. That pyramid has had it's value deducted, making it now

\[
P^{55} = S^{25}+S^{16}+S^9+S^4+S^1.
\]


\begin{example}
A slaughterfest of captures awaits the curious player
in the example study depicted in
Figure \ref{fig:rithmomachia-numerous-captures-with-pyramid}.

\begin{question}
Can you find all captures in
Figure \ref{fig:rithmomachia-numerous-captures-with-pyramid}, if
\begin{itemize}
\item[(a)]
  it is White's turn, and
\item[(b)]
  it is Black's turn?
\end{itemize}
\end{question}
\end{example}
\index{piece!circle}
\index{capture!circle}

\begin{figure}[h!] 
\centering 
%\includegraphics[height=7.0cm,width=13.5cm]{rithmomachia-numerous-captures-with-pyramid.jpg}
\begin{tabular}{|c|c|c|c|c|c|c|c||c} \hline
\circletfill{\,}7 & \quad  & \quad  & \quad  & \quad  & \quad  & \quad  & \trianglepa{\,}42 & $p$ \\ \hline
\quad  & \quad  & \quad  & \quad  & \quad  & \quad  & \quad  & \quad  & $o$ \\ \hline
\quad  & \quad  & \quad  & \squadfill{\,}361 & \quad  & \trianglepa{\,}72 & \trianglepafill{\,}90 & \quad  & $n$ \\ \hline
\quad  & \quad  & \quad  & \quad  & \quad  & \quad  & \quad  & \quad  & $m$ \\ \hline
\quad  & \quad  & \quad  & \quad  & \quad  & \quad  & \quad  & \quad  & $l$ \\ \hline
\quad  & \quad  & \quad  & \squad{\,}289 & \squadfill{\,}49 & \quad  & \quad  & \quad  & $k$ \\ \hline
\quad  & \quad  & \quad  & \quad  & \quad  & \quad  & \quad  & \quad  & $j$ \\ \hline
\quad  & \quad  & \quad  & \quad  & \quad  & \quad  & \quad  & \quad  & $i$ \\ \hline
\quad  & \quad  & \quad  & \quad  & \squad{\,}169 & \quad  & \quad  & \squadfill{\,}120 & $h$ \\ \hline
\quad  & \quad  & \quad  & \quad  & \quad  & \quad  & \quad  & \quad  & $g$ \\ \hline
\quad  & \quad  & \quad  & \quad  & \quad  & \quad  & \quad  & \quad  & $f$ \\ \hline
\quad  & \quad  & \circlet{\,}4 & \quad  & \quad  & \quad  & \quad  & \quad  & $e$ \\ \hline
\quad  & \circlet{\,}4 & \circletfill{\,}7 & \circlet{\,}16 & \quad  & \quad  & \quad  & \quad  & $d$ \\ \hline
\quad  & \quad  & \circlet{\,}2 & \circletfill{\,}9 & \quad  & \quad  & \trianglepa{\,}9 & \quad  & $c$ \\ \hline
\quad  & \quad  & \quad  & \quad  & \quad  & \quad  & \quad  & \quad  & $b$ \\ \hline
\rhombus{\,}$P^{91}$ & \quad  & \quad  & \quad  & \quad  & \quad  & \circletfill{\,}9 & \quad  & $a$ \\ \hline
1 &  2 &  3 &  4 &  5 &  6 &  7 &  8  &  \\
\end{tabular}
\caption{\textcolor{gray}{All types of captures are illustrated in this rithmomachia study. Have fun finding them all! }}
\label{fig:rithmomachia-numerous-captures-with-pyramid} 
\end{figure}
%sage: GSp = board_initial_matrix(pyramid_decomposition = True)
%sage: GS = copy(GSp)
%sage: GS = GS*0
%sage: c,C,p,P,t,T,s,S = var("c,C,p,P,t,T,s,S")
%sage: GS[0,0] = P^91 + S^36 + S^25 + S^16 + S^9 + S^4 + S
%sage: GS[0,0] = P^(91) + S^(36) + S^(25) + S^(16) + S^9 + S^4 + S
%sage: GS[0, 15] = c^7; GS[7, 15] = T^(42)
%sage: GS[6, 2] = T^9; GS[6, 13] = t^(90); GS[6, 0] = c^9
%sage: GS[4, 7] = S^(169); GS[7, 7] = s^(120); GS[4, 10] = s^(49)
%sage: GS[3, 10] = S^(289); GS[3, 13] = s^(361); GS[5, 13] = T^(72)
%sage: GS[2, 2] = C^2; GS[2, 4] = C^4; GS[3, 3] = C^(16); GS[1, 3] = C^4; GS[2, 3] = c^7; GS[3, 2] = c^9
%%sage: pclst = game_state_to_piece_list(GS)
%%sage: piece_list_to_latex(pclst, verbose = False)


\begin{figure}[h!] 
\centering 
\includegraphics[height=4.0cm,width=14.0cm]{Rithmomachia-game-pieces/white-squares-red-border-in-2lines.jpg} 
\caption{Rithmomachia White square pieces}
\label{fig:white-squares-red-border-in-2lines} 
\end{figure}


\subsection{How a game is won}
\label{sec:winning}

As mentioned in \S \ref{sec:setup}, if it's a player's turn yet that player
has no legal move then the player loses that rithmomachia game.
There are several other ways to win a game of rithmomachia, broken down
below into two distinct types.
\index{stalemate}


\vskip.1in
\postersection{How to Win a Game}
A player loses when it is their turn and no legal move can be made, or when
\begin{itemize}
\item
{\it common}: The number and values of the captured pieces satisfy certain conditions.
\item
{\it proper}:
A certain type of configuration of at least 3 enemy pieces is arranged in friendly territory.
\end{itemize}
\postersectionend
\index{territory}

One type of win is called {\it common} (or {\it minor})
and the other type of win is called {\it proper} (or {\it major}).
As explained below, some types of victory discussed here
{\it require} White and
Black to agree upon a certain conventions before the game.
\index{victory!conventions}

\subsubsection{Common}

As mentioned, in addition to common, this is also called a minor victory,
because it is somewhat easier to manage.
\index{victory!common}

\begin{itemize}
\item
%body
Assume that White and Black previously agreed upon a value for the
number of captured pieces, say, $N_0$. As an example, de Boissi\`ere
gives $N_0=4$ (see page 204 of \cite{Ri46}). Another source suggests
a larger proportion of the pieces, say $50\%$ (so $N_0 = 12$).

The first player to capture $N_0$ of the enemy's pieces (irregardless of their
values) is declared the {\it winner by body}
(sometimes called a {\it de Bonis victory}).
\index{victory!common!of the body}

\item
%goods
Assume that White and Black previously agreed upon a number for the
total value of captured pieces, say, $N_1$. For example, de Boissi\`ere
gives $N_1=100$ (see page 205 of \cite{Ri46}). Another source suggests
a larger proportion of the total value, say $50 \%$.
Recall from Figures \ref{fig:whitepieces} and \ref{fig:blackpieces},
the total value of the pieces depends
on whether they are Odd/Black or Even/White pieces,
but in any case that total is in the range $1300$-to-$1800$.

The first player to capture enemy's pieces whose total
equals
or exceeds $N_1$ is declared the {\it winner by goods}
(also called a {\it de Lite victory}).
Another alternative is for White and Black to agree
that such a victory only occurs is $N_1$ is attained exactly by
a combination of captured enemy piece values.
\index{victory!common! of the goods}

\item
%law-suit
Assume that White and Black previously agreed upon two
numbers, $N_2$ and $N_3$, for the total value of a number of
captured pieces. As an example, de Boissi\`ere
gives $N_2=4$ and $N_3=100$ (see page 205 of \cite{Ri46}).

The first player to capture at least $N_2$ enemy's pieces whose total
equals or exceeds $N_3$ is declared the {\it winner by law-suit}.
The winner can capture more than $N_2$ pieces but of them,
there are $N_2$ captured enemy pieces whose total is
at least $N_3$.
\index{victory!common! of the law-suit}

\item
%honor
This type of victory is similar to the
law-suit, but the
total value must be attained exactly. 
The first player to capture $N_2$ enemy's pieces whose total
equals $N_3$ is declared the {\it winner by honor}.
The winner can capture more than $N_3$ pieces but of them,
there are $N_2$ captured enemy pieces whose total is exactly $N_3$.
\index{victory!common! of honor}

\item
%law-suit/honor
This type of victory is similar to the honor and the 
law-suit, but the total number of captures and the
total value must both be attained exactly. 

We use the same notation.
The first player to capture $N_2$ enemy's pieces (no more, no less)
whose total equals $N_3$ is declared the {\it winner by honor and law-suit}.
To win by this method, the player must capture exactly $N_2$
enemy pieces and their total must be exactly $N_3$, no more
and no less.
\index{victory!common! of honor and law-suit}


\end{itemize}
These are all the common victories.


%{\footnotesize{
%\begin{center}
%\begin{tabular}{|c|c|} \hline  
%common type   & description \\ \hline    
%by body  & first player to capture more than $N_0$ enemy pieces wins \\
%by goods  & first player to capture enemy pieces, total value exceeds $N_1$ wins \\
%by lawsuit & capture more than $N_0$ enemy pieces, total value exceeds $N_1$ \\
%by honor & capture $N_0$ enemy pieces, total value equals $N_1$ \\ \hline
%\end{tabular}
%\vskip .1in
%{\bf Table: piece capture types}
%\end{center}
%}}

\vskip .1in
\postersection{Common victories:}
\begin{itemize}%[leftmargin=*, label=\textcolor{RithmoGold}{\faTrophy}]
    \item \textbf{By Body:} First to capture an agreed number of enemy pieces (e.g., $N_0 = 4$).
    \item \textbf{By Goods:} First to capture enemy pieces whose total value equals or exceeds an agreed sum (e.g., $N_1 = 100$).
    \item \textbf{By Law-suit:} Capture $\ge N_2$ pieces whose total value is $\ge N_3$ (e.g., $N_2=4, N_3=100$).
    \item \textbf{By Honor:} Capture $N_2$ pieces whose total value is exactly $N_3$.
\end{itemize}
\postersectionend


\subsubsection{Proper}


\begin{quotation}
\dots this is called a proper victory with good reason.
For just as those who make war win a victory after routing  
the enemy's battle-line and arranging everything in
accordance with their own will, so too the harmony
and arrangement of numbers \dots

\hskip 1in
Claude de Boissi\`ere (\cite{Ri46}, pages 206-207)
\end{quotation}
\index{battlefield harmony}
\index{de Boissi\`ere, Claude}

\vskip .2in

Each of these proper victories involve
placing a fixed number of friendly pieces in the enemy
territory such that the values of these pieces form
eith an arithmetical or geometrical or harmonic
progression. We will call any such progressions a
{\it coordinating progression}. Such progressions were known to
Pythagoras and covered in the book by Nicomachus \cite{Ni00}. 
\index{victory!proper}
\index{coordinating progression}
\index{progression!coordinating}

There are three types of proper wins
(see \cite{Ri46} pages 206-212, \cite{Ne98}, \cite{DG21}.

\begin{itemize}
\item
{\it Small},
\index{victory!proper! small}
  
\item
{\it Big},
\index{victory!proper! big}

\item
{\it Great}.
\index{victory!proper! great}
\end{itemize}
%A quick glance at the literature in the bibliography shows
%the variation of alternative names for these wins.
Each of these win types require a certain (coordinating) arrangement
of $3$ or $4$ pieces, the arrangement being a pattern of $3$ or $4$
values of pieces in enemy territory must all belong to either an
{\it arithmetic} progression, a {\it geometric} progression, or a
{\it harmonic} progression. Definitions and a examples of some of these can be found
in Figure \ref{fig:great-win-harmonies}.
\index{harmonic progression}

Since each side has only about $20$ or so values amongest their pieces,
there are lots of $3$-term or $4$-term arithmetic progressions
which {\it aren't} formed from these values. We used the programming
language Python to extract the $3$-term arithmetic/geometric/harmonic progressions
which {\it are} formed from three piece values.
The algorithm used did not distinguish between pieces {\it themselves} but rather
between the values of the pieces. Of course, the values of pieces
can repeat (like Odd/Black's {\tt c9}s and Even/White's {\tt T6}
and {\tt C6}). These lower bounds are woven in the discussion below.

\begin{itemize}
\item
An {\it arithmetical progression} is a sequence $a,b,c$ satisfying $b-a=c-b$.
%The differences are (reading from top to bottom) as follows,
%\[
%1, 2, 4, 20, 30, 57, 72, 105, 39.
%\]
\index{arithmetical progression}

\begin{itemize}
\item
 White has at least $11$ such progressions:

\[
\{(2, 4, 6), (2, 9, 16), (4, 6, 8), (4, 20, 36), (8, 25, 42), (8, 36, 64),
\]
\[
(9, 45, 81), (9, 81, 153), (15, 20, 25), (20, 42, 64), (49, 169, 289)\},
\]
\item
Black has at least $9$ such progressions:

\[
\{(3, 5, 7), (5, 7, 9), (7, 16, 25), (7, 28, 49), (7, 64, 121),
\]
\[
(12, 56, 100), (12, 66, 120), (16, 36, 56), (28, 64, 100)
\}.
\]
\end{itemize}

%Recall from the above:
%White/Even has 15 repeated values, while Black/Odd has 13.

\begin{question}
Accounting for the repeated values can you determine the
{\it exact} number of arithmetical progressions possible for
White? For Black?
\end{question}

\item
A {\it geometric progression} is a sequence $a,b,c$ 
satisfying $b/a=c/b$.
\index{geometrical progression}

For example, White has at least $11$ such progressions:

\[
\{
(2, 4, 8), (4, 6, 9), (4, 8, 16), (4, 16, 64), (9, 15, 25), (16, 20, 25),
\]
\[
 (16, 36, 81), (25, 45, 81), (36, 42, 49), (64, 72, 81), (81, 153, 289)
\}.
\]
Black has at least $10$ such progressions:

\[
\{
(9, 12, 16), (9, 30, 100), (16, 28, 49), (16, 36, 81), (25, 30, 36),
\]
\[
(36, 66, 121), (36, 90, 225), (49, 56, 64), (64, 120, 225), (81, 90, 100)
\}.
\]

\begin{question}
Accounting for the duplicate piece values can you determine the
{\it exact} number of geometrical progressions possible for
White? For Black?
\end{question}

\item
The {\it harmonic progressions} $a,b,c$ in each row of the harmonic table
satisfy $1/a - 1/b = 1/b - 1/c$, so their reciprocals form an
arithmetical sequence.

White has $2$ of them $(9, 15, 45)$ and $(9, 16, 72)$,
while Black has none.
\index{harmonic!progression}

\begin{question}
Can you determine the smallest number of White must make in a game
of rithmomachia to achieve a small proper win by harmonic progression
with the
sequence $9,15,45$?
\end{question}
\end{itemize}

\vskip .1in
\postersection{Offensive capabilities}
Each player who wants to achieve such a small proper win will
naturally want to have as many or more triples of pieces
which form a coordinating progression. To win, the pieces in
such an ``offensive'' triple (a {\it fireteam}, in the battlefield lingo)
\index{fireteam}
must all invade enemy territory.
The White/Even player has $60$ pairs of values which belong
to at least one coordinating progression. For example, each of the
pairs $(4, 6)$ and $(20, 25)$ belong to at least {\it two}
coordinating progressions.
The Black/Odd player has $57$ pairs of values which belong
to at least one coordinating progression. For example, each of the
pairs $(5, 7)$ and $(28, 49)$ belong to at least {\it two}
coordinating progressions.
From this perspective, the two sides seem to have roughly
equal offensive capabilities.
\index{offensive triple}
\postersectionend

%\begin{center}
\begin{figure}
\begin{tabular}{c|c|c}    
\begin{tabular}{ccc} \hline
   &  arithmetical   &    \\ \hline
2     &     4        &  6 \\ 
{\bf 3}  &  {\bf 5}    &  {\bf 7} \\   \hline
\end{tabular}
&
\begin{tabular}{cccc} \hline
&  geometric   &      &     \\ \hline
2              &  4   &   8 \\ 
{\bf 100}      & {\bf 190}  & {\bf 361} \\ \hline
\end{tabular}
&
\begin{tabular}{cccc} \hline
 & harmonic   &        &     \\ \hline
9            &   15    &  45 \\ 
9      &   16    &  72 \\   \hline
\end{tabular}
\end{tabular}
\caption{\textcolor{gray}{Example coordinating configurations for a {\it proper win}.}}
\label{fig:great-win-harmonies}
\end{figure}
%\end{center}

According to de Boissi\`ere (see \cite{Ri46}, page 206),
such coordinating progressions of pieces are a metaphor for the
cooperation and coordination of the corresponding fireteam on
the battlefield.
\index{battlefield harmony}

\vskip.1in
\postersection{Proper victories:}
%Require placing friendly pieces in enemy territory in an arithmetical/geometrical/harmonic progression. No linear arrangement of the pieces is required in de Boissière's version.
\begin{itemize}%[leftmargin=*, label=\textcolor{RithmoGold}{\faStar}]
\item
\textbf{Small victory:} Place \textbf{three} friendly pieces
in enemy territory whose values $a,b,c$ form an arithmetic, geometric,
or harmonic progression.
  %  \begin{itemize}
  %      \item Arithmetic: $a, b, c$ 
  %      \item Geometric: $a, b, c$ 
  %      \item Harmonic: $a, b, c$ 
  %  \end{itemize}
\item
\textbf{Big victory:} Four pieces in enemy territory; one subset
of three forms a coordinating progression, and another subset of
three forms another.

\item
\textbf{Great victory:} Among the four pieces, there are subsets
of $3$ that form arithmetic, geometric, and harmonic progressions.
\end{itemize}
%Achieving these "harmonies" reflected divine proportions.
\postersectionend
\index{victory!small}
\index{victory!big}
\index{victory!great}

In other words, a player wins by a great victory when they have a big win
but, moreover, there exists among the four pieces in enemy territory,

\begin{itemize}
\item
a subset of three of the four whose values form an
arithmetic progression, 
\item
a subset of three of the four whose values form a
geometric progression, and
\item
a subset of three of the four whose values form a
harmonic progression.
\end{itemize}

\begin{remark}
{\rm  
In some versions of rithmomachia, 
these three pieces must be arranged in a straight line
(vertical, horizontal, or diagonal).
Other versions of rithmomachia allow these pieces to be
arranged in the shape of a right angled triangle.
No such linear or geometric conditions exist for the de
Boissi\`ere version here.
}
\end{remark}
\index{linear condition, none}


%{\small{
%\begin{center}
%\begin{tabular}{|c|c|}  \hline
%progression & small proper victory description \\ \hline
%arithmetical   & 3 friendly pieces in enemy terr, values form an arith. prog.\\
%geometrical  & 3 friendly pieces in enemy terr, values form a geom. prog.\\
%harmonic      &  3 friendly pieces in enemy terr, values form a harmonic prog.\\ \hline
%\end{tabular}
%\vskip .1in
%       {\bf Table: Winning by a small proper victory}
%\end{center}
%}}


\noindent
{\it Example}:
Figure \ref{fig:rithmomachia-white-win-geometrical-harmony2} depicts an
example of White winning by geometrical progression.
White has three pieces, at $i1$, $i4$, $i7$,
%equally spaced along a line
in enemy territory with values $(2, 4, 8)$ in geometrical progression.
\index{victory!!proper!by arithmetic progression}
\index{victory!!proper!by geometric progression}
\index{victory!!proper!by harmonic progression}


\begin{figure}[h!] 
\centering 
%\includegraphics[height=7.0cm,width=13.5cm]{rithmomachia-white-win-geometrical-harmony2.jpg} 
\begin{tabular}{|c|c|c|c|c|c|c|c||c} \hline
\squadfill{\,}49 & \squadfill{\,}121 & \quad  & \quad  & \quad  & \quad  & \squadfill{\,}225 & \squadfill{\,}361 & $p$ \\ \hline
\squadfill{\,}28 & \squadfill{\,}66 & \trianglepafill{\,}36 & \trianglepafill{\,}30 & \trianglepafill{\,}56 & \trianglepafill{\,}64 & \squadfill{\,}120 & \rhombusfill{\,}$p^{190}$ & $o$ \\ \hline
\trianglepafill{\,}16 & \trianglepafill{\,}12 & \circletfill{\,}9 & \circletfill{\,}25 & \circletfill{\,}49 & \circletfill{\,}81 & \trianglepafill{\,}90 & \trianglepafill{\,}100 & $n$ \\ \hline
\quad  & \quad  & \circletfill{\,}3 & \circletfill{\,}5 & \circletfill{\,}7 & \circletfill{\,}9 & \quad  & \quad  & $m$ \\ \hline
\quad  & \quad  & \quad  & \quad  & \quad  & \quad  & \quad  & \quad  & $l$ \\ \hline
\quad  & \quad  & \quad  & \quad  & \quad  & \quad  & \quad  & \quad  & $k$ \\ \hline
\quad  & \quad  & \quad  & \quad  & \quad  & \quad  & \quad  & \quad  & $j$ \\ \hline
\circlet{\,}2 & \quad  & \quad  & \circlet{\,}4 & \quad  & \quad  & \circlet{\,}8 & \quad  & $i$ \\ \hline
\quad  & \quad  & \quad  & \quad  & \quad  & \quad  & \quad  & \quad  & $h$ \\ \hline
\quad  & \quad  & \quad  & \quad  & \quad  & \quad  & \quad  & \quad  & $g$ \\ \hline
\quad  & \quad  & \quad  & \quad  & \quad  & \quad  & \quad  & \quad  & $f$ \\ \hline
\quad  & \quad  & \quad  & \quad  & \quad  & \quad  & \quad  & \quad  & $e$ \\ \hline
\quad  & \quad  & \quad  & \circlet{\,}6 & \quad  & \quad  & \quad  & \quad  & $d$ \\ \hline
\trianglepa{\,}81 & \trianglepa{\,}72 & \circlet{\,}64 & \circlet{\,}36 & \circlet{\,}16 & \circlet{\,}4 & \trianglepa{\,}6 & \trianglepa{\,}9 & $c$ \\ \hline
\squad{\,}153 & \rhombus{\,}$P^{91}$ & \trianglepa{\,}49 & \trianglepa{\,}42 & \trianglepa{\,}20 & \trianglepa{\,}25 & \squad{\,}45 & \squad{\,}15 & $b$ \\ \hline
\squad{\,}289 & \squad{\,}169 & \quad  & \quad  & \quad  & \quad  & \squad{\,}81 & \squad{\,}25 & $a$ \\ \hline
1 &  2 &  3 &  4 &  5 &  6 &  7 &  8  &  \\
\end{tabular} 
\caption{\textcolor{gray}{A Rithmomachia position depicting a White win by geometrical harmony.}}
\label{fig:rithmomachia-white-win-geometrical-harmony2} 
\end{figure}
%sage: GSp = board_initial_matrix(pyramid_decomposition = True)
%sage: GS = copy(GSp)
%sage: GS[2,3] = 0; GS[6, 8] = C^8; GS[4,3] = 0; GS[3, 8] = C^4; GS[5,3] = 0; GS[0, 8] = C^2
%sage: board_plot(GS)
%%sage: pclst = game_state_to_piece_list(GS)
%%sage: piece_list_to_latex(pclst, verbose = False)
\index{geometrical progression}


\noindent
{\it Example}:
The game in \S \ref{sec:example-game} provides an example
of a small geometric (proper) victory of Even/White over Odd/Black.


\begin{question}
Can you find an example of a final position that achieves a big
victory winning position?
\end{question}
  

\begin{question}
Can you find an example of a final position that achieves a great
victory winning position?
\end{question}



\begin{figure}[h!] 
\centering 
\includegraphics[height=15.0cm,width=10.cm]{Rithmomachia-game-pieces/chatGPT_rithmomachia-logo5} 
\caption{Yeah, \rithmomachia!  \small{(thanks chatGPT:-)}}
\label{fig:chatGPT_rithmomachia-logo5}
\end{figure}


%\newpage


\section{Two example games}
\label{sec:example-game}

A couple of example games are given next.

\subsection{An arithmetical win}

In this case, it's easy to describe this game. White has a plan to place
three pieces in a coordinating progression. Black is moves more-or-less
randomly so has little chance of stopping this plan from being implemented.
\index{coordinating progression}

\newpage
%\begin{enumerate}
%\item

White plays first. After the first few moves,
the position of the board and moves were as follows. The last move played
is underlined.
\index{game state}

\vskip .1in

%{\footnotesize{
{\small{
\begin{tabular}{c}
The first few turns \\
\begin{tabular}{r|cc|cc}
 & Even moves &     & Odd moves             & \\    \hline 
1   &  {\tt C002d6e6} &  &  \underline{\tt s121p2p5} &  \\
\end{tabular}
\\
\begin{tabular}{|c|c|c|c|c|c|c|c||c} \hline
\squadfill{\,}49 & \quad & \quad & \quad & \squadfill{\,}121 & \quad & \squadfill{\,}225 & \squadfill{\,}361 & $p$ \\ \hline
\squadfill{\,}28 & \squadfill{\,}66 & \trianglepafill{\,}36 & \trianglepafill{\,}30 & \trianglepafill{\,}56 & \trianglepafill{\,}64 & \squadfill{\,}120 & \rhombusfill{\,}$p^{190}$ & $o$ \\ \hline
\trianglepafill{\,}16 & \trianglepafill{\,}12 & \circletfill{\,}9 & \circletfill{\,}25 & \circletfill{\,}49 & \circletfill{\,}81 & \trianglepafill{\,}90 & \trianglepafill{\,}100 & $n$ \\ \hline
\quad & \quad & \circletfill{\,}3 & \circletfill{\,}5 & \circletfill{\,}7 & \circletfill{\,}9 & \quad & \quad & $m$ \\ \hline
\quad & \quad & \quad & \quad & \quad & \quad & \quad & \quad & $l$ \\ \hline
\quad & \quad & \quad & \quad & \quad & \quad & \quad & \quad & $k$ \\ \hline
\quad & \quad & \quad & \quad & \quad & \quad & \quad & \quad & $j$ \\ \hline
\quad & \quad & \quad & \quad & \quad & \quad & \quad & \quad & $i$ \\ \hline
\quad & \quad & \quad & \quad & \quad & \quad & \quad & \quad & $h$ \\ \hline
\quad & \quad & \quad & \quad & \quad & \quad & \quad & \quad & $g$ \\ \hline
\quad & \quad & \quad & \quad & \quad & \quad & \quad & \quad & $f$ \\ \hline
\quad & \quad & \quad & \quad & \quad & \circlet{\,}2 & \quad & \quad & $e$ \\ \hline
\quad & \quad & \circlet{\,}8 & \circlet{\,}6 & \circlet{\,}4 & \quad & \quad & \quad & $d$ \\ \hline
\trianglepa{\,}81 & \trianglepa{\,}72 & \circlet{\,}64 & \circlet{\,}36 & \circlet{\,}16 & \circlet{\,}4 & \trianglepa{\,}6 & \trianglepa{\,}9 & $c$ \\ \hline
\squad{\,}153 & \rhombus{\,}$P^{91}$ & \trianglepa{\,}49 & \trianglepa{\,}42 &\trianglepa{\,}20 & \trianglepa{\,}25 & \squad{\,}45 & \squad{\,}15 & $b$ \\ \hline
\squad{\,}289 & \squad{\,}169 & \quad & \quad & \quad & \quad & \squad{\,}81 & \squad{\,}25 & $a$ \\ \hline
1 & 2 & 3 & 4 & 5 & 6 & 7 & 8 & \\
\end{tabular} 
\end{tabular}
}}

%\newpage

%\item  
%\noindent
%After a few more turns, the board and moves (there were no captures) were

{\small{
%{\footnotesize{
\begin{tabular}{c}
Round 2 game state \\
\begin{tabular}{r|cc|cc}
   & Even moves & captures   & odd moves & captures \\    \hline 
1. & C{\tt 002d6e6}  & \quad & {\tt s121p2p5} & \quad  \\
2. & {\tt C002e6f6} & \quad &  {\tt c007m5l5} & \quad  \\
\end{tabular}
\\
\begin{tabular}{|c|c|c|c|c|c|c|c||c} \hline
\squadfill{\,}49 & \quad & \quad & \quad & \squadfill{\,}121 & \quad & \squadfill{\,}225 & \squadfill{\,}361 & $p$ \\ \hline
\squadfill{\,}28 & \squadfill{\,}66 & \trianglepafill{\,}36 & \trianglepafill{\,}30 & \trianglepafill{\,}56 & \trianglepafill{\,}64 & \squadfill{\,}120 & \rhombusfill{\,}$p^{190}$ & $o$ \\ \hline
\trianglepafill{\,}16 & \trianglepafill{\,}12 & \circletfill{\,}9 & \circletfill{\,}25 & \circletfill{\,}49 & \circletfill{\,}81 & \trianglepafill{\,}90 & \trianglepafill{\,}100 & $n$ \\ \hline
\quad & \quad & \circletfill{\,}3 & \circletfill{\,}5 & \quad & \circletfill{\,}9 & \quad & \quad & $m$ \\ \hline
\quad & \quad & \quad & \quad & \circletfill{\,}7 & \quad & \quad & \quad & $l$ \\ \hline
\quad & \quad & \quad & \quad & \quad & \quad & \quad & \quad & $k$ \\ \hline
\quad & \quad & \quad & \quad & \quad & \quad & \quad & \quad & $j$ \\ \hline
\quad & \quad & \quad & \quad & \quad & \quad & \quad & \quad & $i$ \\ \hline
\quad & \quad & \quad & \quad & \quad & \quad & \quad & \quad & $h$ \\ \hline
\quad & \quad & \quad & \quad & \quad & \quad & \quad & \quad & $g$ \\ \hline
\quad & \quad & \quad & \quad & \quad & \circlet{\,}2 & \quad & \quad & $f$ \\ \hline
\quad & \quad & \quad & \quad & \quad & \quad & \quad & \quad & $e$ \\ \hline
\quad & \quad & \circlet{\,}8 & \circlet{\,}6 & \circlet{\,}4 & \quad & \quad & \quad & $d$ \\ \hline
\trianglepa{\,}81 & \trianglepa{\,}72 & \circlet{\,}64 & \circlet{\,}36 & \circlet{\,}16 & \circlet{\,}4 & \trianglepa{\,}6 & \trianglepa{\,}9 & $c$ \\ \hline
\squad{\,}153 & \rhombus{\,}$P^{91}$ & \trianglepa{\,}49 & \trianglepa{\,}42 & \trianglepa{\,}20 & \trianglepa{\,}25 & \squad{\,}45 & \squad{\,}15 & $b$ \\ \hline
\squad{\,}289 & \squad{\,}169 & \quad & \quad & \quad & \quad & \squad{\,}81 & \squad{\,}25 & $a$ \\ \hline
1 & 2 & 3 & 4 & 5 & 6 & 7 & 8 & \\
\end{tabular}
\end{tabular}
}}

%\newpage


%\noindent
%After a few more turns (again, there were no captures), the position was like this:


{\small{
%{\footnotesize{
\begin{tabular}{c}
Round 3 game state \\
\begin{tabular}{r|cc|cc}
& Even moves &     & odd moves &  \\    \hline  
1   &  {\tt C002d6e6} & \quad &  {\tt s121p2p5} & \quad \\
2   &  {\tt C002e6f6} & \quad &  {\tt c007m5l5} & \quad \\
3   &  {\tt C002f6g6} & \quad &  \underline{\tt c009m6m5} & \quad \\
\end{tabular}
\\
\begin{tabular}{|c|c|c|c|c|c|c|c||c} \hline
\squadfill{\,}49 & \quad & \quad & \quad & \squadfill{\,}121 & \quad & \squadfill{\,}225 & \squadfill{\,}361 & $p$ \\ \hline
\squadfill{\,}28 & \squadfill{\,}66 & \trianglepafill{\,}36 & \trianglepafill{\,}30 & \trianglepafill{\,}56 & \trianglepafill{\,}64 & \squadfill{\,}120 & \rhombusfill{\,}$p^{190}$ & $o$ \\ \hline
\trianglepafill{\,}16 & \trianglepafill{\,}12 & \circletfill{\,}9 & \circletfill{\,}25 & \circletfill{\,}49 & \circletfill{\,}81 & \trianglepafill{\,}90 & \trianglepafill{\,}100 & $n$ \\ \hline
\quad & \quad & \circletfill{\,}3 & \circletfill{\,}5 & \circletfill{\,}9 & \quad & \quad & \quad & $m$ \\ \hline
\quad & \quad & \quad & \quad & \circletfill{\,}7 & \quad & \quad & \quad & $l$ \\ \hline
\quad & \quad & \quad & \quad & \quad & \quad & \quad & \quad & $k$ \\ \hline
\quad & \quad & \quad & \quad & \quad & \quad & \quad & \quad & $j$ \\ \hline
\quad & \quad & \quad & \quad & \quad & \quad & \quad & \quad & $i$ \\ \hline
\quad & \quad & \quad & \quad & \quad & \quad & \quad & \quad & $h$ \\ \hline
\quad & \quad & \quad & \quad & \quad & \circlet{\,}2 & \quad & \quad & $g$ \\ \hline
\quad & \quad & \quad & \quad & \quad & \quad & \quad & \quad & $f$ \\ \hline
\quad & \quad & \quad & \quad & \quad & \quad & \quad & \quad & $e$ \\ \hline
\quad & \quad & \circlet{\,}8 & \circlet{\,}6 & \circlet{\,}4 & \quad & \quad & \quad & $d$ \\ \hline
\trianglepa{\,}81 & \trianglepa{\,}72 & \circlet{\,}64 & \circlet{\,}36 & \circlet{\,}16 & \circlet{\,}4 & \trianglepa{\,}6 & \trianglepa{\,}9 & $c$ \\ \hline
\squad{\,}153 & \rhombus{\,}$P^{91}$ & \trianglepa{\,}49 & \trianglepa{\,}42 & \trianglepa{\,}20 & \trianglepa{\,}25 & \squad{\,}45 & \squad{\,}15 & $b$ \\ \hline
\squad{\,}289 & \squad{\,}169 & \quad & \quad & \quad & \quad & \squad{\,}81 & \squad{\,}25 & $a$ \\ \hline
1 & 2 & 3 & 4 & 5 & 6 & 7 & 8 & \\
\end{tabular}
\end{tabular}
}}

%\newpage

%\item  
%\noindent
%After the next round of turns, call it round 4:


{\small{
%{\footnotesize{
\begin{tabular}{c}
Round 4 game state\\   
\begin{tabular}{r|cc|cc}
 & Even moves & captures   & odd moves & captures  \\    \hline 
1   &  {\tt C002d6e6} & \quad &  {\tt s121p2p5} & \quad \\
2   &  {\tt C002e6f6} & \quad &  {\tt c007m5l5} & \quad \\
3   &  {\tt C002f6g6} & \quad &  {\tt c009m6m5} & \quad \\
4   &  {\tt C002g6h6} & \quad &  \underline{\tt s121p5p2} & \quad \\
\end{tabular}
\\      
\begin{tabular}{|c|c|c|c|c|c|c|c||c} \hline
\squadfill{\,}49 & \squadfill{\,}121 & \quad & \quad & \quad & \quad & \squadfill{\,}225 & \squadfill{\,}361 & $p$ \\ \hline
\squadfill{\,}28 & \squadfill{\,}66 & \trianglepafill{\,}36 & \trianglepafill{\,}30 & \trianglepafill{\,}56 & \trianglepafill{\,}64 & \squadfill{\,}120 & \rhombusfill{\,}$p^{190}$ & $o$ \\ \hline
\trianglepafill{\,}16 & \trianglepafill{\,}12 & \circletfill{\,}9 & \circletfill{\,}25 & \circletfill{\,}49 & \circletfill{\,}81 & \trianglepafill{\,}90 & \trianglepafill{\,}100 & $n$ \\ \hline
\quad & \quad & \circletfill{\,}3 & \circletfill{\,}5 & \circletfill{\,}9 & \quad & \quad & \quad & $m$ \\ \hline
\quad & \quad & \quad & \quad & \circletfill{\,}7 & \quad & \quad & \quad & $l$ \\ \hline
\quad & \quad & \quad & \quad & \quad & \quad & \quad & \quad & $k$ \\ \hline
\quad & \quad & \quad & \quad & \quad & \quad & \quad & \quad & $j$ \\ \hline
\quad & \quad & \quad & \quad & \quad & \quad & \quad & \quad & $i$ \\ \hline
\quad & \quad & \quad & \quad & \quad & \circlet{\,}2 & \quad & \quad & $h$ \\ \hline
\quad & \quad & \quad & \quad & \quad & \quad & \quad & \quad & $g$ \\ \hline
\quad & \quad & \quad & \quad & \quad & \quad & \quad & \quad & $f$ \\ \hline
\quad & \quad & \quad & \quad & \quad & \quad & \quad & \quad & $e$ \\ \hline
\quad & \quad & \circlet{\,}8 & \circlet{\,}6 & \circlet{\,}4 & \quad & \quad & \quad & $d$ \\ \hline
\trianglepa{\,}81 & \trianglepa{\,}72 & \circlet{\,}64 & \circlet{\,}36 & \circlet{\,}16 & \circlet{\,}4 & \trianglepa{\,}6 & \trianglepa{\,}9 & $c$ \\ \hline
\squad{\,}153 & \rhombus{\,}$P^{91}$ & \trianglepa{\,}49 & \trianglepa{\,}42 & \trianglepa{\,}20 & \trianglepa{\,}25 & \squad{\,}45 & \squad{\,}15 & $b$ \\ \hline
\squad{\,}289 & \squad{\,}169 & \quad & \quad & \quad & \quad & \squad{\,}81 & \squad{\,}25 & $a$ \\ \hline
1 & 2 & 3 & 4 & 5 & 6 & 7 & 8 & \\
\end{tabular}
\end{tabular}
}}
%\newpage

%\item
\noindent
White's 2-circle {\tt C2} is finally in the desired position:



{\small{
%{\footnotesize{
\begin{tabular}{c}
Round 5 game state\\   
\begin{tabular}{r|cc|cc}
    & Even moves         & captures    & odd moves   &  captures \\    \hline 
1   &  {\tt C002d6e6} & \quad &  {\tt s121p2p5} & \quad \\
 \vdots  &  \vdots    & \quad &  \vdots        & \quad \\
%2   &  {\tt C002e6f6} & \quad &  {\tt c007m5l5} & \quad \\
%3   &  {\tt C002f6g6} & \quad &  {\tt c009m6m5} & \quad \\
4   &  {\tt C002g6h6} & \quad &  {\tt s121p5p2} & \quad \\
5   &  {\tt C002h6i6} & \quad &  \underline{\tt t090n7l7} & \quad \\
\end{tabular}
\\
\begin{tabular}{|c|c|c|c|c|c|c|c||c} \hline
\squadfill{\,}49 & \squadfill{\,}121 & \quad & \quad & \quad & \quad & \squadfill{\,}225 & \squadfill{\,}361 & $p$ \\ \hline
\squadfill{\,}28 & \squadfill{\,}66 & \trianglepafill{\,}36 & \trianglepafill{\,}30 &
\trianglepafill{\,}56 & \trianglepafill{\,}64 & \squadfill{\,}120 & \rhombusfill{\,}$p^{190}$ & $o$ \\ \hline
\trianglepafill{\,}16 & \trianglepafill{\,}12 & \circletfill{\,}9 & \circletfill{\,}25 & \circletfill{\,}49 & \circletfill{\,}81 & \quad & \trianglepafill{\,}100 & $n$ \\ \hline
\quad & \quad & \circletfill{\,}3 & \circletfill{\,}5 & \circletfill{\,}9 & \quad & \quad & \quad & $m$ \\ \hline
\quad & \quad & \quad & \quad & \circletfill{\,}7 & \quad & \trianglepafill{\,}90 & \quad & $l$ \\ \hline
\quad & \quad & \quad & \quad & \quad & \quad & \quad & \quad & $k$ \\ \hline
\quad & \quad & \quad & \quad & \quad & \quad & \quad & \quad & $j$ \\ \hline
\quad & \quad & \quad & \quad & \quad & \circlet{\,}2 & \quad & \quad & $i$ \\ \hline
\quad & \quad & \quad & \quad & \quad & \quad & \quad & \quad & $h$ \\ \hline
\quad & \quad & \quad & \quad & \quad & \quad & \quad & \quad & $g$ \\ \hline
\quad & \quad & \quad & \quad & \quad & \quad & \quad & \quad & $f$ \\ \hline
\quad & \quad & \quad & \quad & \quad & \quad & \quad & \quad & $e$ \\ \hline
\quad & \quad & \circlet{\,}8 & \circlet{\,}6 & \circlet{\,}4 & \quad & \quad & \quad & $d$ \\ \hline
\trianglepa{\,}81 & \trianglepa{\,}72 & \circlet{\,}64 & \circlet{\,}36 & \circlet{\,}16 & \circlet{\,}4 & \trianglepa{\,}6 & \trianglepa{\,}9 & $c$ \\ \hline
\squad{\,}153 & \rhombus{\,}$P^{91}$ & \trianglepa{\,}49 & \trianglepa{\,}42 & \trianglepa{\,}20 & \trianglepa{\,}25 & \squad{\,}45 & \squad{\,}15 & $b$ \\ \hline
\squad{\,}289 & \squad{\,}169 & \quad & \quad & \quad & \quad & \squad{\,}81 & \squad{\,}25 & $a$ \\ \hline
1 & 2 & 3 & 4 & 5 & 6 & 7 & 8 & \\
\end{tabular}
\end{tabular}
}}

%\newpage

\vskip .1in


{\small{
%{\footnotesize{
\begin{tabular}{c}
Round 6 game state \\      
\begin{tabular}{r|cc|cc}
    & Even moves     &  captures     & odd moves    & captures \\    \hline 
1   &  {\tt C002d6e6} & \quad &  {\tt s121p2p5} & \quad \\
\vdots &  \vdots    & \vdots  & \vdots    & \vdots \\  
%4   &  {\tt C002g6h6} & \quad &  {\tt s121p5p2} & \quad \\
5   &  {\tt C002h6i6} & \quad &  {\tt t090n7l7} & \quad \\
6   &  {\tt C004d5e5} & \quad &   \underline{\tt c005m4l4}  & \quad \\
\end{tabular}
\\
\begin{tabular}{|c|c|c|c|c|c|c|c||c} \hline
\squadfill{\,}49 & \squadfill{\,}121 & \quad & \quad & \quad & \quad & \squadfill{\,}225 & \squadfill{\,}361 & $p$ \\ \hline
\squadfill{\,}28 & \squadfill{\,}66 & \trianglepafill{\,}36 & \trianglepafill{\,}30 & \trianglepafill{\,}56 & \trianglepafill{\,}64 & \squadfill{\,}120 & \rhombusfill{\,}$p^{190}$ & $o$ \\ \hline
\trianglepafill{\,}16 & \trianglepafill{\,}12 & \circletfill{\,}9 & \circletfill{\,}25 & \circletfill{\,}49 & \circletfill{\,}81 & \quad & \trianglepafill{\,}100 & $n$ \\ \hline
\quad & \quad & \circletfill{\,}3 & \quad & \circletfill{\,}9 & \quad & \quad & \quad & $m$ \\ \hline
\quad & \quad & \quad & \circletfill{\,}5 & \circletfill{\,}7 & \quad & \trianglepafill{\,}90 & \quad & $l$ \\ \hline
\quad & \quad & \quad & \quad & \quad & \quad & \quad & \quad & $k$ \\ \hline
\quad & \quad & \quad & \quad & \quad & \quad & \quad & \quad & $j$ \\ \hline
\quad & \quad & \quad & \quad & \quad & \circlet{\,}2 & \quad & \quad & $i$ \\ \hline
\quad & \quad & \quad & \quad & \quad & \quad & \quad & \quad & $h$ \\ \hline
\quad & \quad & \quad & \quad & \quad & \quad & \quad & \quad & $g$ \\ \hline
\quad & \quad & \quad & \quad & \quad & \quad & \quad & \quad & $f$ \\ \hline
\quad & \quad & \quad & \quad & \circlet{\,}4 & \quad & \quad & \quad & $e$ \\ \hline
\quad & \quad & \circlet{\,}8 & \circlet{\,}6 & \quad & \quad & \quad & \quad & $d$ \\ \hline
\trianglepa{\,}81 & \trianglepa{\,}72 & \circlet{\,}64 & \circlet{\,}36 & \circlet{\,}16 & \circlet{\,}4 & \trianglepa{\,}6 & \trianglepa{\,}9 & $c$ \\ \hline
\squad{\,}153 & \rhombus{\,}$P^{91}$ & \trianglepa{\,}49 & \trianglepa{\,}42 &
\trianglepa{\,}20 & \trianglepa{\,}25 & \squad{\,}45 & \squad{\,}15 & $b$ \\ \hline
\squad{\,}289 & \squad{\,}169 & \quad & \quad & \quad & \quad & \squad{\,}81 & \squad{\,}25 & $a$ \\ \hline
1 & 2 & 3 & 4 & 5 & 6 & 7 & 8 & \\
\end{tabular}
\end{tabular}
}}

%\newpage

\vskip .1in



{\small{
%{\footnotesize{
\begin{tabular}{c}
Round 7 game state\\      
\begin{tabular}{r|cc|cc}
 & Even moves & captures   & odd moves & captures \\ \hline 
1   &  {\tt C002d6e6} & \quad &  {\tt s121p2p5} & \quad \\
\vdots &  \vdots    & \vdots  & \vdots    & \vdots \\  
%5   &  {\tt C002h6i6} & \quad &  {\tt t090n7l7} & \quad \\
6   &  {\tt C004d5e5} & \quad &   {\tt c005m4l4}  & \quad \\
7   &  {\tt C004e5f5} & \quad &   {\tt s225p7p4} &  \quad \\
\end{tabular}
\\
\begin{tabular}{|c|c|c|c|c|c|c|c||c} \hline
\squadfill{\,}49 & \squadfill{\,}121 & \quad & \squadfill{\,}225 & \quad & \quad & \quad & \squadfill{\,}361 & $p$ \\ \hline
\squadfill{\,}28 & \squadfill{\,}66 & \trianglepafill{\,}36 & \trianglepafill{\,}30 & \trianglepafill{\,}56 & \trianglepafill{\,}64 & \squadfill{\,}120 & \rhombusfill{\,}$p^{190}$ & $o$ \\ \hline
\trianglepafill{\,}16 & \trianglepafill{\,}12 & \circletfill{\,}9 & \circletfill{\,}25 & \circletfill{\,}49 & \circletfill{\,}81 & \quad & \trianglepafill{\,}100 & $n$ \\ \hline
\quad & \quad & \circletfill{\,}3 & \quad & \circletfill{\,}9 & \quad & \quad & \quad & $m$ \\ \hline
\quad & \quad & \quad & \circletfill{\,}5 & \circletfill{\,}7 & \quad & \trianglepafill{\,}90 & \quad & $l$ \\ \hline
\quad & \quad & \quad & \quad & \quad & \quad & \quad & \quad & $k$ \\ \hline
\quad & \quad & \quad & \quad & \quad & \quad & \quad & \quad & $j$ \\ \hline
\quad & \quad & \quad & \quad & \quad & \circlet{\,}2 & \quad & \quad & $i$ \\ \hline
\quad & \quad & \quad & \quad & \quad & \quad & \quad & \quad & $h$ \\ \hline
\quad & \quad & \quad & \quad & \quad & \quad & \quad & \quad & $g$ \\ \hline
\quad & \quad & \quad & \quad & \circlet{\,}4 & \quad & \quad & \quad & $f$ \\ \hline
\quad & \quad & \quad & \quad & \quad & \quad & \quad & \quad & $e$ \\ \hline
\quad & \quad & \circlet{\,}8 & \circlet{\,}6 & \quad & \quad & \quad & \quad & $d$ \\ \hline
\trianglepa{\,}81 & \trianglepa{\,}72 & \circlet{\,}64 & \circlet{\,}36 & \circlet{\,}16 & \circlet{\,}4 & \trianglepa{\,}6 & \trianglepa{\,}9 & $c$ \\ \hline
\squad{\,}153 & \rhombus{\,}$P^{91}$ & \trianglepa{\,}49 & \trianglepa{\,}42 & \trianglepa{\,}20 & \trianglepa{\,}25 & \squad{\,}45 & \squad{\,}15 & $b$ \\ \hline
\squad{\,}289 & \squad{\,}169 & \quad & \quad & \quad & \quad & \squad{\,}81 & \squad{\,}25 & $a$ \\ \hline
1 & 2 & 3 & 4 & 5 & 6 & 7 & 8 & \\
\end{tabular} 
\end{tabular}
}}
  
%\newpage

\vskip .1in


{\small{
%{\footnotesize{
\begin{tabular}{c}
Round 8 \\      
\begin{tabular}{r|cc|cc}
 & Even moves & captures   & odd moves & captures \\    \hline 
1   &  {\tt C002d6e6} & \quad &  {\tt s121p2p5} & \quad \\
\vdots &  \vdots    & \vdots  & \vdots    & \vdots \\  
%6   &  {\tt C004d5e5} & \quad &   {\tt c005m4l4}  & \quad \\
7   &  {\tt C004e5f5} & \quad &    {\tt s225p7p4} &  \quad \\
8   &  {\tt C004f5g5} & \quad &   {\tt c009m5m4} & \quad \\
\end{tabular}
\\
\begin{tabular}{|c|c|c|c|c|c|c|c||c} \hline
\squadfill{\,}49 & \squadfill{\,}121 & \quad & \squadfill{\,}225 & \quad & \quad & \quad & \squadfill{\,}361 & $p$ \\ \hline
\squadfill{\,}28 & \squadfill{\,}66 & \trianglepafill{\,}36 & \trianglepafill{\,}30 & \trianglepafill{\,}56 & \trianglepafill{\,}64 & \squadfill{\,}120 & \rhombusfill{\,}$p^{190}$ & $o$ \\ \hline
\trianglepafill{\,}16 & \trianglepafill{\,}12 & \circletfill{\,}9 & \circletfill{\,}25 & \circletfill{\,}49 & \circletfill{\,}81 & \quad & \trianglepafill{\,}100 & $n$ \\ \hline
\quad & \quad & \circletfill{\,}3 & \circletfill{\,}9 & \quad & \quad & \quad & \quad & $m$ \\ \hline
\quad & \quad & \quad & \circletfill{\,}5 & \circletfill{\,}7 & \quad & \trianglepafill{\,}90 & \quad & $l$ \\ \hline
\quad & \quad & \quad & \quad & \quad & \quad & \quad & \quad & $k$ \\ \hline
\quad & \quad & \quad & \quad & \quad & \quad & \quad & \quad & $j$ \\ \hline
\quad & \quad & \quad & \quad & \quad & \circlet{\,}2 & \quad & \quad & $i$ \\ \hline
\quad & \quad & \quad & \quad & \quad & \quad & \quad & \quad & $h$ \\ \hline
\quad & \quad & \quad & \quad & \circlet{\,}4 & \quad & \quad & \quad & $g$ \\ \hline
\quad & \quad & \quad & \quad & \quad & \quad & \quad & \quad & $f$ \\ \hline
\quad & \quad & \quad & \quad & \quad & \quad & \quad & \quad & $e$ \\ \hline
\quad & \quad & \circlet{\,}8 & \circlet{\,}6 & \quad & \quad & \quad & \quad & $d$ \\ \hline
\trianglepa{\,}81 & \trianglepa{\,}72 & \circlet{\,}64 & \circlet{\,}36 & \circlet{\,}16 & \circlet{\,}4 & \trianglepa{\,}6 & \trianglepa{\,}9 & $c$ \\ \hline
\squad{\,}153 & \rhombus{\,}$P^{91}$ & \trianglepa{\,}49 & \trianglepa{\,}42 & \trianglepa{\,}20 & \trianglepa{\,}25 & \squad{\,}45 & \squad{\,}15 & $b$ \\ \hline
\squad{\,}289 & \squad{\,}169 & \quad & \quad & \quad & \quad & \squad{\,}81 & \squad{\,}25 & $a$ \\ \hline
1 & 2 & 3 & 4 & 5 & 6 & 7 & 8 & \\
\end{tabular}
\end{tabular}
}}

%\newpage

\vskip .1in



{\small{
%{\footnotesize{
\begin{tabular}{c}
Round 9 \\      
\begin{tabular}{r|cc|cc}
 & Even moves & even captures    & odd moves & odd captures \\    \hline 
1   &  {\tt C002d6e6} & \quad &  {\tt s121p2p5} & \quad \\
\vdots &  \vdots    & \vdots  & \vdots    & \vdots \\  
8   &  {\tt C004f5g5} & \quad &  {\tt c009m5m4} & \quad \\
9   &  {\tt C004g5h5} & \quad &  {\tt c081n6n7} & \quad \\
\end{tabular}
\\
\begin{tabular}{|c|c|c|c|c|c|c|c||c} \hline
\squadfill{\,}49 & \squadfill{\,}121 & \quad & \squadfill{\,}225 & \quad & \quad & \quad & \squadfill{\,}361 & $p$ \\ \hline
\squadfill{\,}28 & \squadfill{\,}66 & \trianglepafill{\,}36 & \trianglepafill{\,}30 & \trianglepafill{\,}56 & \trianglepafill{\,}64 & \squadfill{\,}120 & \rhombusfill{\,}$p^{190}$ & $o$ \\ \hline
\trianglepafill{\,}16 & \trianglepafill{\,}12 & \circletfill{\,}9 & \circletfill{\,}25 & \circletfill{\,}49 & \quad & \circletfill{\,}81 & \trianglepafill{\,}100 & $n$ \\ \hline
\quad & \quad & \circletfill{\,}3 & \circletfill{\,}9 & \quad & \quad & \quad & \quad & $m$ \\ \hline
\quad & \quad & \quad & \circletfill{\,}5 & \circletfill{\,}7 & \quad & \trianglepafill{\,}90 & \quad & $l$ \\ \hline
\quad & \quad & \quad & \quad & \quad & \quad & \quad & \quad & $k$ \\ \hline
\quad & \quad & \quad & \quad & \quad & \quad & \quad & \quad & $j$ \\ \hline
\quad & \quad & \quad & \quad & \quad & \circlet{\,}2 & \quad & \quad & $i$ \\ \hline
\quad & \quad & \quad & \quad & \circlet{\,}4 & \quad & \quad & \quad & $h$ \\ \hline
\quad & \quad & \quad & \quad & \quad & \quad & \quad & \quad & $g$ \\ \hline
\quad & \quad & \quad & \quad & \quad & \quad & \quad & \quad & $f$ \\ \hline
\quad & \quad & \quad & \quad & \quad & \quad & \quad & \quad & $e$ \\ \hline
\quad & \quad & \circlet{\,}8 & \circlet{\,}6 & \quad & \quad & \quad & \quad & $d$ \\ \hline
\trianglepa{\,}81 & \trianglepa{\,}72 & \circlet{\,}64 & \circlet{\,}36 & \circlet{\,}16 & \circlet{\,}4 & \trianglepa{\,}6 & \trianglepa{\,}9 & $c$ \\ \hline
\squad{\,}153 & \rhombus{\,}$P^{91}$ & \trianglepa{\,}49 & \trianglepa{\,}42 & \trianglepa{\,}20 & \trianglepa{\,}25 & \squad{\,}45 & \squad{\,}15 & $b$ \\ \hline
\squad{\,}289 & \squad{\,}169 & \quad & \quad & \quad & \quad & \squad{\,}81 & \squad{\,}25 & $a$ \\ \hline
1 & 2 & 3 & 4 & 5 & 6 & 7 & 8 & \\
\end{tabular}
\end{tabular}
}}

%\newpage

\vskip .1in
 
\noindent
White/Even's 4-circle {\tt C4} is in enemy territory (in row $i$):

\vskip .1in

{\small{
%{\footnotesize{
\begin{tabular}{c}
Round 10 \\
\begin{tabular}{r|cc|cc}
 & Even moves &    & odd moves &  \\    \hline 
1   &  {\tt C002d6e6} & \quad &  {\tt s121p2p5} & \quad \\
\vdots &  \vdots    & \vdots  & \vdots    & \vdots \\  
%8   &  {\tt C004f5g5} & \quad &  {\tt c009m5m4} & \quad \\
9   &  {\tt C004g5h5} & \quad &  {\tt c081n6n7} & \quad \\
10  &  {\tt C004h5i5} & \quad &  {\tt c081n7n6} & \quad \\
\end{tabular}
\\
\begin{tabular}{|c|c|c|c|c|c|c|c||c} \hline
\squadfill{\,}49 & \squadfill{\,}121 & \quad & \squadfill{\,}225 & \quad & \quad & \quad & \squadfill{\,}361 & $p$ \\ \hline
\squadfill{\,}28 & \squadfill{\,}66 & \trianglepafill{\,}36 & \trianglepafill{\,}30 & \trianglepafill{\,}56 & \trianglepafill{\,}64 & \squadfill{\,}120 & \rhombusfill{\,}$p^{190}$ & $o$ \\ \hline
\trianglepafill{\,}16 & \trianglepafill{\,}12 & \circletfill{\,}9 & \circletfill{\,}25 & \circletfill{\,}49 & \circletfill{\,}81 & \quad & \trianglepafill{\,}100 & $n$ \\ \hline
\quad & \quad & \circletfill{\,}3 & \circletfill{\,}9 & \quad & \quad & \quad & \quad & $m$ \\ \hline
\quad & \quad & \quad & \circletfill{\,}5 & \circletfill{\,}7 & \quad & \trianglepafill{\,}90 & \quad & $l$ \\ \hline
\quad & \quad & \quad & \quad & \quad & \quad & \quad & \quad & $k$ \\ \hline
\quad & \quad & \quad & \quad & \quad & \quad & \quad & \quad & $j$ \\ \hline
\quad & \quad & \quad & \quad & \circlet{\,}4 & \circlet{\,}2 & \quad & \quad & $i$ \\ \hline
\quad & \quad & \quad & \quad & \quad & \quad & \quad & \quad & $h$ \\ \hline
\quad & \quad & \quad & \quad & \quad & \quad & \quad & \quad & $g$ \\ \hline
\quad & \quad & \quad & \quad & \quad & \quad & \quad & \quad & $f$ \\ \hline
\quad & \quad & \quad & \quad & \quad & \quad & \quad & \quad & $e$ \\ \hline
\quad & \quad & \circlet{\,}8 & \circlet{\,}6 & \quad & \quad & \quad & \quad & $d$ \\ \hline
\trianglepa{\,}81 & \trianglepa{\,}72 & \circlet{\,}64 & \circlet{\,}36 & \circlet{\,}16 & \circlet{\,}4 & \trianglepa{\,}6 & \trianglepa{\,}9 & $c$ \\ \hline
\squad{\,}153 & \rhombus{\,}$P^{91}$ & \trianglepa{\,}49 & \trianglepa{\,}42 & \trianglepa{\,}20 & \trianglepa{\,}25 & \squad{\,}45 & \squad{\,}15 & $b$ \\ \hline
\squad{\,}289 & \squad{\,}169 & \quad & \quad & \quad & \quad & \squad{\,}81 & \squad{\,}25 & $a$ \\ \hline
1 & 2 & 3 & 4 & 5 & 6 & 7 & 8 & \\
\end{tabular}
\end{tabular}
}}

%\newpage

\vskip .1in



{\small{
%{\footnotesize{
\begin{tabular}{c}
Round 11 game state \\
\begin{tabular}{r|cc|cc}
 & Even moves & even captures    & odd moves & odd captures \\    \hline 
1   &  {\tt C002d6e6} & \quad &  {\tt s121p2p5} & \quad \\
\vdots &  \vdots    & \vdots  & \vdots    & \vdots \\  
%9   &  {\tt C004g5h5} & \quad &  {\tt c081n6n7} & \quad \\
10  &  {\tt C004h5i5} & \quad &  {\tt c081n7n6} & \quad \\
11  &  {\tt C006d4e4} & \quad &  {\tt t090l7j7} & \quad \\
\end{tabular}
\\
\begin{tabular}{|c|c|c|c|c|c|c|c||c} \hline
\squadfill{\,}49 & \squadfill{\,}121 & \quad & \squadfill{\,}225 & \quad & \quad & \quad & \squadfill{\,}361 & $p$ \\ \hline
\squadfill{\,}28 & \squadfill{\,}66 & \trianglepafill{\,}36 & \trianglepafill{\,}30 & \trianglepafill{\,}56 & \trianglepafill{\,}64 & \squadfill{\,}120 & \rhombusfill{\,}$p^{190}$ & $o$ \\ \hline
\trianglepafill{\,}16 & \trianglepafill{\,}12 & \circletfill{\,}9 & \circletfill{\,}25 & \circletfill{\,}49 & \circletfill{\,}81 & \quad & \trianglepafill{\,}100 & $n$ \\ \hline
\quad & \quad & \circletfill{\,}3 & \circletfill{\,}9 & \quad & \quad & \quad & \quad & $m$ \\ \hline
\quad & \quad & \quad & \circletfill{\,}5 & \circletfill{\,}7 & \quad & \quad & \quad & $l$ \\ \hline
\quad & \quad & \quad & \quad & \quad & \quad & \quad & \quad & $k$ \\ \hline
\quad & \quad & \quad & \quad & \quad & \quad & \trianglepafill{\,}90 & \quad & $j$ \\ \hline
\quad & \quad & \quad & \quad & \circlet{\,}4 & \circlet{\,}2 & \quad & \quad & $i$ \\ \hline
\quad & \quad & \quad & \quad & \quad & \quad & \quad & \quad & $h$ \\ \hline
\quad & \quad & \quad & \quad & \quad & \quad & \quad & \quad & $g$ \\ \hline
\quad & \quad & \quad & \quad & \quad & \quad & \quad & \quad & $f$ \\ \hline
\quad & \quad & \quad & \circlet{\,}6 & \quad & \quad & \quad & \quad & $e$ \\ \hline
\quad & \quad & \circlet{\,}8 & \quad & \quad & \quad & \quad & \quad & $d$ \\ \hline
\trianglepa{\,}81 & \trianglepa{\,}72 & \circlet{\,}64 & \circlet{\,}36 & \circlet{\,}16 & \circlet{\,}4 & \trianglepa{\,}6 & \trianglepa{\,}9 & $c$ \\ \hline
\squad{\,}153 & \rhombus{\,}$P^{91}$ & \trianglepa{\,}49 & \trianglepa{\,}42 & \trianglepa{\,}20 & \trianglepa{\,}25 & \squad{\,}45 & \squad{\,}15 & $b$ \\ \hline
\squad{\,}289 & \squad{\,}169 & \quad & \quad & \quad & \quad & \squad{\,}81 & \squad{\,}25 & $a$ \\ \hline
1 & 2 & 3 & 4 & 5 & 6 & 7 & 8 & \\
\end{tabular}
\end{tabular}
}}

\newpage


\vskip .1in

\noindent
Presumably, White's strategy is clear now: to win by (small proper) arithmetical progression.


{\small{
%{\footnotesize{
\begin{tabular}{c}
Round 12 game state \\
\begin{tabular}{r|cc|cc}
 & Even moves & even captures    & odd moves & odd captures \\    \hline 
1   &  {\tt C002d6e6} & \quad &  {\tt s121p2p5} & \quad \\
\vdots &  \vdots    & \vdots  & \vdots    & \vdots \\  
%10  &  C004h5i5 & \quad &  c081n7n6 & \quad \\
11  &  {\tt C006d4e4} & \quad &  {\tt t090l7j7} & \quad \\
12  &  {\tt C006e4f4} & \quad &  {\tt c003m3l3} & \quad \\
\end{tabular}
\\
\begin{tabular}{|c|c|c|c|c|c|c|c||c} \hline
\squadfill{\,}49 & \squadfill{\,}121 & \quad & \squadfill{\,}225 & \quad & \quad & \quad & \squadfill{\,}361 & $p$ \\ \hline
\squadfill{\,}28 & \squadfill{\,}66 & \trianglepafill{\,}36 & \trianglepafill{\,}30 & \trianglepafill{\,}56 & \trianglepafill{\,}64 & \squadfill{\,}120 & \rhombusfill{\,}$p^{190}$ & $o$ \\ \hline
\trianglepafill{\,}16 & \trianglepafill{\,}12 & \circletfill{\,}9 & \circletfill{\,}25 & \circletfill{\,}49 & \circletfill{\,}81 & \quad & \trianglepafill{\,}100 & $n$ \\ \hline
\quad & \quad & \quad & \circletfill{\,}9 & \quad & \quad & \quad & \quad & $m$ \\ \hline
\quad & \quad & \circletfill{\,}3 & \circletfill{\,}5 & \circletfill{\,}7 & \quad & \quad & \quad & $l$ \\ \hline
\quad & \quad & \quad & \quad & \quad & \quad & \quad & \quad & $k$ \\ \hline
\quad & \quad & \quad & \quad & \quad & \quad & \trianglepafill{\,}90 & \quad & $j$ \\ \hline
\quad & \quad & \quad & \quad & \circlet{\,}4 & \circlet{\,}2 & \quad & \quad & $i$ \\ \hline
\quad & \quad & \quad & \quad & \quad & \quad & \quad & \quad & $h$ \\ \hline
\quad & \quad & \quad & \quad & \quad & \quad & \quad & \quad & $g$ \\ \hline
\quad & \quad & \quad & \circlet{\,}6 & \quad & \quad & \quad & \quad & $f$ \\ \hline
\quad & \quad & \quad & \quad & \quad & \quad & \quad & \quad & $e$ \\ \hline
\quad & \quad & \circlet{\,}8 & \quad & \quad & \quad & \quad & \quad & $d$ \\ \hline
\trianglepa{\,}81 & \trianglepa{\,}72 & \circlet{\,}64 & \circlet{\,}36 & \circlet{\,}16 & \circlet{\,}4 & \trianglepa{\,}6 & \trianglepa{\,}9 & $c$ \\ \hline
\squad{\,}153 & \rhombus{\,}$P^{91}$ & \trianglepa{\,}49 & \trianglepa{\,}42 & \trianglepa{\,}20 & \trianglepa{\,}25 & \squad{\,}45 & \squad{\,}15 & $b$ \\ \hline
\squad{\,}289 & \squad{\,}169 & \quad & \quad & \quad & \quad & \squad{\,}81 & \squad{\,}25 & $a$ \\ \hline
1 & 2 & 3 & 4 & 5 & 6 & 7 & 8 & \\
\end{tabular} 
\end{tabular}
}}

%\newpage
\vskip .1in

%\noindent
%Round 13:


{\small{
%{\footnotesize{
\begin{tabular}{c}
Round 13 game state \\
\begin{tabular}{r|cc|cc}
 & Even moves & even captures    & odd moves & odd captures \\    \hline 
1   &  {\tt C002d6e6} & \quad &  {\tt s121p2p5} & \quad \\
\vdots &  \vdots    & \vdots  & \vdots    & \vdots \\  
%11  &  {\tt C006d4e4} & \quad &  {\tt t090l7j7} & \quad \\
12  &  {\tt C006e4f4} & \quad &  {\tt c003m3l3} & \quad \\
13  &  {\tt C006f4g4} & \quad &  {\tt c009m4m3} & \quad \\
\end{tabular}
\\
\begin{tabular}{|c|c|c|c|c|c|c|c||c} \hline
\squadfill{\,}49 & \squadfill{\,}121 & \quad & \squadfill{\,}225 & \quad & \quad & \quad & \squadfill{\,}361 & $p$ \\ \hline
\squadfill{\,}28 & \squadfill{\,}66 & \trianglepafill{\,}36 & \trianglepafill{\,}30 & \trianglepafill{\,}56 & \trianglepafill{\,}64 & \squadfill{\,}120 & \rhombusfill{\,}$p^{190}$ & $o$ \\ \hline
\trianglepafill{\,}16 & \trianglepafill{\,}12 & \circletfill{\,}9 & \circletfill{\,}25 & \circletfill{\,}49 & \circletfill{\,}81 & \quad & \trianglepafill{\,}100 & $n$ \\ \hline
\quad & \quad & \circletfill{\,}9 & \quad & \quad & \quad & \quad & \quad & $m$ \\ \hline
\quad & \quad & \circletfill{\,}3 & \circletfill{\,}5 & \circletfill{\,}7 & \quad & \quad & \quad & $l$ \\ \hline
\quad & \quad & \quad & \quad & \quad & \quad & \quad & \quad & $k$ \\ \hline
\quad & \quad & \quad & \quad & \quad & \quad & \trianglepafill{\,}90 & \quad & $j$ \\ \hline
\quad & \quad & \quad & \quad & \circlet{\,}4 & \circlet{\,}2 & \quad & \quad & $i$ \\ \hline
\quad & \quad & \quad & \quad & \quad & \quad & \quad & \quad & $h$ \\ \hline
\quad & \quad & \quad & \circlet{\,}6 & \quad & \quad & \quad & \quad & $g$ \\ \hline
\quad & \quad & \quad & \quad & \quad & \quad & \quad & \quad & $f$ \\ \hline
\quad & \quad & \quad & \quad & \quad & \quad & \quad & \quad & $e$ \\ \hline
\quad & \quad & \circlet{\,}8 & \quad & \quad & \quad & \quad & \quad & $d$ \\ \hline
\trianglepa{\,}81 & \trianglepa{\,}72 & \circlet{\,}64 & \circlet{\,}36 & \circlet{\,}16 & \circlet{\,}4 & \trianglepa{\,}6 & \trianglepa{\,}9 & $c$ \\ \hline
\squad{\,}153 & \rhombus{\,}$P^{91}$ & \trianglepa{\,}49 & \trianglepa{\,}42 & \trianglepa{\,}20 & \trianglepa{\,}25 & \squad{\,}45 & \squad{\,}15 & $b$ \\ \hline
\squad{\,}289 & \squad{\,}169 & \quad & \quad & \quad & \quad & \squad{\,}81 & \squad{\,}25 & $a$ \\ \hline
1 & 2 & 3 & 4 & 5 & 6 & 7 & 8 & \\
\end{tabular}
\end{tabular}
}}

%\newpage

%\vskip .1in


\vskip .1in

{\small{
%{\footnotesize{
\begin{tabular}{c}
Round 14 game state \\
\begin{tabular}{r|cc|cc}
 & Even moves &    & odd moves &  \\    \hline 
1   &  {\tt C002d6e6} & \quad &  {\tt s121p2p5} & \quad \\
\vdots &  \vdots    & \vdots  & \vdots    & \vdots \\  
12  &  {\tt C006e4f4} & \quad &  {\tt c003m3l3} & \quad \\
13  &  {\tt C006f4g4} & \quad &  {\tt c009m4m3} & \quad \\
14  &  {\tt C006g4h4} & \quad & \underline{\tt c007l5k5} & \quad \\
\end{tabular}
\\
\begin{tabular}{|c|c|c|c|c|c|c|c||c} \hline
\squadfill{\,}49 & \squadfill{\,}121 & \quad & \squadfill{\,}225 & \quad & \quad & \quad & \squadfill{\,}361 & $p$ \\ \hline
\squadfill{\,}28 & \squadfill{\,}66 & \trianglepafill{\,}36 & \trianglepafill{\,}30 & \trianglepafill{\,}56 & \trianglepafill{\,}64 & \squadfill{\,}120 & \rhombusfill{\,}$p^{190}$ & $o$ \\ \hline
\trianglepafill{\,}16 & \trianglepafill{\,}12 & \circletfill{\,}9 & \circletfill{\,}25 & \circletfill{\,}49 & \circletfill{\,}81 & \quad & \trianglepafill{\,}100 & $n$ \\ \hline
\quad & \quad & \circletfill{\,}9 & \quad & \quad & \quad & \quad & \quad & $m$ \\ \hline
\quad & \quad & \circletfill{\,}3 & \circletfill{\,}5 & \quad & \quad & \quad & \quad & $l$ \\ \hline
\quad & \quad & \quad & \quad & \circletfill{\,}7 & \quad & \quad & \quad & $k$ \\ \hline
\quad & \quad & \quad & \quad & \quad & \quad & \trianglepafill{\,}90 & \quad & $j$ \\ \hline
\quad & \quad & \quad & \quad & \circlet{\,}4 & \circlet{\,}2 & \quad & \quad & $i$ \\ \hline
\quad & \quad & \quad & \circlet{\,}6 & \quad & \quad & \quad & \quad & $h$ \\ \hline
\quad & \quad & \quad & \quad & \quad & \quad & \quad & \quad & $g$ \\ \hline
\quad & \quad & \quad & \quad & \quad & \quad & \quad & \quad & $f$ \\ \hline
\quad & \quad & \quad & \quad & \quad & \quad & \quad & \quad & $e$ \\ \hline
\quad & \quad & \circlet{\,}8 & \quad & \quad & \quad & \quad & \quad & $d$ \\ \hline
\trianglepa{\,}81 & \trianglepa{\,}72 & \circlet{\,}64 & \circlet{\,}36 & \circlet{\,}16 & \circlet{\,}4 & \trianglepa{\,}6 & \trianglepa{\,}9 & $c$ \\ \hline
\squad{\,}153 & \rhombus{\,}$P^{91}$ & \trianglepa{\,}49 & \trianglepa{\,}42 & \trianglepa{\,}20 & \trianglepa{\,}25 & \squad{\,}45 & \squad{\,}15 & $b$ \\ \hline
\squad{\,}289 & \squad{\,}169 & \quad & \quad & \quad & \quad & \squad{\,}81 & \squad{\,}25 & $a$ \\ \hline
1 & 2 & 3 & 4 & 5 & 6 & 7 & 8 & \\
\end{tabular}
\end{tabular}
}}

%\newpage


\vskip .1in

\noindent
Round 15: The strategy for White/Even is very simple: march three
pieces to enemy territory. If their values form a coordinating progression
then White/Even can declare a small proper victory over
Black/Odd.

After 15 moves, finally White/Even now has an arithmetical progression
on row $j$, which is in enemy territory.
\index{arithmetical progression}

\vskip .1in
{\small{
\begin{center}
\begin{tabular}{r|cc|cc}
 & {Even moves} & even captures    & Odd moves & odd captures \\    \hline 
1   &  {\tt C002d6e6} & \quad &  {\tt s121p2p5} & \quad \\
2   &  {\tt C002e6f6} & \quad &  {\tt c007m5l5} & \quad \\
3   &  {\tt C002f6g6} & \quad &  {\tt c009m6m5} & \quad \\
4   &  {\tt C002g6h6} & \quad &  {\tt s121p5p2} & \quad \\
5   &  {\tt C002h6i6} & \quad &  {\tt t090n7l7} & \quad \\
6   &  {\tt C004d5e5} & \quad &  {\tt c005m4l4} & \quad \\
7   &  {\tt C004e5f5} & \quad &  {\tt s225p7p4} & \quad \\
8   &  {\tt C004f5g5} & \quad &  {\tt c009m5m4} & \quad \\
9   &  {\tt C004g5h5} & \quad &  {\tt c081n6n7} & \quad \\
10  &  {\tt C004h5i5} & \quad &  {\tt c081n7n6} & \quad \\
11  &  {\tt C006d4e4} & \quad &  {\tt t090l7j7} & \quad \\
12  &  {\tt C006e4f4} & \quad &  {\tt c003m3l3} & \quad \\
13  &  {\tt C006f4g4} & \quad &  {\tt c009m4m3} & \quad \\
14  &  {\tt C006g4h4} & \quad &  {\tt c007l5k5} & \quad \\
15  &  {\tt C006h4i4} & \quad &  1-0 & \quad \\
%15  &  C006h4i4 & \quad &  c007k5j5 & \quad \\
\end{tabular}
\end{center}
}}


\vskip .1in

\noindent
Here 1-0 means the first player (White or Even) has scored a win
(so earns the 1) over Black or Odd (who sadly scores 0).

%\newpage
game state:

\vskip .1in
{\footnotesize{
\begin{center}
\begin{tabular}{|c|c|c|c|c|c|c|c||c} \hline
\squadfill{\,}49 & \squadfill{\,}121 & \quad & \squadfill{\,}225 & \quad & \quad & \quad & \squadfill{\,}361 & $p$ \\ \hline
\squadfill{\,}28 & \squadfill{\,}66 & \trianglepafill{\,}36 & \trianglepafill{\,}30 & \trianglepafill{\,}56 & \trianglepafill{\,}64 & \squadfill{\,}120 & \rhombusfill{\,}$p^{190}$ & $o$ \\ \hline
\trianglepafill{\,}16 & \trianglepafill{\,}12 & \circletfill{\,}9 & \circletfill{\,}25 & \circletfill{\,}49 & \circletfill{\,}81 & \quad & \trianglepafill{\,}100 & $n$ \\ \hline
\quad & \quad & \circletfill{\,}9 & \quad & \quad & \quad & \quad & \quad & $m$ \\ \hline
\quad & \quad & \circletfill{\,}3 & \circletfill{\,}5 & \quad & \quad & \quad & \quad & $l$ \\ \hline
\quad & \quad & \quad & \quad & \circletfill{\,}7 & \quad & \quad & \quad & $k$ \\ \hline
\quad & \quad & \quad & \quad & \quad & \quad & \trianglepafill{\,}90 & \quad & $j$ \\ \hline
\quad & \quad & \quad & \circlet{\,}6 & \circlet{\,}4 & \circlet{\,}2 & \quad & \quad & $i$ \\ \hline
\quad & \quad & \quad & \quad & \quad & \quad & \quad & \quad & $h$ \\ \hline
\quad & \quad & \quad & \quad & \quad & \quad & \quad & \quad & $g$ \\ \hline
\quad & \quad & \quad & \quad & \quad & \quad & \quad & \quad & $f$ \\ \hline
\quad & \quad & \quad & \quad & \quad & \quad & \quad & \quad & $e$ \\ \hline
\quad & \quad & \circlet{\,}8 & \quad & \quad & \quad & \quad & \quad & $d$ \\ \hline
\trianglepa{\,}81 & \trianglepa{\,}72 & \circlet{\,}64 & \circlet{\,}36 & \circlet{\,}16 & \circlet{\,}4 & \trianglepa{\,}6 & \trianglepa{\,}9 & $c$ \\ \hline
\squad{\,}153 & \rhombus{\,}$P^{91}$ & \trianglepa{\,}49 & \trianglepa{\,}42 & \trianglepa{\,}20 & \trianglepa{\,}25 & \squad{\,}45 & \squad{\,}15 & $b$ \\ \hline
\squad{\,}289 & \squad{\,}169 & \quad & \quad & \quad & \quad & \squad{\,}81 & \squad{\,}25 & $a$ \\ \hline
1 & 2 & 3 & 4 & 5 & 6 & 7 & 8 & \\
\end{tabular}
\end{center}
}}
\index{victory!small, proper}
%\end{enumerate}

%The game is over once White can declare a small proper victory after
%playing {\tt C006h4i4}.


%\newpage


\begin{figure}[h!] 
\centering 
\includegraphics[height=2.0cm,width=15.0cm]{Rithmomachia-game-pieces/black-triangles-red-border-in-line} 
\caption{Rithmomachia Black/Odd triangles pieces in a line}
\label{fig:black-triangles-red-border-in-line}
\end{figure}

\vskip 3in

\subsection{A harmonic win}

The moves of the second game analyzed here is recorded as follows:


\[
\begin{array}{l|ll|ll|} \hline
  \quad & Even\ moves & captures & Odd\ moves & captures \\ \hline
  1 & T^9{\tt c8e8} &  & t^{90}{\tt n7l7} &   \\
  2 & T^6{\tt c7e7} &  & t^{90}{\tt l7j7} &   \\
  3 & T^9{\tt e8g8} &  & s^{120}{\tt o7l7} &   \\
  4 & S^{15}{\tt b8e8} &  & s^{225}{\tt p7m7} &   \\
  5 & T^9{\tt g8i8} &  & s^{121}{\tt p2p5} &   \\
  6 & T^6{\tt e7e5} &  & t^{16}{\tt n1l1} &   \\
  7 & S^{15}{\tt e8h8} &  & t^{16}{\tt l1j1} &   \\
  8 & T^6{\tt e5g5} &  & s^{28}{\tt o1l1} &   \\
  9 & S^{15}{\tt h8h5} &  & s^{49}{\tt p1m1} &   \\
  10 & S^{15}{\tt h5k5} &  & t^{90}{\tt j7h7} &   \\
  11 & S^{45}{\tt b7e7} & S^{45}\times t^{90} & s^{120}{\tt l7i7} &   \\
  12 & S^{45}{\tt e7h7} &  & s^{225}{\tt m7j7} &   \\
  13 & S^{45}{\tt h7h4} &  & s^{120}{\tt i7f7} &   \\
  14 & S^{45}{\tt h4k4} &  & 1-0 &   \\
  \quad & \quad & \quad & \quad & \quad \\ \hline
\end{array}
\]
The Python program for playing rithmomachia was revised since the
game in the previous section. Now Black's strategy is more
in line with defending against a possible White invasion.

After White's first move, the position is depicted here:

\begin{center}
    \adjustimage{max size={0.9\linewidth}{0.9\paperheight}}{rithmomachia-2025-05-27-move1w.png}
\end{center}
After the second Black move, the position is depicted here:

\begin{center}
    \adjustimage{max size={0.9\linewidth}{0.9\paperheight}}{rithmomachia-2025-05-27-move2b.png}
\end{center}
Imagine, in this position, that the triangle $T^6$ was moved to the side, so there is a
clear shot from $T^{45}$ to $t^{90}$. 
You can see that White can't move the $S^{45}$ piece ahead in the {\bf $225$}
file without fear of being captured (by division) by $t^{90}$ on the next move.
In some sense, Black's blockage using $t^{90}$ has successfully stymied White's plans for now.
However, White evades the blockage by
making a few moves to switch files and invade via the $42$ file.
The final position is here:
\begin{center}
    \adjustimage{max size={0.9\linewidth}{0.9\paperheight}}{rithmomachia-2025-05-27-move14w}
\end{center}


\section{Last but not least}


Rithmomachia is a board game emphasizing pattern recognition,
logical thinking, and numerical relationships between the pieces.
It has an interesting history, and presents us with
a pleasant diversion involving logic, strategy, twinged with philosophical
ideas.

\subsection{Open questions}

A few open questions remain.

\begin{question}
Assuming best play, does one side had an advantage over the other?
\end{question}  

There are arguments on both sides:
\begin{itemize}
\item  
We know that Odd/Black has higher total piece value ($1752$)
than Even/White ($1312$), so Odd/Back has the advantage from this
point of view.
\item
  On the other hand,
\begin{itemize}
\item[(a)]
there are at most $118$ ways for White to capture Black
by addition, and 
\item[(b)] there are at most $141$ ways for White to
  capture Black by subtraction.

\vskip .1in
\noindent
Compare that with: 
\item[(c)] there are at most $82$ ways for
Black to capture White by addition, and 
\item[(d)] there are at most $119$ ways for
  Black to capture White by subtraction.
\end{itemize}
So it seems Even/White has the advantage from this perspective.
\end{itemize}

\vskip .1in
\postersection{Defense}
As we've seen in the example games, a proper win can be quickly
achieved by invading the enemy territory with a fireteam.
To counter White's infiltration strategy, Black can perform this checklist:

\begin{itemize}
\item
scan for two White pieces in Black territory that are part of a
potential coordinating sequence.
\item
Identify if a third White piece can be moved into
Black territory to complete this sequence.
\item
  If such a threat exists, Black must either
  (a) capture one of White's two pieces already in the formation,
  (b) block the path of the incoming third White piece.
\end{itemize}
Easier said than done!
\postersectionend

\vskip .1in
As far as we know, this next one is also an open question.

\begin{question}
Will Odd/Black always win with perfect play from both sides in rithmomachia? 
\end{question}



\begin{figure}[h!] 
\centering 
\includegraphics[height=10.5cm,width=8.5cm]{Rithmomachia-game-pieces/Gemini_rithmomachia-battle6.png} 
\caption{\rithmomachia pieces can get mad too. 
  Watch out especially for that hot head \circlet{\,}{\tt 8}! (Thanks gemini:-)}
\label{fig:Gemini_rithmomachia-battle6} 
\end{figure}

\subsection{Summary of Rithmomachia Rules and Setup}

\begin{itemize}
    \item \textbf{Board:} The game is played on an 8-by-16 rectangular board. 

    \item \textbf{Players and Pieces:} Each player controls 24 pieces (white/even and black/odd), comprising:
      8 \emph{circles}, 8 \emph{triangles}, and 8 \emph{squares}. One of these squares is a
      special multi-valued piece called a {\it pyramid}. 

    \item \textbf{Movement:}
    \begin{itemize}
        \item Circles move orthogonally one.
        \item Triangles move exactly two spaces orthogonally.
        \item Squares and pyramids move exactly three spaces orthogonally.
        \item No piece may jump over another; movement must be unobstructed.
    \end{itemize}

    \item \textbf{Capturing:}
    \begin{itemize}
        \item \emph{capture by numbering:} The attacking piece is one move away.
        \item \emph{capture by addition/subtraction:} One move away for each of the attacking pieces.
        \item \emph{capture by multiplication/division:} Attacking piece at distance $d$, if its value
          equals $d$ (divided by/times) the enemy’s value.
        \item \emph{Siege capture:} An enemy piece is flanked on all sides by friendly pieces.
    \end{itemize}

    \item \textbf{Victory Conditions:}
    \begin{itemize}
        \item \textbf{Proper victory:} 
          Three of your pieces in arithmetic, geometric, or harmonic progression
          on the opponent's side of the board.
        \item \textbf{Common victory:} Some pre-agreed proportion of
          your opponents pieces (by count or by total value).
\end{itemize}
           

    \item \textbf{Philosophical and Mathematical Themes:}
    \begin{itemize}
        \item Game emphasizes harmony, proportion, and the balance of opposites (odd/even).
        \item Influenced by Pythagorean number mysticism and Boethian arithmetic.
        \item Intended as a didactic tool for exploring numerical relationships and moral virtues.
    \end{itemize}
\end{itemize}


\subsection{Play a game!}

If you're still curious, then play a game.
A board is simply an $8\times 16$ grid of squares -- easy to
make from cardboard and even two chess/checker boards
side-by-side will work.
Pieces aren't hard to make with cardboard, scissors, and a
couple of different colored pens/crayons.
%The board can be made by taping together a few pieces of graph paper.

%
%\begin{center}
%    \adjustimage{max size={0.9\linewidth}{0.8\paperheight}}{Rithmomachia-game-pieces/rithmomachia-set7-PXL_20250530b}
%\end{center}

Hopefully, you'll have a little fun bringing this forgotten
game back to life.  If finding a human opponent is challenging,
consider exploring computer versions, though ({\it caveat emptor})
be mindful that rule variations may exist.

You can even reach out to the
author for a free copy of their Python/{\tt SageMath}
program to get started. The good news is that it's free, while the
bad news is that there is a learning curve to installing and
using the mathematical software program {\tt SageMath} that my
program currently relies on.
The example in \S \ref{sec:example-game} was played using that program.

\vskip .3in
\noindent
{\it Answer to Question \ref{question:1}}:
$u = (2a+1)^2$, $v = (2b+1)^2$, $w = (2c+1)^2$, $x = (2d+1)^2$.

\vskip .3in
\begin{thebibliography}{99}


\bibitem[CR22]{CR22} 
B. Codenotti, G. Resta, {\bf Rithmomachia: alla riscoperta di un gioco medievale},
Messaggerie Scacchistiche, Italy, 2022. ({\it in Italian})

\bibitem[Co84]{Co84} 
M.E. Coughtrie,
{\bf Rhythmomachia : a propaedeutic game of the middle ages},
PhD thesis (Philosophy Dept, Univ. Cape Town), 1984.

\bibitem[dB56]{dB56}
Claude de Boissi\`ere,
{\bf The Very Excellent and Ancient Pythagorean Game, called Rithmomachia}
(in French), 1556.

\bibitem[DG21]{DG21} 
R. Durnil and T. Guardia
{\it Rithmomachia: An Academic Proposal of Rules},
Espacio Matem\'atico Vol. 2 No. 2 (2021) 130--143.

\bibitem[Fo03]{Fo03} Menso Folkerts,
{\it ``Rithmomachia,'' a Mathematical Game from the Middle Ages},
in {\bf Essays on Early Medieval Mathematics} (by Menso Folkerts),
Routledge, 2003.

\bibitem[MacT]{MacT} Herman the Lame biography at
\newline
\url{https://mathshistory.st-andrews.ac.uk/Biographies/Hermann_of_Reichenau/}

\bibitem[Ma83]{Ma83}
Michael Masi, {\bf Boethian Number Theory: A Translation of the
  {\it De Institutione Arithmetica} (with Introduction and Notes)},
(Studies in Classical Antiquity, 6.) Amsterdam, 1983. 
  
\bibitem[Mo01]{Mo01}
Ann Moyer,
{\bf The Philosophers' Game: Rithmomachia in Medieval and Renaissance Europe},
University of Michigan Press, 2001.

\bibitem[Ne98]{Ne98}
M. Necek, {\it Rithmomachia: The Lost Mathematical Treasure of the Dark Ages},
honors thesis, 1998.

\bibitem[Ni00]{Ni00} Nichomachus,
{\bf Introduction to Arithmetic}, Greece, circa 100 CE.
\newline
(translated in \cite{Ma83})

\bibitem[PG22]{PG22}
I. Papadaki, A. Gagatsis,
{\it Rythmomachia: A pedagogical arithmetic game},
Mediterranean Journal for Research in Mathematics
Education, 19 (2022) 82--94.

\bibitem[Pl00]{Pl00} Plutarch, {\bf Moralia, V}, Greece, circa 100 CE.
\newline
\url{https://penelope.uchicago.edu/Thayer/E/Roman/Texts/Plutarch/Moralia/Isis_and_Osiris*/B.html}.  
  
\bibitem[Ri46]{Ri46} J.F.C. Richards,
{\it Boissi\`ere's Pythagorean game,} 
Scripta Mathematica 12(1946)177-217.

\bibitem[SE11]{SE11} 
D. E. Smith and C. C. Eaton, {\it Rithmomachia, The Great Medieval Number Game},
American Mathematical Monthly 18 (1911) 79-80.


\end{thebibliography}

%\vskip .2in
%\begin{center}
%    \adjustimage{max size={1.0\linewidth}{0.9\paperheight}}{rithmomachia-initial-state_matplotlib4}
%\end{center}


\printindex

\newpage

\begin{sidewaysfigure}
 \includegraphics[height=11.5cm,width=20.5cm]{rithmomachia-initial-state_matplotlib4}
 \caption{The rithmomachia board.}
%\label{swim}
\end{sidewaysfigure}



\end{document}
